\documentclass[11pt,a4paper]{article}
\usepackage[T2A]{fontenc}
\usepackage[utf8]{inputenc}
\usepackage[russian]{babel}
\usepackage{amsmath,amssymb,amsthm,mathtools}
\usepackage{tikz}
\usetikzlibrary{arrows.meta,patterns}
\usepackage{needspace}
\usepackage{capt-of}
\definecolor{videolink}{RGB}{0,68,155}
\usepackage{hyperref}
\hypersetup{colorlinks=true,linkcolor=black,urlcolor=videolink,
  pdftitle={Спектральная теория операторов. Подробный конспект},
  pdfsubject={Лекции и семинары А. А. Толченникова}}
\setlength{\emergencystretch}{2em}
\renewcommand{\theHsection}{\arabic{part}.\arabic{section}}
\theoremstyle{plain}
\newtheorem{theorem}{Теорема}[section]
\newtheorem{lemma}[theorem]{Лемма}
\newtheorem{proposition}[theorem]{Утверждение}
\newtheorem{corollary}[theorem]{Следствие}
\newtheorem*{auxlemma}{Лемма}
\theoremstyle{definition}
\newtheorem{definition}[theorem]{Определение}
\newtheorem{example}[theorem]{Пример}
\newtheorem{exercise}[theorem]{Задача}
\newtheorem*{remark}{Замечание}
\newtheorem*{reminder}{Напоминание}

\newcommand{\CC}{\mathbb C}
\newcommand{\RR}{\mathbb R}
\newcommand{\ZZ}{\mathbb Z}
\newcommand{\Hh}{\mathcal H}
\newcommand{\Dom}{\mathcal D}
\newcommand{\id}{I}
\newcommand{\dd}{\,\mathrm d}
\newcommand{\ddiff}[2][]{\frac{\delta^{#1}#2}{\delta x^{#1}}}
\newcommand{\ii}{\mathrm i}
\newcommand{\ee}{\mathrm e}
\newcommand{\abs}[1]{\lvert#1\rvert}
\newcommand{\norm}[1]{\lVert#1\rVert}
\newcommand{\sch}[1]{C_c^\infty(#1)}
\DeclareMathOperator{\Ran}{Ran}
\DeclareMathOperator{\dist}{dist}
\DeclareMathOperator{\ad}{ad}
\DeclareMathOperator{\tr}{tr}
\DeclareMathOperator{\supp}{supp}
\newcommand{\video}[1]{\href{https://www.youtube.com/watch?v=#1}{видеозапись}}
\newcommand{\at}[3]{\href{https://www.youtube.com/watch?v=#1&t=#2s}{#3}}
\newcommand{\record}[2]{\noindent\emph{#1. \video{#2}.}\par\nobreak\medskip}
\newcommand{\topic}[2]{\Needspace{5\baselineskip}\subsection*{#1}\noindent\emph{#2.}\par\nobreak\smallskip}
\newcommand{\sessionsection}[1]{%
  \Needspace{8\baselineskip}\refstepcounter{section}%
  \section*{#1}\addcontentsline{toc}{section}{#1}}

\title{Спектральная теория операторов\\Подробный конспект лекций и семинаров}
\author{По занятиям А.\,А.~Толченникова}
\date{Осень 2026 года}

\begin{document}
\maketitle
\noindent Записи курса опубликованы каналом FMMF MSU.
\href{https://www.youtube.com/playlist?list=PLa4X-npIEnXk}{Плейлист лекций}
содержит 12 записей, включая две части первой лекции.
\href{https://www.youtube.com/playlist?list=PLaHTNEue0Egw}{Плейлист семинаров}
содержит 11 записей. Ссылки возле подразделов ведут к соответствующим
местам записей. Скалярное произведение линейно по первому аргументу.

\tableofcontents

\part{Лекции}

\sessionsection{Лекция 1. Функциональное исчисление}
\noindent\emph{4 сентября. \href{https://www.youtube.com/watch?v=gP2CCEuVufo}{Часть 1},
\href{https://www.youtube.com/watch?v=VzC1xw0kzKY}{часть 2}.}\par\medskip
Во вступлении намечены функциональное исчисление, теория
неограниченных и самосопряжённых операторов, операторные
полугруппы и приложения.

\topic{Функция от ограниченного оператора}{\at{gP2CCEuVufo}{156}{Часть 1, 2:36}}
\begin{example}[Диагональная матрица]
Если $A$ действует в одномерном пространстве умножением на
$\lambda$, то $f(A)=f(\lambda)\id$. Для диагональной матрицы
каждое из её координатных подпространств одномерно, поэтому
\[
 A=\operatorname{diag}(\lambda_1,\ldots,\lambda_n)
 \quad\Longrightarrow\quad
 f(A)=\operatorname{diag}(f(\lambda_1),\ldots,f(\lambda_n)).
\]
При блочной диагональной записи это правило означает
$f(A_1\oplus A_2)=f(A_1)\oplus f(A_2)$.
Если $A=S\operatorname{diag}(\lambda_1,\ldots,\lambda_n)S^{-1}$,
то в исходном базисе
$f(A)=S\operatorname{diag}(f(\lambda_1),\ldots,f(\lambda_n))S^{-1}$.
В базисе собственных векторов это правило означает, что
$f(A)v=f(\lambda)v$ для $Av=\lambda v$. Пусть $\lambda_j$
теперь перечисляют различные собственные значения, а $P_j$
~--- ортогональные проекторы на собственные подпространства
самосопряжённой матрицы. Тогда
\[
 P_jP_k=\delta_{jk}P_j,\qquad \sum_jP_j=\id,\qquad
 A=\sum_j\lambda_jP_j,\qquad f(A)=\sum_j f(\lambda_j)P_j.
\]
Отсюда $f(A)+g(A)=(f+g)(A)$ и
$f(A)g(A)=\sum_j f(\lambda_j)g(\lambda_j)P_j=(fg)(A)$.
В частности, любые две функции от одной и той же матрицы
коммутируют. Значения $f$ нужны только на множестве $\{\lambda_j\}$.
На доске отдельно выписаны также $1(A)=\id$,
$f(\lambda\id)=f(\lambda)\id$ и $x(A)=A$ для
координатной функции $x$ \at{gP2CCEuVufo}{673}{часть 1, 11:13}.
\end{example}

\begin{definition}[Функциональное исчисление]
Пусть $\mathcal F$~--- унитальная алгебра функций, содержащая
координатную функцию $x\mapsto x$, а $\mathcal A$~---
алгебра операторов с единицей. Функциональное исчисление для
$A\in\mathcal A$ задаётся гомоморфизмом
$\mu_A:\mathcal F\to\mathcal A$, удовлетворяющим
$\mu_A(1)=\id$ и $\mu_A(x)=A$. Пишем $f(A)=\mu_A(f)$.
Когда алгебры снабжены топологиями, от исчисления требуют
непрерывности в соответствующих топологиях.
\end{definition}
Для $p(x)=\sum_{j=0}^{m}a_jx^j$ эти требования однозначно
дают $p(A)=\sum_{j=0}^{m}a_jA^j$. Для целой $f$
пишут $f(z)=\sum_{j\ge0}a_jz^j$ и подставляют $A$.
Если $A$ ограничен, сходимость по норме следует из
$\norm{a_jA^j}\le |a_j|\norm A^j$ и бесконечного радиуса
сходимости.

\begin{theorem}[Голоморфное функциональное исчисление]
\label{thm:holomorphic-functional-calculus}
Пусть $A$~--- ограниченный оператор в гильбертовом
пространстве $\Hh$, $U$~--- открытая окрестность
$\sigma(A)$, $f\in\mathcal O(U)$. Пусть $\gamma$~---
положительно ориентированная кусочно-гладкая граница
открытого множества $D$, причём
$\sigma(A)\subset D$ и $\overline D$ компактно и лежит
в $U$. Допускается несколько компонент границы.
Тогда формула
\[
 R_\lambda(A)=(\lambda\id-A)^{-1},\qquad
 f(A)=\frac{1}{2\pi\ii}\int_\gamma f(\lambda)R_\lambda(A)\dd\lambda .
\]
не зависит от выбора допустимой границы и определяет
унитальный гомоморфизм с указанным значением на координате.
\end{theorem}

\begin{center}
\begin{minipage}{\linewidth}
\centering
\begin{tikzpicture}[x=.020cm,y=-.020cm,>=Stealth]
 % Внешняя граница U, а внутри неё контур интегрирования.
 \draw (99,179)
  .. controls (100,134) and (161,80) .. (205,56)
  .. controls (239,36) and (275,20) .. (315,22)
  .. controls (358,24) and (402,30) .. (431,44)
  .. controls (456,54) and (443,72) .. (460,91)
  .. controls (476,111) and (499,122) .. (489,149)
  .. controls (472,195) and (405,220) .. (336,239)
  .. controls (289,252) and (246,260) .. (204,255)
  .. controls (152,251) and (93,233) .. cycle;
 \draw (119,176)
  .. controls (128,132) and (170,91) .. (213,67)
  .. controls (248,46) and (278,37) .. (316,39)
  .. controls (373,40) and (437,55) .. (451,102)
  .. controls (471,161) and (416,190) .. (368,207)
  .. controls (306,228) and (251,234) .. (204,240)
  .. controls (182,244) and (157,234) .. (140,224)
  .. controls (125,215) and (114,196) .. cycle;
 \draw[->] (140,224) .. controls (157,234) and (182,244) .. (204,240);
 % Отдельная заштрихованная область спектра.
 \draw[pattern=north east lines] (141,172)
  .. controls (152,134) and (184,99) .. (220,77)
  .. controls (244,62) and (269,56) .. (291,64)
  .. controls (309,72) and (325,68) .. (343,67)
  .. controls (382,62) and (421,77) .. (433,106)
  .. controls (449,144) and (410,169) .. (373,183)
  .. controls (323,201) and (263,211) .. (215,213)
  .. controls (178,214) and (131,207) .. cycle;
 \draw[->] (106,154)--(417,118);
 \draw[->] (241,213)--(214,13);
 \node[fill=white,inner sep=2pt] at (292,157) {$\sigma(A)$};
 \node at (181,224) {$\gamma$};
 \node at (273,267) {$U$};
\end{tikzpicture}
\captionof{figure}{Спектр $\sigma(A)$ и положительно ориентированный
контур $\gamma$ внутри окрестности голоморфности $U$.
По доске в \at{gP2CCEuVufo}{1540}{части 1, 25:40}.}\label{fig:holomorphic-calculus-contour}
\end{minipage}
\end{center}

\begin{proof}
На лекции свобода выбора контура используется при
доказательстве мультипликативности
\at{gP2CCEuVufo}{1962}{часть 1, 32:42}. Сначала обоснуем
независимость интеграла от выбора допустимого контура.
Подынтегральная функция $f(\lambda)R_\lambda(A)$
голоморфна по операторной норме на $U\setminus\sigma(A)$.

Возьмём произвольные допустимые границы $\delta_j=\partial E_j$, $j=1,2$.
Внутри $E_1\cap E_2$ выберем открытое множество $E_0$
с кусочно-гладкой границей $\delta_0$ так, чтобы
$\sigma(A)\subset E_0$ и $\overline{E_0}\subset E_1\cap E_2$.
Это возможно по компактности спектра. Для любых $x,y\in\Hh$
функция $\lambda\mapsto f(\lambda)(R_\lambda(A)x,y)$
голоморфна между $\delta_j$ и $\delta_0$.
По теореме Коши
\[
 \int_{\delta_j}f(\lambda)(R_\lambda(A)x,y)\dd\lambda
 =
 \int_{\delta_0}f(\lambda)(R_\lambda(A)x,y)\dd\lambda .
\]
В границу $E_j\setminus\overline{E_0}$ контур $\delta_0$
входит с обратной ориентацией. Переход к пределу в суммах
Римана позволяет вынести скалярное произведение за интеграл.
Равенство для всех $x,y$ даёт совпадение операторных
интегралов по $\delta_j$ и $\delta_0$ при $j=1,2$.

Для $a,b\in\CC$ равенство
$(af+bg)(A)=af(A)+bg(A)$ следует из линейности интеграла
на одном и том же контуре. Функции $1$ и $z$ определены
на всей плоскости. Для них можно взять окружность
$|\lambda|=R$ с $R$ достаточно большим, чтобы она
охватывала исходную область и $R>\norm A$. На окружности
\[
 (\lambda\id-A)^{-1}
 =\lambda^{-1}(\id-\lambda^{-1}A)^{-1}
 =\sum_{k=0}^\infty \lambda^{-k-1}A^k .
\]
Ряд равномерно сходится в операторной норме, поскольку
нормы слагаемых ограничены членами геометрического ряда
с отношением $\norm A/R<1$. Ряд для $\lambda R_\lambda$
также равномерно сходится, и оба ряда можно интегрировать
почленно. Для целого $m$ интеграл
$(2\pi\ii)^{-1}\int_{|\lambda|=R}\lambda^m\dd\lambda$
равен $1$ при $m=-1$ и $0$ во всех остальных случаях.
В ряду для $R_\lambda$ степень $-1$ получается при $k=0$,
а в ряду для $\lambda R_\lambda$~--- при $k=1$. Отсюда
$1(A)=A^0=\id$ и $z(A)=A^1=A$.

Для доказательства мультипликативности интегрируем
по положительно ориентированным границам двух вложенных
областей $\sigma(A)\subset D_1$, $\overline{D_1}\subset D_2$,
$\overline{D_2}\subset U$ (рис.~\ref{fig:holomorphic-calculus-nested}).
Функцию $f(A)$ вычисляем по внешней границе
$\gamma_2=\partial D_2$, а $g(A)$~--- по внутренней
$\gamma_1=\partial D_1$.

\begin{center}
\begin{minipage}{\linewidth}
\centering
\begin{tikzpicture}[>=Stealth]
 \draw (-.05,0) ellipse (2.75 and 1.55);
 \draw (-.12,.08) ellipse (2.22 and 1.24);
 \draw (-.23,.11) ellipse (1.66 and .94);
 \draw[pattern=north east lines] (-.35,.13) ellipse (1.12 and .68);
 \node[fill=white,inner sep=2pt] at (-.35,.13) {$\sigma(A)$};
 \node[fill=white,inner sep=1pt] at (-.62,.90) {$D_1$};
 \node[fill=white,inner sep=1pt] at (.41,1.13) {$D_2$};
 \node[fill=white,inner sep=1pt] at (-.62,1.425) {$U$};
\end{tikzpicture}
\captionof{figure}{Вложенные области $\sigma(A)\subset D_1$,
$\overline{D_1}\subset D_2$, $\overline{D_2}\subset U$.
В доказательстве мультипликативности $\mu$ пробегает
$\partial D_1$, а $\lambda$~--- $\partial D_2$.
По доске в \at{gP2CCEuVufo}{1970}{части 1, 32:50}.}\label{fig:holomorphic-calculus-nested}
\end{minipage}
\end{center}

Вычисляем, как на доске \at{gP2CCEuVufo}{2200}{часть 1, 36:40}:
\[
 \begin{aligned}
 (\lambda\id-A)(R_\lambda-R_\mu)
 &=\id-(\lambda\id-A)R_\mu\\
 &=\underbrace{\id-(\mu\id-A)R_\mu}_{0}
    +(\mu-\lambda)R_\mu .
 \end{aligned}
\]
Умножение слева на $R_\lambda$ даёт резольвентное тождество
\[
 R_\lambda R_\mu
 =\frac{R_\lambda-R_\mu}{\mu-\lambda}.
\]
Поэтому
\[
 \begin{aligned}
 f(A)g(A)
 &=\frac{1}{(2\pi\ii)^2}\int_{\gamma_2}\int_{\gamma_1}
    f(\lambda)g(\mu)
    \frac{R_\lambda-R_\mu}{\mu-\lambda}\dd\mu\,\dd\lambda .
 \end{aligned}
\]
В слагаемом с $R_\lambda$ интеграл по $\mu$ равен нулю:
точка $\mu=\lambda$ находится снаружи $\gamma_1$.
В слагаемом с $-R_\mu$ меняем порядок интегрирования.
Для фиксированного $\mu$ внутренний теперь интеграл содержит
$f(\lambda)/(\lambda-\mu)$, и по формуле Коши он равен
$2\pi\ii f(\mu)$. В результате остаётся
$(2\pi\ii)^{-1}\int_{\gamma_1}f(\mu)g(\mu)R_\mu\dd\mu$.
\end{proof}

Для $A=\xi\id$ резольвента равна
$(\lambda-\xi)^{-1}\id$. Поэтому операторный интеграл
на каждом векторе $x$ даёт $f(\xi)x$ по интегральной
формуле Коши \at{gP2CCEuVufo}{1620}{часть 1, 27:00}.
Интеграл понимается как интеграл Римана в операторной
норме. В частности, для фиксированной допустимой
границы
\[
 \norm{f(A)}\le
 \frac{\operatorname{length}(\gamma)}{2\pi}
 \sup_{\lambda\in\gamma}\norm{R_\lambda(A)}
 \sup_{\lambda\in\gamma}|f(\lambda)|.
\]
Таким образом, равномерная сходимость голоморфных функций
на компактах в $U$ влечёт сходимость соответствующих
операторов по норме. Локальное разложение резольвенты,
которое обосновывает её голоморфность, подробно выводится
в лекции~6. При переходе от соглашения той лекции к
$R_\lambda(A)=(\lambda\id-A)^{-1}$ резольвента меняет знак.

\topic{Спектральная мера и формула Фурье}{\at{gP2CCEuVufo}{3054}{Часть 1, 50:54}}
Для самосопряжённого оператора рассматривается запись
$A=\int \lambda\dd E_A(\lambda)$, где $E_A$~---
проекторнозначная спектральная мера. Функция от оператора
задаётся интегралом $f(A)=\int f(\lambda)\dd E_A(\lambda)$.
В конечномерном примере
$A=\sum_j\lambda_j P_j$ формула принимает вид
$f(A)=\sum_j f(\lambda_j)P_j$. На доске для преобразования
Фурье используется симметричная нормировка
\at{gP2CCEuVufo}{3183}{часть 1, 53:03}. При
$f(x)=(2\pi)^{-1/2}\int \widehat f(t)\ee^{\ii tx}\dd t$
подстановка оператора даёт
\[
 f(A)=\frac{1}{\sqrt{2\pi}}\int_{\RR}\widehat f(t)\ee^{\ii tA}\dd t.
\]
Для ограниченного $A$ экспонента задаётся рядом
$\ee^{\ii tA}=\sum_{k\ge0}(\ii t)^k A^k/k!$.
Почленное дифференцирование даёт
$-\ii\partial_t\ee^{\ii tA}=A\ee^{\ii tA}$. Поэтому
отображение $u_0\mapsto\ee^{\ii tA}u_0$ является
оператором решения задачи Коши
$-\ii u'(t)=Au(t)$, $u(0)=u_0$. Для неограниченного
оператора на лекции намечается тот же подход через
задачу Коши при подходящих условиях на $A$
\at{gP2CCEuVufo}{3416}{часть 1, 56:56}.

Для сходимости интеграла по норме достаточно
$\int_{\RR}|\widehat f(t)|\,\norm{\ee^{\ii tA}}\dd t<\infty$.
У самосопряжённого $A$ экспонента унитарна, и достаточно
$\widehat f\in L^1(\RR)$. В частности, подходят функции
Шварца. При произвольном ограниченном $A$ требование к
$\widehat f$ зависит от роста $\norm{\ee^{\ii tA}}$.

Многочлены разбираются в смысле распределений. При той же
нормировке
$\widehat f(t)=(2\pi)^{-1/2}\int_{\RR} f(x)\ee^{-\ii tx}\dd x$
функциям $1$, $x$ и $x^2$ отвечают соответственно
$\sqrt{2\pi}\,\delta$, $\sqrt{2\pi}\,\ii\delta'$ и
$-\sqrt{2\pi}\,\delta''$.
Действие производных $\delta$ задаётся формулами
$\langle\delta',\psi\rangle=-\psi'(0)$ и
$\langle\delta'',\psi\rangle=\psi''(0)$. Подставляя
$\psi(t)=\ee^{\ii tA}$, получаем $\id$, $A$ и
$-(\ii A)^2=A^2$. Это показывает, как запись Фурье
воспроизводит подстановку оператора в многочлен
\at{gP2CCEuVufo}{3682}{часть 1, 1:01:22}.

\topic{Несколько некоммутирующих аргументов}{\at{gP2CCEuVufo}{3977}{Часть 1, 1:06:17}}
Если $f(x,y)=xy=yx$, то при замене $x$ на $A$,
$y$ на $B$ возникают произведения $AB$ и $BA$.
В обсуждении рассмотрена и полусумма $(AB+BA)/2$.
Чтобы задать конкретный порядок действий, вводят
индексы Фейнмана \at{gP2CCEuVufo}{3908}{часть 1, 1:05:08}.
\begin{definition}[Индексы Фейнмана]
Индексы над операторами задают порядок их действия.
Оператор с меньшим индексом действует раньше и в
произведении стоит правее. Например, при $F(x,y)=xy$
\[
 F(\overset{1}A,\overset{2}B)=BA,\qquad
 F(\overset{2}A,\overset{1}B)=AB .
\]
Значим относительный порядок индексов. В выкладках можно
заменять их любыми числами, сохраняющими этот порядок.
\end{definition}
Для разложимого тензора $f(x)g(y)$ правило состоит в
составлении произведения из $f(A)$ и $g(B)$ в указанном
порядке, а для конечной суммы тензоров применяется
линейность. Для продолжения на пополнение тензорного
произведения используют проективную норму
\[
 \norm{w}_\pi=\inf_{w=\sum_k f_k\otimes g_k}
                 \sum_k\norm{f_k}\norm{g_k}.
\]
Соседние в порядке действия множители, относящиеся к одному
оператору, можно объединять. При раздвигании индексов одну переменную
разделяют на две, чтобы отделить соответствующие операторы
от остальных множителей. В качестве исторического примера на лекции
обсуждается упорядочивание по времени в выражении
$\exp(\int_0^T\overset{s}{A(s)}\dd s)$, где меткой порядка
служит сам момент времени (\at{gP2CCEuVufo}{4063}{1:07:43}).

\topic{Разностная производная и единственность}{\at{VzC1xw0kzKY}{25}{Часть 2, 0:25}}
\begin{definition}[Разностная производная]
Обозначим $\ddiff{f}(x,y)=(f(x)-f(y))/(x-y)$, $x\ne y$.
На диагонали полагают $\ddiff{f}(x,x)=f'(x)$, когда
производная определена. Условие на класс одноместных
символов, выписанное на доске, требует
$\ddiff{f}\in\mathcal F\widehat\otimes_\pi\mathcal F$
для рассматриваемых $f\in\mathcal F$.
\end{definition}
\begin{theorem}[Единственность исчисления функций одного аргумента]
Пусть выполнено условие на разностную производную.
Если два непрерывных унитальных гомоморфизма
$\mu_1,\mu_2:\mathcal F\to\mathcal A$ принимают одинаковое
значение на координатной функции, то $\mu_1=\mu_2$.
\end{theorem}
\begin{proof}
Положим $A=\mu_1(x)$, $B=\mu_2(x)$. Пишем
$f(A)=\mu_1(f)$ и $f(B)=\mu_2(f)$, сохраняя разные
буквы для двух исчислений. Для произвольного
$f\in\mathcal F$ применяем тождество
$f(x)-f(y)=(x-y)\ddiff{f}(x,y)$ с индексами $2$ у $A$
и $1$ у $B$ \at{VzC1xw0kzKY}{297}{часть 2, 4:57}:
\[
 f(\overset{2}{A})-f(\overset{1}{B})
 =\bigl(\overset{2}{A}-\overset{1}{B}\bigr)
      \ddiff{f}(\overset{2}{A},\overset{1}{B}).
\]
Раздвигаем индексы в разностной производной: у $B$
заменяем $1$ на $0$, а у $A$ заменяем $2$ на $3$.
Порядок множителей сохраняется. Получаем запись с
автономной скобкой, выделенной на доске прямоугольником:
\[
 f(A)-f(B)
 =\boxed{\bigl(\overset{2}{A}-\overset{1}{B}\bigr)}
      \ddiff{f}(\overset{3}{A},\overset{0}{B}).
\]
Сначала действуют множители с индексом~$0$, затем
вся скобка с индексами~$1,2$, затем множители с индексом~$3$.
Поэтому скобка
действует как самостоятельный оператор $A-B$.
По условию $A=B$, так что скобка равна нулю.
Получаем $f(A)=f(B)$ для каждого $f\in\mathcal F$.
\end{proof}

В обсуждении после доказательства лектор поясняет,
как читать запись после раздвигания индексов
\at{VzC1xw0kzKY}{540}{часть 2, 9:00}.
Четырём индексам отвечают четыре аргумента функции
$F(x,y,z,w)=(x-y)\ddiff{f}(z,w)$, куда подставляются
$\overset{2}{A},\overset{1}{B},\overset{3}{A},
\overset{0}{B}$ соответственно. Сначала $F$ приближают
конечными суммами разложимых тензоров. Для слагаемого
$p(x)q(y)r(z)s(w)$ применяют к вектору последовательно
$s(B)$, $q(B)$, $p(A)$ и $r(A)$, поэтому получается
$r(A)p(A)q(B)s(B)$. В частности, для
$\ddiff{f}(z,w)=\sum_j g_j(z)h_j(w)$ после вычитания
слагаемых от $x$ и $y$ выходят операторы
$g_j(A)(A-B)h_j(B)$. При раздвигании индексов переход
от одной записи к другой на тензорах использует
мультипликативность $\mu_1$ и $\mu_2$, а переход к
пределу~--- непрерывность. Условие на $\ddiff{f}$
относится ко всей алгебре $\mathcal F$.

\begin{proposition}[Коммутатор с функцией оператора]
При том же условии на $\ddiff{f}$ и при обозначении
$[A,B]=AB-BA$ справедлива формула
\[
 [A,f(B)]=\sum_j g_j(B)[A,B]h_j(B),
 \qquad \ddiff{f}=\sum_jg_j\otimes h_j .
\]
\end{proposition}
\begin{proof}
В первой части лекции \at{gP2CCEuVufo}{4182}{1:09:42}
лектор записывает разность как функцию трёх аргументов:
\[
 \begin{aligned}
 [A,f(B)]
 &=\overset{1}{A}f(\overset{0}{B})
     -f(\overset{2}{B})\overset{1}{A}
   =F(\overset{0}{B},\overset{2}{B},\overset{1}{A}),\\
 F(x,y,z)&=z\bigl(f(x)-f(y)\bigr)=z(x-y)\ddiff{f}(x,y).
 \end{aligned}
\]
В разностной производной раздвигаем индексы $0$ и $2$
наружу до $-1$ и $3$. Индексы множителей автономной
скобки остаются $0,1,2$:
\[
 [A,f(B)]
 =\boxed{\overset{1}{A}
       \bigl(\overset{0}{B}-\overset{2}{B}\bigr)}
   \ddiff{f}(\overset{-1}{B},\overset{3}{B}) .
\]
Первое слагаемое скобки действует как $AB$, второе~---
как $BA$. Значит, скобка действует как $[A,B]$
между множителями разностной производной. Так как
$\ddiff{f}(x,y)=\ddiff{f}(y,x)$, можем поставить индекс $3$
над первым аргументом разностной производной и $-1$ над
вторым. На тензоре $g_j(x)h_j(y)$ получаем
$g_j(B)[A,B]h_j(B)$. Суммирование и переход к пределу
дают формулу утверждения.
\end{proof}

\sessionsection{Лекция 2. Операции некоммутативного анализа}
\record{11 сентября}{8kQuUxC-VVg}

\topic{Индексы и автономные скобки}{\at{8kQuUxC-VVg}{195}{3:15}}
В начале лекции обсуждается зависимость квантования от
порядка. Для символа $xp$ можно последовательно применить
дифференцирование и умножение, либо выполнить те же действия
в обратном порядке. В этом примере квантование Вейля
соответствует полусумме двух операторов.

\begin{definition}[Генератор и фейнмановский набор]
Пусть $\mathcal F_1$~--- алгебра скалярных функций одной
переменной с поточечным умножением, содержащая $1$ и
координатную функцию $x\mapsto x$, а $\mathcal A$~---
алгебра операторов. Оператор $A\in\mathcal A$ называется
\emph{генератором} (порождающим элементом), если существует
непрерывный гомоморфизм алгебр
$\mu_A:\mathcal F_1\to\mathcal A$ такой, что
$\mu_A(x)=A$ \at{8kQuUxC-VVg}{270}{4:30}.
Гомоморфизм понимается в смысле функционального исчисления
лекции~1, в частности $\mu_A(1)=\id$. Его значение
на $f\in\mathcal F_1$ обозначают через $f(A)$.

Набор генераторов $A_1,\ldots,A_n$ с проставленными
индексами Фейнмана называется \emph{фейнмановским набором}
\at{8kQuUxC-VVg}{387}{6:27}. Для каждого $A_j$ его
исчисление $\mu_{A_j}$ позволяет вычислить $f_j(A_j)$
для функции $f_j$ одной переменной.
От пространства символов $\mathcal F_n$ требуют плотности
алгебраической тензорной степени $\mathcal F_1^{\otimes n}$.
Для $f(x_1,\ldots,x_n)=\prod_{j=1}^n f_j(x_j)$ и
индексов $1,\ldots,n$ по порядку берут
\[
 \mu_{A_1,\ldots,A_n}(f)=f_n(A_n)\cdots f_1(A_1).
\]
Сначала исчисление продолжают линейно на тензорную
степень, затем непрерывно на $\mathcal F_n$ при условии,
что такое продолжение существует.
\end{definition}
Пример на доске~--- $\mathcal F_n=\mathcal O(\CC^n)$
с равномерной сходимостью на компактах, $\mathcal A=\mathcal B(\Hh)$.
В этой топологии многочлены приближают целую функцию по её
ряду, а мономы разложимы. Общий приём доказательства
тождеств состоит в проверке на разложимых символах и
применении линейности и непрерывности.

\begin{lemma}[Склейка и раздвигание индексов]
Пусть в упорядоченной подстановке два соседних по порядку
действия генератора равны $A$. Тогда их можно слить в одну
подстановку $A$, ограничив символ на диагональ
соответствующих двух переменных. Обратная операция
раздвигает индекс на два соседних.
\end{lemma}
\begin{proof}
Для $f(x,y)=h(x)g(y)$ получается $g(A)h(A)$.
Так как $\mu_A$~--- гомоморфизм, это равно
$(gh)(A)$. Скалярный символ $(i^*f)(x)=f(x,x)$
как раз равен $g(x)h(x)$. Если есть другие переменные,
их множители стоят по обе стороны склеиваемой пары.
При $A^1,B^2,A^3$ между двумя копиями $A$ стоит $B$,
и довод о соседних множителях применить нельзя.
\end{proof}

Если два индекса $k,\ell$ стоят над одним и тем же
$A$, а все остальные индексы $s_1,\ldots,s_m$ лежат
вне отрезка с концами $k,\ell$, можно поставить над
обеими копиями $A$ любой общий индекс $\alpha$ из
этого отрезка, а затем ограничить символ на диагональ:
\[
 \begin{aligned}
 &f(\overset{k}{A},\overset{\ell}{A},
       \overset{s_1}{B_1},\ldots,\overset{s_m}{B_m})\\
 &\qquad=(i^*f)(\overset{\alpha}{A},
       \overset{s_1}{B_1},\ldots,\overset{s_m}{B_m}),
 \qquad (i^*f)(x,y_1,\ldots,y_m)=f(x,x,y_1,\ldots,y_m).
 \end{aligned}
\]
Это обобщение лектор выписывает в
\at{8kQuUxC-VVg}{1153}{19:13}.

\begin{lemma}[Выделение множителя, или автономные скобки]
Пусть
$F(x_1,\ldots,x_n)=G(x_1,\ldots,x_k)H(x_{k+1},\ldots,x_n)$.
Пусть индексы генераторов первой группы все лежат в
интервале $(a,b)$, а индексы второй группы находятся вне
его. Тогда в подстановке к $zH(x_{k+1},\ldots,x_n)$
вместо $z$ можно подставить оператор
$G(\overset{j_1}{A_1},\ldots,\overset{j_k}{A_k})$
с любым индексом из $(a,b)$.
\end{lemma}
\begin{proof}
Для разложимых $G$ и $H$ все множители первой группы
стоят в произведении единым блоком. Его внутренний
порядок задают $j_1,\ldots,j_k$. Индекс каждого множителя
$H$ лежит по одну сторону от $(a,b)$, поэтому замена
блока одним оператором сохраняет порядок. К общим
символам переходим по непрерывности.
\end{proof}

\begin{example}[Автономная скобка]
В выражении
\[
 \bigl(\overset{2}{A}-\overset{3}{A}\bigr)
       \sin\bigl(\overset{1}{B}+\overset{4}{B}\bigr)
 =\overset{2}{\boxed{A-A}}\,
       \sin\bigl(\overset{1}{B}+\overset{4}{B}\bigr)=0
\]
индексы $2,3$ образуют один блок между $1$ и $4$.
Если индексы двух групп чередуются, выделить такую
скобку уже нельзя \at{8kQuUxC-VVg}{1829}{30:29}.
\end{example}

\topic{Левое и правое регулярные представления}{\at{8kQuUxC-VVg}{1966}{32:46}}
\begin{definition}[Регулярные представления]
В унитальной алгебре $\mathcal A$ положим
$L_A X=AX$ и $R_A X=XA$ для $X\in\mathcal A$.
Тогда $L_A L_B=L_{AB}$, $R_AR_B=R_{BA}$, а
$L_AR_B=R_BL_A$, поскольку обе композиции
вычисляются как $X\mapsto AXB$.
\end{definition}

\begin{theorem}[Критерий генератора через регулярные представления]
\at{8kQuUxC-VVg}{2124}{35:24}.
При условии на разностную производную эквивалентны
следующие утверждения:
\begin{enumerate}
 \item $A$~--- генератор в $\mathcal A$.
 \item $L_A$~--- генератор в $\mathcal L(\mathcal A)$.
 \item $R_A$~--- генератор в $\mathcal L(\mathcal A)$.
\end{enumerate}
При этом $f(L_A)=L_{f(A)}$ и $f(R_A)=R_{f(A)}$.
\end{theorem}
\begin{proof}[Доказательство прямых импликаций]
Если $A$~--- генератор, положим $\mu_{L_A}=L\circ\mu_A$.
Это композиция двух непрерывных гомоморфизмов, и
$\mu_{L_A}(x)=L_{\mu_A(x)}=L_A$. Так получаем
$(1)\Rightarrow(2)$. Для $(1)\Rightarrow(3)$ берём
$\mu_{R_A}=R\circ\mu_A$. Коммутативность $\mathcal F_1$
даёт $f(A)g(A)=g(A)f(A)$, поэтому
\[
 R_{f(A)}R_{g(A)}=R_{g(A)f(A)}=R_{(fg)(A)}.
\]
Следовательно, $\mu_{R_A}$ также непрерывный
гомоморфизм со значением $R_A$ на $x$.
\end{proof}

\begin{auxlemma}[Коммутирование левых и правых регулярных представлений]
\at{8kQuUxC-VVg}{2318}{38:38}.
Пусть $A,B\in\mathcal A$.
\begin{enumerate}
 \item[(а)] $L_A R_B=R_B L_A$.
 \item[(б)] Если $L_A$~--- генератор и для $f\in\mathcal F_1$
 выполнено условие на разностную производную, то
 $f(L_A)R_B=R_B f(L_A)$.
\end{enumerate}
\end{auxlemma}
\begin{proof}
Для пункта (а) при любом $X\in\mathcal A$ имеем
$L_A R_B X=AXB=R_B L_A X$.
Для пункта (б) подставляем
$\overset{1}{L_A},\overset{3}{L_A},\overset{2}{R_B}$
в $F(x,y,z)=(f(y)-f(x))z$:
\[
 \begin{aligned}
 f(L_A)R_B-R_B f(L_A)
 &=F(\overset{1}{L_A},\overset{3}{L_A},\overset{2}{R_B})\\
 &=\overset{2}{R_B}\bigl(\overset{3}{L_A}-\overset{1}{L_A}\bigr)
    \ddiff{f}(\overset{3}{L_A},\overset{1}{L_A})\\
 &=\boxed{\overset{2}{R_B}
       \bigl(\overset{3}{L_A}-\overset{1}{L_A}\bigr)}
    \ddiff{f}(\overset{4}{L_A},\overset{0}{L_A}).
 \end{aligned}
\]
Раздвигание индексов разностной производной до $0$ и $4$
выделяет автономную скобку с порядком действий $1,2,3$.
Она равна $L_A R_B-R_B L_A=0$ по пункту (а).
\end{proof}

\begin{proof}[Доказательство обратных импликаций]
Докажем $(2)\Rightarrow(1)$. Пусть дано исчисление
$\mu_{L_A}$. По пункту (б) леммы $f(L_A)$ коммутирует
с $R_X$ для любого $X\in\mathcal A$. Тогда
\[
 f(L_A)X=f(L_A)R_X\id
         =R_Xf(L_A)\id=\bigl(f(L_A)\id\bigr)X.
\]
Поэтому полагаем $\mu_A(f)=f(L_A)\id$ и получаем
$f(L_A)=L_{\mu_A(f)}$. Отображение $\mu_A$ линейно
и непрерывно, так как это композиция $\mu_{L_A}$
с применением оператора к $\id$. Имеем
$\mu_A(1)=\id$ и $\mu_A(x)=L_A\id=A$. Проверим мультипликативность:
\[
 \begin{aligned}
 \mu_A(f)\mu_A(g)
   &=\bigl(f(L_A)\id\bigr)\bigl(g(L_A)\id\bigr)
   =f(L_A)\bigl(g(L_A)\id\bigr)\\
   &=(fg)(L_A)\id=\mu_A(fg).
 \end{aligned}
\]
Так построено функциональное исчисление для $A$
\at{8kQuUxC-VVg}{2719}{45:19}.

Для $(3)\Rightarrow(1)$ тот же ход использует
коммутирование $f(R_A)$ с $L_X$. Поэтому
$f(R_A)X=f(R_A)L_X\id=X(f(R_A)\id)$. Полагаем
$\mu_A(f)=f(R_A)\id$. Тогда из
$f(R_A)=R_{\mu_A(f)}$ получаем
$\mu_A(fg)=\mu_A(g)\mu_A(f)$. Поменяв $f$ и $g$ и
учтя $fg=gf$, получаем мультипликативность.
Линейность, непрерывность и значения на $1$ и $x$
проверяются применением к $\id$. Единственность
исчисления даёт указанные формулы для $f(L_A)$ и $f(R_A)$.
\end{proof}

\begin{theorem}[Функции от нескольких регулярных представлений]
Для разложимого $F(x,y)=g(x)h(y)$ имеем
\[
 F(L_A,R_B)X=g(A)Xh(B)
      =(zF)(\overset3 A,\overset1 B,\overset2 X).
\]
В последней записи $z$ отвечает аргументу $X$.
Для общего двухместного символа обе части понимаются
по линейности и непрерывности.
\end{theorem}
\begin{proof}
На произведении $g(x)h(y)$ по предыдущей теореме имеем
\[
 g(L_A)h(R_B)X=L_{g(A)}R_{h(B)}X=g(A)Xh(B).
\]
Поскольку $L_A$ и $R_B$ коммутируют, для этих двух
представлений относительный порядок безразличен.
\end{proof}
Для нескольких левых представлений берём разложимый
символ $f=\bigotimes_{j=1}^n f_j$. Тогда
\[
 f(\overset1{L_{A_1}},\ldots,\overset n{L_{A_n}})
   =L_{\,f_n(A_n)\cdots f_1(A_1)} .
\]
Для правых представлений та же цепочка даёт
$R_{\,f_1(A_1)\cdots f_n(A_n)}$, потому что
$R_CR_D=R_{DC}$. Иными словами, индексы аргументов
под правым представлением разворачиваются.
Для смешанного набора с индексами $j_1,\ldots,j_n$ у
левых представлений и $\alpha_1,\ldots,\alpha_m$ у
правых запись с доски имеет вид
\[
 \begin{aligned}
 &F(\overset{j_1}{L_{A_1}},\ldots,\overset{j_n}{L_{A_n}},
    \overset{\alpha_1}{R_{B_1}},\ldots,
    \overset{\alpha_m}{R_{B_m}})(X)\\
 &\quad=(zF)\bigl(
    \overset{j_1+c}{A_1},\ldots,\overset{j_n+c}{A_n},
    \overset{-\alpha_1-d}{B_1},\ldots,
    \overset{-\alpha_m-d}{B_m},\overset{0}{X}\bigr).
 \end{aligned}
\]
Здесь $z$~--- последний аргумент, а $c,d$ выбирают так,
чтобы новые индексы группы $A$ были положительными, а
группы $B$~--- отрицательными. На разложимом символе
левая группа умножает $X$ слева в своём порядке, правая
умножает справа в обратном порядке. Группы можно
переставлять между собой благодаря $L_AR_B=R_BL_A$
\at{8kQuUxC-VVg}{3770}{1:02:50}.

\topic{Гомоморфизмы и переплетение}{\at{8kQuUxC-VVg}{3960}{1:06:00}}
\begin{theorem}[Сохранение функционального исчисления]
Для генераторов с одноместным исчислением, удовлетворяющим
условию на разностные производные, справедливы следующие свойства.
\begin{enumerate}
\item Если $\kappa$~--- непрерывный унитальный гомоморфизм
операторных алгебр, то $\kappa(A)$~--- генератор с
$f(\kappa(A))=\kappa(f(A))$.
\item Если $AB=BA$, то $f(A)B=Bf(A)$.
\item Если $AB=BC$ для генераторов $A,C$ одной алгебры,
то $f(A)B=Bf(C)$.
\item Если $A$ и $C$ действуют в пространствах $V$ и $W$
соответственно, а $S:W\to V$ удовлетворяет $AS=SC$, то
$f(A)S=Sf(C)$ в совместном блочном исчислении.
\item Если $Av=\lambda v$, то $f(A)v=f(\lambda)v$.
\end{enumerate}
\end{theorem}
\begin{proof}
Для первого свойства берём композицию
$\mu_{\kappa(A)}=\kappa\circ\mu_A$. Она непрерывна,
сохраняет произведение и единицу и переводит $x$ в
$\kappa(A)$. По единственности
$f(\kappa(A))=\mu_{\kappa(A)}(f)=\kappa(f(A))$.
Во втором пункте $[A,B]=0$, поэтому в коммутаторной
формуле из лекции~1 каждое слагаемое содержит нулевую
автономную скобку. Для третьего пункта лектор ссылается
на тот же приём \at{8kQuUxC-VVg}{4371}{1:12:51}.
Берём
$F(x,y,z)=z(f(x)-f(y))$ с индексами $x=\overset{3}{A}$,
$y=\overset{1}{C}$, $z=\overset{2}{B}$. Вынося
$(x-y)z$ и раздвигая внешние индексы разностной
производной, получаем скобку $AB-BC$ и равенство
\[
 f(A)B-Bf(C)=\sum_j g_j(A)(AB-BC)h_j(C)=0.
\]
Для четвёртого пункта перемножаем блочные матрицы
в $V\oplus W$:
\[
 \begin{pmatrix}A&0\\0&C\end{pmatrix}
 \begin{pmatrix}0&S\\0&0\end{pmatrix}
 =\begin{pmatrix}0&AS\\0&0\end{pmatrix}
 =\begin{pmatrix}0&SC\\0&0\end{pmatrix}
 =\begin{pmatrix}0&S\\0&0\end{pmatrix}
 \begin{pmatrix}A&0\\0&C\end{pmatrix}.
\]
По второму пункту можно заменить диагональную матрицу
на $\operatorname{diag}(f(A),f(C))$ и получить равенство
$f(A)S=Sf(C)$ в верхнем правом блоке.
Для последнего свойства выбираем $W=\CC$, $C=\lambda\id$
и $S(c)=cv$. Тогда $AS(c)=cAv=c\lambda v=SC(c)$,
поэтому по предыдущему пункту
$f(A)v=f(A)S(1)=Sf(C)(1)=S(f(\lambda))=f(\lambda)v$.
\end{proof}

\begin{figure}[ht]
\centering
\begin{tikzpicture}[>=Stealth]
 \node (vl) at (0,1.25) {$V$};
 \node (vr) at (3,1.25) {$V$};
 \node (wl) at (0,0) {$W$};
 \node (wr) at (3,0) {$W$};
 \draw[->] (vl)-- node[above] {$A$} (vr);
 \draw[->] (wl)-- node[below] {$C$} (wr);
 \draw[->] (wl)-- node[left] {$S$} (vl);
 \draw[->] (wr)-- node[right] {$S$} (vr);
 \begin{scope}[xshift=5.1cm]
  \node (evl) at (0,1.25) {$V$};
  \node (evr) at (3,1.25) {$V$};
  \node (cl) at (0,0) {$\CC$};
  \node (cr) at (3,0) {$\CC$};
  \draw[->] (evl)-- node[above] {$A$} (evr);
  \draw[->] (cl)-- node[below] {$\lambda\id$} (cr);
  \draw[->] (cl)-- node[left] {$S$} (evl);
  \draw[->] (cr)-- node[right] {$S$} (evr);
 \end{scope}
\end{tikzpicture}
\caption{Переплетение $A$ и $C$ (слева) и применение к
собственному вектору, где $S(c)=cv$ (справа).
По доске в \at{8kQuUxC-VVg}{4490}{1:14:50}
и \at{8kQuUxC-VVg}{4774}{1:19:34}.}
\end{figure}

\topic{Дифференцирование}{\at{8kQuUxC-VVg}{4800}{1:20:00}}
\begin{definition}[Дифференцирование]
Линейное отображение $\mathcal D:\mathcal A\to\mathcal A$
называется дифференцированием, если
$\mathcal D(XY)=(\mathcal DX)Y+X(\mathcal DY)$.
Внутреннее дифференцирование имеет вид
$\mathcal DX=[B,X]$. Правило Лейбница в этом случае
получается вставкой и вычитанием $XBY$:
\[
 [B,XY]=BXY-XBY+XBY-XYB=[B,X]Y+X[B,Y].
\]
\end{definition}

\begin{lemma}[Дифференцирование в левом регулярном представлении]
\label{lem:derivation-left-regular}
\at{8kQuUxC-VVg}{4997}{1:23:17}.
При переходе к левому регулярному представлению
дифференцирование $\mathcal D$ действует коммутатором:
$L_{\mathcal D A}=[\mathcal D,L_A]$ (рис.~\ref{fig:derivation-left-regular}).
\end{lemma}

\begin{center}
\begin{minipage}{\linewidth}
\centering
\begin{tikzpicture}[>=Stealth]
 \node (al) at (0,1.35) {$\mathcal A$};
 \node (ar) at (3.5,1.35) {$\mathcal A$};
 \node (ll) at (0,0) {$\mathcal L(\mathcal A)$};
 \node (lr) at (3.5,0) {$\mathcal L(\mathcal A)$};
 \draw[->] (al)--node[above] {$\mathcal D$} (ar);
 \draw[->] (al)--node[left] {$L$} (ll);
 \draw[->] (ar)--node[right] {$L$} (lr);
 \draw[->] (ll)--node[below] {$[\mathcal D,\,\cdot\,]$} (lr);
\end{tikzpicture}
\captionof{figure}{Диаграмма к лемме~\ref{lem:derivation-left-regular}.
По доске в \at{8kQuUxC-VVg}{5145}{1:25:45}.}
\label{fig:derivation-left-regular}
\end{minipage}
\end{center}

\begin{proof}
Для любого $X\in\mathcal A$ имеем
\[
 [\mathcal D,L_A]X=\mathcal D(AX)-A\mathcal DX
       =(\mathcal DA)X=L_{\mathcal DA}X.
\]
\end{proof}

\begin{theorem}[Формула дифференцирования сложной функции]
При выполнении условий исчисления для $\ddiff{f}$
и $\mathcal D\in\mathcal L(\mathcal A)$
\[
 \mathcal D f(A)=\overset{2}{\mathcal D(A)}\,
        \ddiff{f}(\overset{1}{A},\overset{3}{A}).
\]
В терминах левых и правых умножений та же формула имеет вид
\[
 \mathcal D f(A)=\ddiff{f}(L_A,R_A)(\mathcal DA).
\]
В тензорной записи это равенство имеет вид
\[
 \mathcal D f(A)=\sum_j g_j(A)(\mathcal DA)h_j(A),
 \qquad \ddiff{f}=\sum_jg_j\otimes h_j.
\]
\end{theorem}
\begin{proof}
Для внутреннего дифференцирования это коммутаторная
формула из лекции~1. В общем случае применяем её к
$[\mathcal D,f(L_A)]$, после чего используем предыдущую
лемму и равенство $f(L_A)=L_{f(A)}$:
\[
 \begin{aligned}
 L_{\mathcal D f(A)}
 &=[\mathcal D,L_{f(A)}]=[\mathcal D,f(L_A)]\\
 &=\overset{2}{[\mathcal D,L_A]}\,
      \ddiff{f}(\overset{1}{L_A},\overset{3}{L_A})\\
 &=\overset{2}{L_{\mathcal D A}}\,
      \ddiff{f}(\overset{1}{L_A},\overset{3}{L_A})\\
 &=L_{\vphantom{\mathcal D A}\smash[t]{\overset{2}{\mathcal D(A)}\,
      \ddiff{f}(\overset{1}{A},\overset{3}{A})}}.
 \end{aligned}
\]
Таким образом, \[
 \mathcal D f(A)=\overset{2}{\mathcal D(A)}\,
        \ddiff{f}(\overset{1}{A},\overset{3}{A}).
\]
Теперь докажем тензорную форму:
\[
 \begin{aligned}
 \mathcal D f(A)
 &=\overset{2}{\mathcal D(A)}\,
        \ddiff{f}(\overset{1}{A},\overset{3}{A})\\
 &=\sum_j h_j(A)\,\mathcal D(A)\,g_j(A)\\
 &=\sum_j g_j(A)\,\mathcal D(A)\,h_j(A),
 \qquad \ddiff{f}=\sum_j g_j\otimes h_j .
 \end{aligned}
\]
\end{proof}

\sessionsection{Лекция 3. Формула Ньютона и построение квазимод}
\record{19 сентября}{Ks_j6QaMPJ4}

\topic{Формулы Ньютона и Тейлора}{\at{Ks_j6QaMPJ4}{92}{1:32}}
\begin{definition}[Разностные производные высших порядков]
Полагаем $\ddiff[0]{f}=f$. Последующие функции определяются
рекурсивно:
\[
 \ddiff[k+1]{f}(y,x,x_1,\ldots,x_k)
 =\frac{\ddiff[k]{f}(y,x_1,\ldots,x_k)-
          \ddiff[k]{f}(x,x_1,\ldots,x_k)}{y-x}.
\]
На совпадающих аргументах используются непрерывные продолжения,
если они определены. В частности,
$\ddiff[k]{f}(x,\ldots,x)=f^{(k)}(x)/k!$.
\end{definition}

Для записи с доски обозначим $Q=C-A$. Если
$g(x_0,\ldots,x_k)=\prod_{j=0}^k g_j(x_j)$, положим
\[
 I_k(g;B_0,\ldots,B_k)
 =g_k(B_k)Q\,g_{k-1}(B_{k-1})Q\cdots Qg_0(B_0).
\]
Таким образом, индексы над $B_j$ нечётные,
$1,3,\ldots,2k+1$, а над $k$ копиями $Q$ чётные,
$2,4,\ldots,2k$. Для общих символов $I_k$ продолжается
линейно и непрерывно в условиях исчисления.

\begin{theorem}[Формула Ньютона с остатком]
Пусть $A,C$~--- генераторы, а все участвующие разностные
производные $f$ являются допустимыми символами. Тогда при
$N\ge1$
\[
 \begin{aligned}
 f(C)-f(A)
 &=\sum_{k=1}^{N-1}I_k(\ddiff[k]{f};A,\ldots,A)+R_N,\\
 R_N&=I_N(\ddiff[N]{f};C,A,\ldots,A).
 \end{aligned}
\]
\end{theorem}
\begin{proof}
При $N=1$ сумма пуста, и формула означает выделение
$C-A$ из $f(C)-f(A)$. Для перехода к $N+1$ сравним
$R_N$ с $I_N(\ddiff[N]{f};A,\ldots,A)$. Отличается только первый
слот $C$ или $A$, причём соответствующий множитель
действует первым. На уровне скалярных символов
\[
 \ddiff[N]{f}(c,x_1,\ldots,x_N)-\ddiff[N]{f}(a,x_1,\ldots,x_N)
   =(c-a)\ddiff[N+1]{f}(c,a,x_1,\ldots,x_N).
\]
Поставим над $C$ индекс $1$, над копией $A$ в первом
слоте второго слагаемого индекс $1.5$, а над ближайшим
старым $Q$ индекс $2$. Раздвигаем $1$ и $1.5$, оставляя
операторы разностной производной по разные стороны от
нового множителя. Линейный множитель $c-a$ становится
автономным $Q$, действующим между первой и второй
переменными $\ddiff[N+1]{f}$. Последующие $N$ множителей $Q$
сохраняют порядок. Перенумерация индексов даёт $R_{N+1}$.
При $N=1$ это в точности переход от
$R_1-I_1(\ddiff{f};A,A)$ к $R_2$, выписанный на доске.
\end{proof}

\begin{theorem}[Формула Тейлора с остатком]
Обозначим через $Q_k$ результат упорядоченной подстановки
$\overset1 C,\overset2 A$ в $(y-x)^k$. По биному
\[
 Q_k=\sum_{j=0}^k\binom{k}{j}(-A)^{k-j}C^j .
\]
Пусть $F_N(y,z,x)=z \ddiff[N]{f}(y,x,\ldots,x)$.
Тогда в тех же условиях
\[
 f(C)-f(A)=\sum_{k=1}^{N-1}\frac{f^{(k)}(A)}{k!}Q_k+
             F_N(\overset1 C,\overset2{Q_N},\overset3 A).
\]
\end{theorem}
\begin{proof}
Пишем скалярную формулу
\[
 f(y)-f(x)
  =\sum_{k=1}^{N-1}\frac{f^{(k)}(x)}{k!}(y-x)^k
   +(y-x)^N \ddiff[N]{f}(y,x,\ldots,x).
\]
Подстановка $y=\overset1 C$, $x=\overset2 A$
совпадает в обеих частях. Для каждого члена суммы
раздвигаем индекс при $A$, отделяя $f^{(k)}(A)$ от
автономной скобки $Q_k$ слева. В остатке раздвигаем
оба индекса. Копию $C$ из $\ddiff[N]{f}$ помещаем раньше
блока $(y-x)^N$, одну общую копию $A$ для диагональных
аргументов~--- позже. Блок остаётся между ними как
$Q_N$, что даёт указанную подстановку в $F_N$.
\end{proof}
При $C=A+\varepsilon B$ уже для $k=2$ имеем
\[
 C^2-2AC+A^2=\varepsilon[B,A]+\varepsilon^2B^2 .
\]
Для $k\ge1$ линейный по $\varepsilon$ член соответствует
$(k-1)$-кратному коммутатору справа с $A$. Если
$[[B,A],A]=0$, в первом порядке остаются только
$\varepsilon f'(A)B$ и
$\tfrac{\varepsilon}{2}f''(A)[B,A]$.
Для ВКБ-подстановки таким же образом используется то, что
коммутатор умножения на функцию и дифференцирования
есть оператор умножения. В формуле Ньютона каждый
множитель $Q$ равен $\varepsilon B$, поэтому $k$-й
член действительно имеет порядок $\varepsilon^k$.
У формулы Тейлора преимущество состоит в том, что
в явных членах появляются обычные производные
$f^{(k)}$, однако $Q_k$ может иметь уже порядок
$\varepsilon$ при $k>1$. Например,
$Q_3=\varepsilon[[B,A],A]+O(\varepsilon^2)$:
это следует из $Q_{k+1}=Q_kC-AQ_k$.

\topic{ВКБ-подстановка}{\at{Ks_j6QaMPJ4}{2250}{37:30}}
\begin{exercise}[Построение квазимоды]
Для гамильтониана $H(x,p)$ с квантуемым импульсом
$P=-\ii h\partial_x$ найти фазу $S$ и амплитуду $\varphi$
в выражении $u=\ee^{\ii S/h}\varphi$, чтобы
$\widehat Hu=O(h^2)$ в формальном разложении.
Гамильтониан квантуется с действием производных до
умножения. Если рассматривается уравнение на заданное
собственное значение, его переносят в гамильтониан.
\end{exercise}
Обозначим через $U$ умножение на $\ee^{\ii S/h}$.
В одномерном случае по правилу Лейбница
\[
 U^{-1}PUv
 =\ee^{-\ii S/h}(-\ii h)\partial_x(\ee^{\ii S/h}v)
 =-\ii h\,\partial_xv+S_xv .
\]
$U$ коммутирует с умножением на $x$, поэтому сопряжённый
символ получается подстановкой $S_x-\ii h\partial_x$
вместо $P$. Здесь $S_x$ действует умножением, поэтому
$[\partial_x,S_x]v=\partial_x(S_xv)-S_x\partial_xv=S_{xx}v$
и $[[\partial_x,S_x],S_x]=0$. По предшествующей формуле
\[
 U^{-1}\widehat H U
 =H(x,S_x)-\ii h\left(
    H_p(x,S_x)\partial_x+\tfrac12 H_{pp}(x,S_x)S_{xx}
    \right)+O(h^2).
\]
Здесь члены разложения действуют на амплитуду $\varphi$.
В нескольких переменных в первом слагаемом в скобках
появляется скалярное произведение с $\partial_x\varphi$,
а во втором~--- свёртка двух матриц. По отдельности
зануляем члены порядков $1$ и $h$. Получаем
уравнение Гамильтона--Якоби
\[
 H(x,S_x)=0.
\]
Член порядка $h$ даёт уравнение переноса
\[
 H_p(x,S_x)\cdot\partial_x\varphi+
 \tfrac12\tr\!\left(H_{pp}(x,S_x)S_{xx}\right)\varphi=0 .
\]
При выполнении обоих уравнений в рассматриваемом разложении
$\widehat H u$ имеет порядок $O(h^2)$.

\topic{Лагранжева поверхность и характеристики}{\at{Ks_j6QaMPJ4}{3963}{1:06:03}}
В фазовом пространстве $(x,p)\in\RR^{2n}$ обозначим
$\theta=\sum_jp_j\dd x_j$ и
$\omega=\sum_j\dd x_j\wedge\dd p_j=-\dd\theta$.
\begin{definition}
$n$-мерную подповерхность $\Lambda$ называют лагранжевой,
если $\omega|_\Lambda=0$. Гамильтоново векторное поле
$X_H$ определяется равенством $\omega(X_H,V)=\dd H(V)$
для всех векторов $V$. В координатах $X_H=(H_p,-H_x)$.
\end{definition}
\begin{lemma}
График $p=S_x$ лагранжев, а уравнение Гамильтона--Якоби
равносильно включению этого графика в $H=0$.
На лагранжевой поверхности, лежащей в $H=0$, поле $X_H$
касательно к поверхности.
\end{lemma}
\begin{proof}
На графике $\theta=\dd S$, поэтому
$\omega|_\Lambda=-\dd(\dd S)=0$. Проекция $\pi_x$ графика
обратима. Подстановка $p=S_x$ в $H=0$ даёт уравнение
Гамильтона--Якоби. Если $V$ касателен к $\Lambda$,
$\dd H(V)=0$, и значит $\omega(X_H,V)=0$. Пространство
$T\Lambda$ равно своему симплектическому ортогоналу,
поскольку оно изотропно и имеет половинную размерность.
Поэтому $X_H\in T\Lambda$.
\end{proof}
Для построения поверхности берут начальную
$(n-1)$-мерную поверхность в $H=0$, на которой начальная фаза
согласована с формой $p\dd x$, и выпускают из неё характеристики,
поперечные к ней. В координатах $(t,\alpha)$
\[
 \dot x=H_p,\qquad \dot p=-H_x,\qquad
 S(x(t,\alpha))=S^0(\alpha)+
       \int_0^t p(s,\alpha)\cdot\dot x(s,\alpha)\dd s .
\]
На начальной поверхности $x^0(\alpha),p^0(\alpha)$ требуется
$H(x^0,p^0)=0$ и $p^0\dd x^0=\dd S^0$.
Вдоль потока $\dot H=H_x\cdot H_p-H_p\cdot H_x=0$.
Гамильтонов поток сохраняет $\omega$, потому что
$\mathcal L_{X_H}\omega=\dd(\iota_{X_H}\omega)=\dd^2H=0$.
Поэтому попарные значения $\omega$ на перенесённых
начальных касательных векторах нулевые. На паре с
$X_H$ форма тоже нулевая в силу сохранения $H$.
Поперечность к начальной поверхности обеспечивает
размерность $n$ у полученной $\Lambda$. Если её
проекция на пространство $x$ локально обратима, то
$\dd\theta=0$ на $\Lambda$ и можно восстановить $S$
интегрированием $\theta$. Вдоль траектории
$\partial_t S=p\cdot\dot x$, что и даёт приведённую
формулу с начальным значением.

\begin{figure}[ht]
\centering
\begin{tikzpicture}[scale=.95,>=Stealth]
 \coordinate (o) at (0,0);
 \draw[->] (o)--(-1.5,-.95) node[below] {$x_1$};
 \draw[->] (o)--(2,-.05) node[right] {$x_n$};
 \draw[->] (o)--(-.4,1.45) node[above] {$p_1$};
 \draw[->] (o)--(.04,1.65) node[above] {$p_n$};
 \draw[thick] (-1.25,.48) .. controls (-.58,1.1) and (1.25,1.91) .. (2.2,1.46)
   .. controls (2.3,1.3) and (2.08,1.07) .. (2.19,.95)
   .. controls (1.1,.95) and (-.36,.48) .. cycle;
 \draw (.22,.69) -- (.57,1.17) -- (1.16,1.23) -- (.92,.77) -- cycle;
 \fill (.68,.93) circle (1pt);
 \node[below] at (.77,.91) {$x^0$};
 \node at (1.45,1.73) {$p=\partial S/\partial x$};
\end{tikzpicture}
\caption{График градиента фазы в фазовом пространстве.
По доске в \at{Ks_j6QaMPJ4}{4035}{1:07:15}.}
\end{figure}

\topic{Якобиан, амплитуда и фокальная точка}{\at{Ks_j6QaMPJ4}{6900}{1:55:00}}
Пусть $M=\partial x/\partial(t,\alpha)$, $J=\det M$.
После нахождения $S$ на участке с обратимой проекцией
уравнение для $x$ имеет вид $\dot x=F(x)$, где
$F(x)=H_p(x,S_x(x))$. Каждый столбец $M$
удовлетворяет линеаризованному уравнению, поэтому
$\dot M=(\partial_x F)M$. Формула дифференцирования
определителя по элементам матрицы даёт
\[
 \frac{\dot J}{J}=
 \frac{1}{\det M}\sum_{j,k}\operatorname{cof}_{jk}(M)\dot M_{jk}
 =\tr(M^{-1}\dot M)=\tr(\partial_xF).
\]
В последнем переходе используется цикличность следа.
По правилу цепочки
$\tr(\partial_xF)=\tr H_{px}+\tr(H_{pp}S_{xx})$.
Производная амплитуды вдоль характеристики равна
$H_p\cdot\partial_x\varphi$. Уравнение переноса принимает
вид $\dot\varphi+\tfrac12(\dot J/J-\tr H_{px})\varphi=0$.
Умножение на $\sqrt{|J|}$ и свёртывание двух членов
в производную произведения дают
\[
 \begin{aligned}
 \partial_t(\sqrt{\abs J}\,\varphi)
     &=\tfrac12\tr H_{px}\sqrt{\abs J}\,\varphi,\\
 \varphi(t,\alpha)&=\varphi(0,\alpha)
    \sqrt{\frac{\abs{J(0,\alpha)}}{\abs{J(t,\alpha)}}}
    \exp\!\left(\tfrac12\int_0^t\tr H_{px}\dd s\right).
 \end{aligned}
\]
В фокальной точке $J=0$ и проекция на $x$ перестаёт быть
локально обратимой. На рисунке одна точка $x^0$ имеет
несколько прообразов с различными импульсами. Тогда
представление единственной фазой $S(x)$ перестаёт описывать
всю поверхность около такой точки.
\begin{lemma}[Лемма Арнольда, локальные координаты]
Для лагранжевой поверхности в окрестности каждой её
точки из координат $x_j,p_j$ можно выбрать локальные
координаты вида $(x_I,p_J)$, где $I$ и $J$~---
дополнительные множества индексов.
\end{lemma}
Когда вырождается проекция $\Lambda$ на $x$, переходят к
проекции на импульсы $p$ или на смешанную координатную
плоскость. При $n=1$ это выбор между $x$ и $p$, а при
$n=2$ возможны, например, пары $(x_1,p_2)$ и $(x_2,p_1)$.
Замене координаты $x_j$ на $p_j$ соответствует
преобразование Фурье по $x_j$. Лемма обеспечивает
существование такого локального выбора
\at{Ks_j6QaMPJ4}{8686}{2:24:46}.

\begin{figure}[ht]
\centering
\begin{tikzpicture}[>=Stealth,scale=.9]
 \begin{scope}
  \draw[->] (0,0)--(-1.15,-.8) node[below] {$x_1$};
  \draw[->] (0,0)--(2.4,0) node[right] {$x_n$};
  \draw[->] (0,0)--(-.38,1.85) node[above] {$p_1$};
  \draw[->] (0,0)--(.04,2) node[above] {$p_n$};
  \draw[thick] (.08,.34) .. controls (.24,1.1) and (.15,1.72) .. (.43,1.95);
  \draw[->] (.16,.55) .. controls (.7,.53) and (1.52,1.16) .. (2.08,.82);
  \draw[->] (.24,.86) .. controls (.86,1.04) and (1.67,1.61) .. (2.18,1.08);
  \draw[->] (.27,1.22) .. controls (.9,1.8) and (1.69,1.97) .. (2.18,1.35);
  \draw (.43,1.95) .. controls (1.25,2.2) and (1.82,2.62) .. (2.18,1.35)
   .. controls (2.37,1.12) and (2.48,.88) .. (2.51,.67)
   .. controls (1.42,.89) and (.57,.32) .. (.08,.34);
  \node at (1.39,2.15) {$\Lambda^n$};
  \node at (1.62,.42) {$\Lambda^{n-1}$};
 \end{scope}
 \begin{scope}[xshift=5.2cm]
  \draw[->] (-.8,0)--(2.6,0) node[right] {$x$};
  \draw[->] (0,-.3)--(0,1.95) node[above] {$p$};
  \draw[densely dashed] (.75,0)--(.75,1.6);
  \draw[thick] (.2,.43) .. controls (2.56,.4) and (2.88,1.7) .. (.3,1.6);
  \fill (.75,.45) circle (1.2pt) (.75,1.59) circle (1.2pt);
  \node[below] at (.75,0) {$x^0$};
 \end{scope}
\end{tikzpicture}
\caption{Характеристики на лагранжевой поверхности и потеря
однозначности проекции. По доске в
\at{Ks_j6QaMPJ4}{8560}{2:22:40}.}
\end{figure}

\sessionsection{Лекция 4. Графики, замыкание и сопряжение}
\record{25 сентября}{HFKXR_j7aA4}

\topic{Линеал, область определения и ограниченность}{\at{HFKXR_j7aA4}{122}{2:02}}
\begin{definition}
Линеал~--- линейное подмножество гильбертова пространства
$\Hh$. Замкнутый линеал называют подпространством.
Линейный оператор $T\colon\Dom(T)\to\Hh$ задан на линеале
$\Dom(T)\subset\Hh$. Его образ обозначаем $\operatorname{Im}T$,
ядро~--- $\ker T$. Если $\overline{\Dom(T)}=\Hh$, оператор
называют плотно определённым.
\end{definition}
Если $\norm{Tx}\le c\norm x$ на $\Dom(T)$, оператор однозначно
продолжается на $\overline{\Dom(T)}$ по непрерывности. Для
$x_n\in\Dom(T)$, $x_n\to x$, та же оценка даёт фундаментальность
$Tx_n$ и независимость предела от выбора $x_n$. Вне замыкания
однозначности нет. Поэтому далее рассматриваем плотно определённые
операторы. При наличии оценки продолжение принадлежит $\mathcal B(\Hh)$
(\at{HFKXR_j7aA4}{342}{5:42}).

\begin{definition}[Норма, ассоциированная с оператором]
На $\Dom(T)$ вводим скалярное произведение
$(x,y)_T=(x,y)+(Tx,Ty)$ и соответствующую норму
$\norm x_T=(\norm x^2+\norm{Tx}^2)^{1/2}$.
\end{definition}
Из
$\sqrt{a^2+b^2}\le a+b\le\sqrt 2\sqrt{a^2+b^2}$
следует эквивалентность $\norm x_T$ и $\norm x+\norm{Tx}$.
Непрерывность $T$ в исходной норме равносильна эквивалентности
$\norm x$ + $\norm{Tx}$ и $\norm x$. Действительно, из $\norm{Tx}\le c\norm x$
получаем $\norm x\le\norm x + \norm{Tx}\le(c+1)\norm x$, а из
$\norm x + \norm {Tx}\le b\norm x$ следует $\norm{Tx}\le (b-1)\norm x$.

\begin{lemma}[Ограниченный обратный на образе]
Обратный оператор $T^{-1}\colon\operatorname{Im}T\to\Dom(T)$
существует тогда и только тогда, когда $\ker T=\{0\}$.
Он существует и ограничен тогда и только тогда, когда
$\norm{Tx}\ge c\norm x$ для всех $x\in\Dom(T)$ при некотором $c>0$.
\end{lemma}
\begin{proof}
Равенство $\ker T=\{0\}$ обеспечивает единственность прообраза
$y=Tx$. Подстановка $y=Tx$ в оценку $\norm{T^{-1}y}\le M\norm y$
даёт $\norm x\le M\norm{Tx}$. Обратно, оценка снизу влечёт
$\ker T=\{0\}$ и $\norm{T^{-1}y}\le c^{-1}\norm y$ при $y=Tx$.
\end{proof}

\topic{Действия с операторами и их графики}{\at{HFKXR_j7aA4}{1590}{26:30}}
$T_1=T_2$, если области определения совпадают и $T_1x=T_2x$
на общей области. Если $\Dom(T_1)\subset\Dom(T_2)$ и $T_1x=T_2x$
на $\Dom(T_1)$, пишем $T_1\subset T_2$ и называем $T_2$ расширением
$T_1$, а $T_1$~--- сужением $T_2$. В $L^2(\RR)$ дифференциальное
выражение $\mathrm d/\mathrm dx$ задаёт разные операторы на
$\sch{\RR}$ и $\mathcal S(\RR)$, причём первый есть сужение второго.
Для суммы и произведения
\[
 \Dom(T+S)=\Dom(T)\cap\Dom(S),\qquad
 \Dom(TS)=\{x\in\Dom(S): Sx\in\Dom(T)\}.
\]
Перестановочность $TS=ST$ требует совпадения и областей, и действий.

\topic{Замкнутый график}{\at{HFKXR_j7aA4}{2120}{35:20}}
На вещественной прямой рассмотрим функцию
\[
 f(x)=\begin{cases}0,&x\le0,\\ 1/x,&x>0.\end{cases}
\]
Она разрывна в нуле. Если $x_n\to0$ и $f(x_n)$ сходится,
среди $x_n$ лишь конечное число положительных членов, иначе
по ним $f(x_n)\to+\infty$. Поэтому $f(x_n)\to0=f(0)$.
В остальных точках это следует из непрерывности. Так получается
замкнутый график разрывной функции.

\begin{center}
\begin{minipage}{\linewidth}
\centering
\begin{tikzpicture}[>=Stealth,scale=.85]
 \draw[->] (-2.2,0)--(3,0) node[right] {$x$};
 \draw[->] (0,-.65)--(0,2.38);
 \draw[thick] (-2.05,0)--(0,0);
 \fill (0,0) circle (1.5pt);
 \draw[thick,domain=.45:2.72,samples=50] plot (\x,{1/\x});
\end{tikzpicture}
\captionof{figure}{Замкнутый график разрывной функции
(\at{HFKXR_j7aA4}{2192}{36:32}).}
\end{minipage}
\end{center}

\begin{definition}[График оператора]
$G(T)=\{(x,Tx):x\in\Dom(T)\}$~--- график $T$. Это линеал в
$\Hh\oplus\Hh$.
\end{definition}

\begin{definition}[Замкнутый оператор]
Оператор $T$ замкнут, если $\Dom(T)$ полно
в норме $\norm{\cdot}_T$, ассоциированной с $T$.
\end{definition}
Включения $T_1\subset T_2$ и $G(T_1)\subset G(T_2)$ равносильны.
Отображение $x\mapsto(x,Tx)$~--- линейный изоморфизм
$\Dom(T)$ на $G(T)$. Скалярное произведение прямой суммы
на графике совпадает с $(x,y)_T$.
\begin{lemma}[Критерии замкнутости]
Оператор $T$ замкнут тогда и только тогда, когда выполнено
любое из следующих условий.
\begin{enumerate}
\item $G(T)$ замкнут в $\Hh\oplus\Hh$.
\item Если $x_n\in\Dom(T)$, $x_n\to x$ и $Tx_n\to y$,
то $x\in\Dom(T)$ и $Tx=y$.
\end{enumerate}
\end{lemma}
\begin{proof}
Изоморфизм $x\mapsto(x,Tx)$ сохраняет норму $\norm{\cdot}_T$, поэтому
полнота $\Dom(T)$ равносильна полноте $G(T)$. В полном пространстве
$\Hh\oplus\Hh$ линеал полон тогда и только тогда, когда он замкнут.
Условие~2 выражает замкнутость графика через сходимость пар
$(x_n,Tx_n)\to(x,y)$.
\end{proof}

\begin{lemma}[Ограниченное возмущение]
Если $T$ замкнут и $S\in\mathcal B(\Hh)$, то $T+S$
с областью $\Dom(T)$ замкнут.
\end{lemma}
\begin{proof}
Область определения та же, поскольку $S$ задан на всём $\Hh$.
Сравним нормы $\norm{\cdot}_T$ и $\norm{\cdot}_{T+S}$ по неравенству
треугольника.
\[
 \begin{aligned}
 \norm x+\norm{(T+S)x}
    &\le(1+\norm S)\norm x+\norm{Tx},\\
 \norm x+\norm{Tx}
    &\le(1+\norm S)\norm x+\norm{(T+S)x}.
 \end{aligned}
\]
Эти нормы эквивалентны, поэтому полнота $\Dom(T)$ в норме
$\norm{\cdot}_T$ даёт полноту в норме $\norm{\cdot}_{T+S}$.
\end{proof}

\begin{lemma}[Ядро замкнутого оператора]
Если $T$ замкнут, то $\ker T$~--- замкнутый линеал
(\at{HFKXR_j7aA4}{3750}{1:02:30}).
\end{lemma}
\begin{proof}
Пусть $x_n\in\ker T$ и $x_n\to x$. Тогда $Tx_n=0\to0$,
поэтому из замкнутости следует $x\in\Dom(T)$ и $Tx=0$.
\end{proof}

\topic{Замыкаемость}{\at{HFKXR_j7aA4}{3870}{1:04:30}}
Если $T_1$~--- замкнутое расширение $T$, то
$G(T)\subset G(T_1)$ и $\overline{G(T)}\subset G(T_1)$.
Поэтому в $\overline{G(T)}$ нет точек $(0,y)$ с $y\ne0$.
\begin{lemma}[График как линейное подмножество]
Линеал $M\subset\Hh\oplus\Hh$ является графиком линейного оператора
тогда и только тогда, когда $M\cap(\{0\}\oplus\Hh)=\{0\}$.
\end{lemma}
\begin{center}
\begin{minipage}{\linewidth}
\centering
\begin{tikzpicture}[>=Stealth]
 \draw[->] (-.78,0)--(2.4,0) node[right] {$\Hh$};
 \draw[->] (0,-.47)--(0,1.87) node[above] {$\Hh$};
 \draw (-.63,-.48)--(2.2,1.68) node[above] {$M$};
 \draw[->] (1.53,1.17)--(0,1.17) node[midway,above] {$\pi_2$};
 \draw[->] (1.53,1.17)--(1.53,0) node[midway,right] {$\pi_1$};
 \fill (1.53,1.17) circle (1.4pt);
\end{tikzpicture}
\captionof{figure}{Проекции линеала $M$ на слагаемые прямой суммы
(\at{HFKXR_j7aA4}{4246}{1:10:46}).}
\end{minipage}
\end{center}
\begin{proof}
Если $M$~--- график, то $(0,y)\in M$ влечёт $y=T0=0$.
Обратно, условие означает, что проекция $\pi_1$ на первое
слагаемое инъективна на $M$. Поэтому
$\pi_1\colon M\to\pi_1(M)$ обратима и
$T=\pi_2\circ\pi_1^{-1}$ на $\pi_1(M)$. По построению $G(T)=M$.
\end{proof}

\begin{definition}[Замыкаемый оператор]
$T$ замыкаем, если $\overline{G(T)}$ является графиком оператора.
Этот оператор обозначают $\overline T$ и называют замыканием $T$.
\end{definition}
Итак, $T$ замыкаем тогда и только тогда, когда у него есть
замкнутое расширение. Замыкание $\overline T$ содержится в каждом
таком расширении. По лемме условие замыкаемости равносильно
отсутствию в $\overline{G(T)}$ точек $(0,y)$ с $y\ne0$.

\topic{Сопряжённый оператор}{\at{HFKXR_j7aA4}{4539}{1:15:39}}
\begin{definition}[Сопряжённый оператор]
Пусть $\Dom(T)$ плотна в $\Hh$. Для $y\in\Hh$ рассмотрим
функционал $\ell_y(x)=(Tx,y)$ на $\Dom(T)$. В $\Dom(T^*)$
входят те $y$, для которых найдётся $c_y$ с оценкой
$\abs{\ell_y(x)}\le c_y\norm x$ при всех $x\in\Dom(T)$.
По плотности $\ell_y$ однозначно
продолжается на $\Hh$. По теореме Рисса существует единственный
$h\in\Hh$ такой, что $\ell_y(x)=(x,h)$ для всех $x\in\Dom(T)$.
Полагаем $T^*y=h$.
\end{definition}
Область $\Dom(T^*)$ всегда содержит нуль, а её плотность
зависит от $T$. Плотность исходной области $\Dom(T)$ нужна для
однозначности $h$.
\Needspace{12\baselineskip}
\begin{lemma}[График сопряжённого]
\label{lem:adjoint-graph}
Пусть $T$ плотно определён. Для поворота $W$ на $90^{\circ}$
в $\Hh\oplus\Hh$ имеем
\[
 W\binom{x}{y}=\binom{-y}{x},\qquad
 G(T^*)=(W(G(T)))^\perp=W(G(T)^\perp).
\]
\end{lemma}
\begin{proof}
\leavevmode
\begin{enumerate}
\item $G(T^*)=(W(G(T)))^\perp$.

Вектор $\binom{y}{h}$ принадлежит $G(T^*)$ тогда и только тогда,
когда $(Tx,y)=(x,h)$ для всех $x\in\Dom(T)$. Это равносильно
равенству
\[
  \begin{aligned}
  0=-(Tx,y)+(x,h)&=\left(\binom{-Tx}{x},\binom{y}{h}\right) \\ &=\left(W\binom{x}{Tx},\binom{y}{h}\right)
  \end{aligned}
 \qquad\text{для всех }x\in\Dom(T).
\]
Таким образом, $\binom{y}{h}\in G(T^*)$ тогда и только тогда, когда он ортогонален всем векторам из $W(G(T))$.

\item $(W(G(T)))^\perp=W(G(T)^\perp)$.
Следует из унитарности $W$.
\end{enumerate}
\end{proof}
\begin{corollary}
\label{cor:adjoint-is-closed}
$T^*$ замкнут.
\end{corollary}
\begin{proof}
Ортогональное дополнение в прямой сумме замкнуто.
\end{proof}
\Needspace{5\baselineskip}
\begin{corollary}
\label{cor:closure-of-graph}
$\overline{G(T)}=(WG(T^*))^\perp$.
\end{corollary}
\begin{proof}
$WG(T^*)=W^2(G(T)^\perp)=G(T)^\perp$, поскольку $W^2=-\id$.
Берём ортогональное дополнение и используем
$(G(T)^\perp)^\perp=\overline{G(T)}$.
\end{proof}
\begin{theorem}[Критерий замыкаемости через сопряжение]
\label{thm:closability-via-adjoint}
Плотно определённый $T$ замыкаем тогда и только тогда,
когда $\Dom(T^*)$ плотна. В этом случае
$T^{**}=\overline T$.
\end{theorem}
\begin{proof}
Если $\Dom(T^*)$ плотна, определён оператор $T^{**}$.
По \mbox{лемме~\ref{lem:adjoint-graph}}
и \mbox{следствию~\ref{cor:closure-of-graph}}
имеем $G(T^{**})=(WG(T^*))^\perp=\overline{G(T)}$.
Значит, $T$ замыкаем и $\overline T=T^{**}$.

Если $\Dom(T^*)$ не плотна, берём ненулевой $y$ из
$\Dom(T^*)^\perp$. Для каждого $z\in\Dom(T^*)$ имеем
$\bigl((-T^*z,z),(0,y)\bigr)=(z,y)=0$. Следовательно,
$(0,y)\in(WG(T^*))^\perp=\overline{G(T)}$, и $T$ незамыкаем.
\end{proof}

\sessionsection{Лекция 5. Поле регулярности, дефект и резольвента}
\record{2 октября}{dRCTK4VqBTk}

Напомним теорему~\ref{thm:closability-via-adjoint}. Плотно определённый
оператор $T$ замыкаем тогда и только тогда, когда $T^*$ плотно
определён. В этом случае $\overline T=T^{**}$.

\topic{Теорема о замкнутом графике}{\at{dRCTK4VqBTk}{60}{1:00}}
\begin{theorem}[О замкнутом графике]
\label{thm:closed-graph}
Если линейный оператор $T$ в гильбертовом пространстве
$\Hh$ замкнут и $\Dom(T)=\Hh$, то $T$ ограничен.
\end{theorem}
\begin{proof}
\leavevmode
\begin{enumerate}
\item Докажем, что $T^*$ ограничен и $\Dom(T^*)=\Hh$.
Для $y\in\Dom(T^*)$ с $\norm y=1$ рассмотрим функционал
$\ell_y(x)=(Tx,y)=(x,T^*y)$ на $\Hh$. Для каждого $x\in\Hh$
\[
 \abs{\ell_y(x)}\le\norm{Tx}\norm y=\norm{Tx}.
\]
По теореме Банаха~--- Штейнгауза из поточечной ограниченности
семейства ограниченных линейных операторов на банаховом пространстве
следует равномерная ограниченность их норм
(\at{dRCTK4VqBTk}{315}{5:15}). Здесь $\norm{\ell_y}=\norm{T^*y}$,
поэтому $\norm{T^*y}\le c\norm y$ для всех $y\in\Dom(T^*)$.

По следствию~\ref{cor:adjoint-is-closed} оператор $T^*$ замкнут.
Поэтому $\Dom(T^*)$ полна по норме, ассоциированной с $T^*$.
Эквивалентная ей норма суммы удовлетворяет оценкам
\[
 \norm y\le \norm y+\norm{T^*y}\le (1+c)\norm y .
\]
Значит, $\Dom(T^*)$ полна в норме $\Hh$ и потому замкнута.
По теореме~\ref{thm:closability-via-adjoint} она также плотна,
откуда $\Dom(T^*)=\Hh$.

\item По первому пункту $\abs{(T^*x,y)}\le c\norm x\norm y$
для всех $x,y\in\Hh$. Следовательно, $\Dom(T^{**})=\Hh$.
Подставляя $x=T^{**}y$ в $(T^*x,y)=(x,T^{**}y)$, получаем
\[
 \norm{T^{**}y}^2=(T^*T^{**}y,y)\le c\norm{T^{**}y}\norm y .
\]
Значит, $\norm{T^{**}y}\le c\norm y$ и $T=T^{**}$ ограничен.
\end{enumerate}
\end{proof}
В конце отмечают, что второй пункт можно получить из первого,
применив его к замкнутому всюду определённому оператору $T^*$
(\at{dRCTK4VqBTk}{1441}{24:01}).

\topic{Точки регулярности и дефект}{\at{dRCTK4VqBTk}{1680}{28:00}}
\begin{definition}[Поле регулярности]
Поле, или множество, регулярности $\Pi(T)$ состоит из тех $\lambda\in\CC$,
для которых существует $c_\lambda>0$ с
\[
 \norm{(T-\lambda\id)x}\ge c_\lambda\norm x,\qquad
 x\in\Dom(T).
\]
\end{definition}
\begin{definition}[Дефектное пространство и дефектное число]
Дефектное пространство $\mathcal N_\lambda(T)$ есть
$\operatorname{Im}(T-\lambda\id)^\perp$. Его размерность
$d_\lambda(T)$ называется дефектным числом. Оно может быть
бесконечным (\at{dRCTK4VqBTk}{1811}{30:11}).
\end{definition}

\begin{lemma}[Обратный в точке регулярности]
\label{lem:regular-point-inverse}
$\lambda\in\Pi(T)$ тогда и только тогда, когда
$(T-\lambda\id)^{-1}$ существует и ограничен на
$\operatorname{Im}(T-\lambda\id)$.
\end{lemma}
\begin{proof}
Оценка снизу влечёт инъективность $T-\lambda\id$.
При $y=(T-\lambda\id)x$ она равносильна оценке
$\norm{(T-\lambda\id)^{-1}y}\le c_\lambda^{-1}\norm y$.
\end{proof}

\begin{lemma}[Открытость поля регулярности]
\label{lem:regularity-open}
$\Pi(T)$ открыто.
\end{lemma}
\begin{proof}
Для $|\mu-\lambda|<c_\lambda$ имеем
\[
 \norm{(T-\mu\id)x}
 \ge\norm{(T-\lambda\id)x}-|\mu-\lambda|\norm x
 \ge(c_\lambda-|\mu-\lambda|)\norm x.
\]
Значит, весь указанный круг лежит в $\Pi(T)$.
\end{proof}

\begin{lemma}[Поле регулярности при замыкании оператора]
\label{lem:regularity-under-closure}
\at{dRCTK4VqBTk}{2481}{41:21}.
Если $T$ замыкаем, то $\Pi(\overline T)=\Pi(T)$. Для
$\lambda\in\Pi(T)$ выполняются равенства
\[
 \operatorname{Im}(\overline T-\lambda\id)
   =\overline{\operatorname{Im}(T-\lambda\id)},
 \qquad d_\lambda(\overline T)=d_\lambda(T).
\]
\end{lemma}
\begin{proof}
\leavevmode
\begin{enumerate}
\item $\Pi(\overline T)=\Pi(T)$.

Если $\lambda\in\Pi(\overline T)$, то оценка для
$\overline T-\lambda\id$ верна и на $\Dom(T)$, где $\overline T x=Tx$.
Поэтому $\lambda\in\Pi(T)$.

Обратно, пусть $\lambda\in\Pi(T)$ и $x\in\Dom(\overline T)$.
По определению замыкания существуют $x_n\in\Dom(T)$ такие, что
$x_n\to x$ и $Tx_n\to\overline T x$. Переходя к пределу в оценке
$\norm{(T-\lambda\id)x_n}\ge c_\lambda\norm{x_n}$, получаем
$\norm{(\overline T-\lambda\id)x}\ge c_\lambda\norm x$.
Значит, $\lambda\in\Pi(\overline T)$.

\Needspace{7\baselineskip}
\item Равенство образов и дефектных чисел при $\lambda\in\Pi(T)$.

Пусть $y\in\overline{\operatorname{Im}(T-\lambda\id)}$. Возьмём
$x_n\in\Dom(T)$ так, что $y_n=(T-\lambda\id)x_n\to y$.
\[
 \norm{x_n-x_m}\le c_\lambda^{-1}\norm{y_n-y_m}.
\]
Значит, $x_n$ фундаментальна и сходится к некоторому $x\in\Hh$.
Пары $(x_n,y_n)$ лежат в $G(T-\lambda\id)$ и сходятся к $(x,y)$.
Поэтому $(x,y)\in\overline{G(T-\lambda\id)}=G(\overline T-\lambda\id)$
(\at{dRCTK4VqBTk}{3153}{52:33}). Равенство графиков следует из того,
что при $x_n\to x$ сходимость $(T-\lambda\id)x_n\to y$ равносильна
$Tx_n\to y+\lambda x$. В частности, $y=(\overline T-\lambda\id)x$.

В обратную сторону для $y=(\overline T-\lambda\id)x$ та же запись
замыкания графика даёт $y_n=(T-\lambda\id)x_n\to y$.
Значит, $y\in\overline{\operatorname{Im}(T-\lambda\id)}$.
Дефектные числа равны, поскольку замыкание линеала сохраняет его
ортогональное дополнение.
\end{enumerate}
\end{proof}

\Needspace{5\baselineskip}
\begin{lemma}[Замкнутость образа в точке регулярности]
\label{lem:regular-range-closed}
Если $T$ замкнут и $\lambda\in\Pi(T)$, то
$\operatorname{Im}(T-\lambda\id)$ замкнут.
\end{lemma}
\begin{proof}
Подставляем $T=\overline T$ в лемму~\ref{lem:regularity-under-closure}.
\end{proof}

\begin{corollary}[Ортогональное разложение в точке регулярности]
\label{cor:regularity-decomposition}
Для замкнутого плотно определённого $T$ и $\lambda\in\Pi(T)$
\[
 \Hh=\operatorname{Im}(T-\lambda\id)\oplus\ker(T^*-\bar\lambda\id),
\]
где слагаемые ортогональны.
\end{corollary}
\begin{proof}
Для любого плотно определённого $A$ условие
$y\perp\operatorname{Im}A$ означает, что $(Ax,y)=0=(x,0)$ для
всех $x\in\Dom(A)$. По определению это равносильно $y\in\ker A^*$.
Значит, $\Hh=\overline{\operatorname{Im}A}\oplus\ker A^*$.
Подставляем $A=T-\lambda\id$. Тогда $A^*=T^*-\bar\lambda\id$,
поэтому $\mathcal N_\lambda(T)=\ker(T^*-\bar\lambda\id)$.
В точке регулярности образ $A$ замкнут по
лемме~\ref{lem:regular-range-closed}.
\end{proof}
Множество замкнутых линейных операторов в $\Hh$ обозначают
$\mathcal C(\Hh)$ (\at{dRCTK4VqBTk}{3557}{59:17}).

\begin{lemma}[Сравнение размерностей]
\label{lem:orthogonal-vector-in-larger-space}
Пусть $F,G$~--- замкнутые подпространства сепарабельного
гильбертова пространства и $\dim F<\dim G$. Тогда
$G\cap F^\perp\ne\{0\}$.
\end{lemma}
\begin{proof}
$G$ может быть бесконечномерным, а $F$ конечномерно.
Ограничим ортогональный проектор на $F$ на подпространство $G$.
Он действует из пространства большей размерности в пространство
меньшей размерности, поэтому его ядро $G\cap F^\perp$ ненулевое.
\end{proof}

\begin{theorem}[Крейна--Красносельского]
\label{thm:krein-krasnoselskii}
\at{dRCTK4VqBTk}{4483}{1:14:43}.
Для замкнутого плотно определённого $T$ в сепарабельном
гильбертовом пространстве $d_\lambda(T)$ постоянно на
каждой связной компоненте $\Pi(T)$.
\end{theorem}
\begin{proof}
Пусть $\lambda,\mu\in\Pi(T)$ и $d_\lambda(T)<d_\mu(T)$.
По лемме~\ref{lem:orthogonal-vector-in-larger-space} существует
$\xi\ne0$ из $\mathcal N_\mu(T)$, ортогональный
$\mathcal N_\lambda(T)$. По следствию~\ref{cor:regularity-decomposition}
имеем $\xi=(T-\lambda\id)u$ для некоторого $u\in\Dom(T)$.
Ортогональность $\xi$ к $\operatorname{Im}(T-\mu\id)$ означает, что
\[
 \bigl((T-\lambda\id)u,(T-\mu\id)x\bigr)=0
 \qquad\text{для всех } x\in\Dom(T).
\]
Берём $x=u$ и пишем $(T-\mu\id)u=\xi+(\lambda-\mu)u$.
Так как скалярное произведение линейно по первому аргументу,
\[
 \norm\xi^2=(\bar\mu-\bar\lambda)(\xi,u),\qquad
 \norm\xi^2\le\abs{\mu-\lambda}\norm\xi\norm u
     \le\frac{\abs{\mu-\lambda}}{c_\lambda}\norm\xi^2 .
\]
При $\abs{\mu-\lambda}<c_\lambda$ получаем противоречие.
Если $d_\mu(T)<d_\lambda(T)$, меняем $\lambda$ и $\mu$ местами.
Чтобы обе оценки работали в одной окрестности $\lambda$, достаточно
взять $\abs{\mu-\lambda}<c_\lambda/2$. Тогда из доказательства
леммы~\ref{lem:regularity-open} можно взять
$c_\mu=c_\lambda-\abs{\mu-\lambda}>c_\lambda/2$.
Итак, $d_\lambda(T)$ локально постоянно на $\Pi(T)$, а значит,
постоянно на каждой её компоненте связности.
\end{proof}

\topic{Числовая область значений и спектр}{\at{dRCTK4VqBTk}{5837}{1:37:17}}
\begin{definition}
Числовая область значений $\Theta(T)$ состоит из $(Tx,x)$
при $x\in\Dom(T)$, $\norm x=1$.
\end{definition}
\begin{lemma}[Оценка через числовую область значений]
\label{lem:numerical-range-regularity}
Для $T\in\mathcal C(\Hh)$ выполнено
$\CC\setminus\overline{\Theta(T)}\subset\Pi(T)$.
\end{lemma}
\begin{proof}
Пусть $\delta=\dist(\lambda,\Theta(T))>0$. При $\norm x=1$
\[
 \norm{(T-\lambda\id)x}\ge\abs{((T-\lambda\id)x,x)}
     =\abs{(Tx,x)-\lambda}\ge\delta .
\]
По однородности $\norm{(T-\lambda\id)x}\ge\delta\norm x$
для всех $x\in\Dom(T)$.
\end{proof}

\begin{example}[Симметрический оператор]
Если $(Tx,x)=(x,Tx)$ для всех $x\in\Dom(T)$, то
$\Theta(T)\subset\RR$. По лемме~\ref{lem:numerical-range-regularity}
обе открытые полуплоскости лежат в $\Pi(T)$.
На каждой из них дефект постоянен. Эти значения называются
индексами дефекта:
\[
 d_+(T)=d_{\ii}(T),\qquad d_-(T)=d_{-\ii}(T).
\]
\end{example}

\begin{example}[Полуограниченный снизу оператор]
Если $(Tx,x)\ge\alpha\norm x^2$ для всех $x\in\Dom(T)$, то
$\Theta(T)\subset[\alpha,+\infty)$ и
$\CC\setminus[\alpha,+\infty)\subset\Pi(T)$.
В плоскости с вырезанным лучом обе полуплоскости лежат в одной
компоненте связности, поэтому $d_+(T)=d_-(T)$.
\end{example}

\begin{center}
\begin{minipage}{\linewidth}
\centering
\begin{tikzpicture}[>=Stealth]
 \draw[->] (-2,0)--(2.25,0);
 \draw[->] (0,-.63)--(0,.9);
 \draw[very thick] (-.92,0)--(2.1,0);
 \fill (-.92,0) circle (1.4pt) node[above left] {$\alpha$};
 \node[above] at (1.03,.03) {$\Theta(T)$};
\end{tikzpicture}
\captionof{figure}{Числовая область значений лежит на луче
$[\alpha,+\infty)$ (\at{dRCTK4VqBTk}{6506}{1:48:26}).}
\end{minipage}
\end{center}

\topic{Резольвентное множество}{\at{dRCTK4VqBTk}{6672}{1:51:12}}
\begin{definition}[Резольвента и спектр]
Пусть $T$~--- замкнутый оператор. Точка $\lambda$ лежит
в резольвентном множестве $\rho(T)$, если обратный оператор
$(T-\lambda\id)^{-1}$ существует и определён на всей $\Hh$.
В этом случае его называют резольвентой $R_\lambda(T)$.
Дополнение $\CC\setminus\rho(T)$ есть спектр $\sigma(T)$.
\end{definition}
График обратного получается из $G(T-\lambda\id)$ перестановкой
координат. Он замкнут, поэтому по теореме~\ref{thm:closed-graph}
резольвента ограничена. Две композиции имеют разные области определения:
\[
 (T-\lambda\id)R_\lambda(T)x=x\quad (x\in\Hh),\qquad
 R_\lambda(T)(T-\lambda\id)x=x\quad (x\in\Dom(T)).
\]

\begin{corollary}
Для замкнутого плотно определённого $T$ справедливы два свойства.
\begin{enumerate}
\item $\rho(T)=\{\lambda\in\Pi(T):d_\lambda(T)=0\}$.
\item $\rho(T)$ открыто, а $\sigma(T)$ замкнут.
\end{enumerate}
\end{corollary}
\begin{proof}
В точке $\lambda\in\Pi(T)$ обратный оператор существует, ограничен
и задан на замкнутом подпространстве $\operatorname{Im}(T-\lambda\id)$.
Оно совпадает с $\Hh$ тогда и только тогда, когда $d_\lambda(T)=0$.
По теореме~\ref{thm:krein-krasnoselskii} множество $\rho(T)$
есть объединение компонент связности $\Pi(T)$, на которых
$d_\lambda(T)=0$. Компоненты открытого множества $\Pi(T)\subset\CC$
открыты, поэтому $\rho(T)$ открыто.
\end{proof}

\Needspace{16\baselineskip}
\sessionsection{Лекция 6. Спектр и компактная резольвента}
\record{9 октября}{pe5XWL9wtg4}

\begin{reminder}
Множество $\rho(T)$ открыто и состоит из компонент связности $\Pi(T)$
с нулевым дефектным числом, $\sigma(T)=\CC\setminus\rho(T)$ замкнуто.
\end{reminder}

\topic{Точечный, остаточный и непрерывный спектры}{\at{pe5XWL9wtg4}{270}{4:30}}
Для замкнутого $T$ резольвентное множество $\rho(T)$
состоит из тех $\lambda$, для которых $(T-\lambda\id)^{-1}$
ограничен и определён на всём $\Hh$. Поэтому
\[
 \lambda\in\sigma(T)
 \quad\Longleftrightarrow\quad
 T-\lambda\id\colon\Dom(T)\to\Hh\text{ не биективен}.
\]
\begin{definition}
\begin{enumerate}
\item[а)] $\lambda\in\sigma_p(T)$ (точечный спектр), если
  $\ker(T-\lambda\id)\ne\{0\}$. Это собственные значения $T$,
  а $\ker(T-\lambda\id)$~--- собственное подпространство.
\item[б)] $\lambda\in\sigma_r(T)$ (остаточный спектр), если
  $\operatorname{Im}(T-\lambda\id)$ замкнут и отличен от $\Hh$.
\item[в)] $\lambda\in\sigma_c(T)$ (непрерывный спектр), если
  $\operatorname{Im}(T-\lambda\id)$ не замкнут.
\end{enumerate}
\end{definition}
Случаи б) и в) исключают друг друга. Точечный спектр
может пересекаться с каждым из них.

Для ограниченного оператора спектр непуст и компактен.
Для замкнутого неограниченного оператора спектр может быть
пустым, совпадать со всей $\CC$ или быть произвольным
неограниченным замкнутым множеством.

\topic{Свойства резольвенты}{\at{pe5XWL9wtg4}{613}{10:13}}
$R_\lambda(T)=(T-\lambda\id)^{-1}$.
\begin{proposition}[Второе резольвентное тождество Гильберта]
Пусть $T$ и $S$~--- замкнутые операторы в $\Hh$,
$\Dom(S)\subseteq\Dom(T)$ и
$\lambda\in\rho(T)\cap\rho(S)$. Тогда
\[
 R_\lambda(T)-R_\lambda(S)=R_\lambda(T)(S-T)R_\lambda(S).
\]
\end{proposition}
\begin{proof}
Для каждого $x\in\Hh$ вектор $R_\lambda(S)x$
принадлежит $\Dom(S)\subseteq\Dom(T)$, поэтому к нему
можно применить $S-T$. Используем
$(S-\lambda\id)R_\lambda(S)=\id$ на $\Hh$ и
$R_\lambda(T)(T-\lambda\id)=\id$ на $\Dom(T)$.
\[
 \begin{aligned}
 R_\lambda(T)(S-T)R_\lambda(S)x
 &=R_\lambda(T)(S-\lambda\id)R_\lambda(S)x\\
 &\quad-R_\lambda(T)(T-\lambda\id)R_\lambda(S)x\\
 &=R_\lambda(T)x-R_\lambda(S)x.
 \end{aligned}\qedhere
\]
\end{proof}

\begin{proposition}[Первое резольвентное тождество Гильберта]
Для $\lambda,\lambda_0\in\rho(T)$
\[
 R_\lambda(T)-R_{\lambda_0}(T)
 =(\lambda-\lambda_0)R_\lambda(T)R_{\lambda_0}(T).
\]
\end{proposition}
\begin{proof}
Во втором тождестве берём $S=T+(\lambda-\lambda_0)\id$
на $\Dom(T)$. Тогда
\[
 R_\lambda(S)=(S-\lambda\id)^{-1}
 =(T-\lambda_0\id)^{-1}=R_{\lambda_0}(T),\qquad
 S-T=(\lambda-\lambda_0)\id.\qedhere
\]
\end{proof}
Резольвенты коммутируют. При $\lambda\ne\lambda_0$
\[
 R_\lambda(T)R_{\lambda_0}(T)
 =\frac{R_\lambda(T)-R_{\lambda_0}(T)}{\lambda-\lambda_0}
 =\frac{R_{\lambda_0}(T)-R_\lambda(T)}{\lambda_0-\lambda}
 =R_{\lambda_0}(T)R_\lambda(T).
\]

\begin{proposition}[Круг резольвентных точек и аналитичность]
Пусть $T$ замкнут и $\lambda_0\in\rho(T)$. Если
$|\lambda-\lambda_0|<\norm{R_{\lambda_0}(T)}^{-1}$,
то $\lambda\in\rho(T)$ и
\[
 R_\lambda(T)=\sum_{n=0}^\infty
   (\lambda-\lambda_0)^n R_{\lambda_0}(T)^{n+1},
\]
где ряд сходится в операторной норме.
\end{proposition}
\begin{proof}
При $x\in\Dom(T)$ имеем
\[
 \norm{x}=\norm{R_{\lambda_0}(T)(T-\lambda_0\id)x}
 \le\norm{R_{\lambda_0}(T)}\,\norm{(T-\lambda_0\id)x}.
\]
В определении поля регулярности подходит константа
$c_{\lambda_0}=\norm{R_{\lambda_0}(T)}^{-1}$.
По доказательству леммы~\ref{lem:regularity-open} открытый
круг радиуса $c_{\lambda_0}$ с центром $\lambda_0$ лежит
в $\Pi(T)$. Круг связен, поэтому дефект в нём тот же,
что в $\lambda_0$, то есть нулевой. Следовательно,
весь круг лежит в $\rho(T)$ (\at{pe5XWL9wtg4}{1427}{23:47}).

Из первого тождества получаем
\[
 R_\lambda(T)\bigl[\id-(\lambda-\lambda_0)R_{\lambda_0}(T)\bigr]
 =R_{\lambda_0}(T).
\]
Так как $|\lambda-\lambda_0|\norm{R_{\lambda_0}(T)}<1$,
по формуле суммы геометрической прогрессии
\[
 \bigl[\id-(\lambda-\lambda_0)R_{\lambda_0}(T)\bigr]^{-1}
 =\sum_{n=0}^\infty(\lambda-\lambda_0)^n R_{\lambda_0}(T)^n.
\]
Домножаем справа на этот обратный
(\at{pe5XWL9wtg4}{1706}{28:26}) и получаем
\[
 R_\lambda(T)=\sum_{n=0}^\infty(\lambda-\lambda_0)^n
        R_{\lambda_0}(T)^{n+1}.
\]
Ряд сходится в операторной норме.
\end{proof}

\topic{Связь спектров оператора и его резольвенты}{\at{pe5XWL9wtg4}{1888}{31:28}}
\begin{theorem}[Связь спектра с резольвентой]
\label{thm:resolvent-spectrum}
Пусть $T$ замкнут, $\lambda_0\in\rho(T)$ и $\lambda\ne\lambda_0$.
\begin{enumerate}
\item $\lambda\in\rho(T)$ тогда и только тогда, когда
  $(\lambda-\lambda_0)^{-1}\in\rho(R_{\lambda_0}(T))$.
\item $\lambda$~--- собственное значение $T$ тогда и только тогда,
  когда $(\lambda-\lambda_0)^{-1}$~--- собственное значение
  $R_{\lambda_0}(T)$. Их кратности совпадают.
\end{enumerate}
\end{theorem}
\begin{proof}
На $\Dom(T)$ имеем
\[
 T-\lambda\id=
 \bigl[R_{\lambda_0}(T)-(\lambda-\lambda_0)^{-1}\id\bigr]
 (T-\lambda_0\id)(\lambda_0-\lambda).
\]
При раскрытии скобок справа получаем
$(\lambda_0-\lambda)\id+T-\lambda_0\id=T-\lambda\id$.
Оператор $(T-\lambda_0\id)(\lambda_0-\lambda)$ биективно
отображает $\Dom(T)$ на $\Hh$, поскольку
$\lambda_0\in\rho(T)$ и $\lambda_0-\lambda\ne0$.
\begin{enumerate}
\item Поэтому $T-\lambda\id$ биективен ровно тогда, когда
  $R_{\lambda_0}(T)-(\lambda-\lambda_0)^{-1}\id$ биективен.
\item Та же биекция переводит ядро $T-\lambda\id$ на ядро
  $R_{\lambda_0}(T)-(\lambda-\lambda_0)^{-1}\id$.
\end{enumerate}
\end{proof}
Это аналог теоремы об отображении спектра.

\topic{Дискретный спектр и компактная резольвента}{\at{pe5XWL9wtg4}{2571}{42:51}}
\begin{definition}
Замкнутый оператор $T$ имеет чисто дискретный спектр,
если $\sigma(T)$ состоит из собственных значений конечной
кратности без конечных точек накопления. Пишем
$\sigma(T)=\{\lambda_j\}_{j=1}^N$, где $N$ может быть
бесконечным. При $N=\infty$ для каждого компакта
$K\subset\CC$ существует $n(K)$ такое, что
$\lambda_j\notin K$ при $j>n(K)$.
\end{definition}
\begin{lemma}
Пусть $T$ замкнут и при некотором $\lambda_0\in\rho(T)$
оператор $R_{\lambda_0}(T)$ компактен. Тогда спектр
$T$ чисто дискретен.
\end{lemma}
\begin{proof}
Ненулевые точки спектра компактного оператора $R_{\lambda_0}(T)$
суть собственные значения $\lambda_1,\lambda_2,\ldots$ конечной
кратности. При бесконечном числе таких точек $\lambda_j\to0$.
По теореме~\ref{thm:resolvent-spectrum} собственные значения
\[
 \mu_j=\lambda_0+\frac1{\lambda_j}
\]
исчерпывают спектр $T$ и имеют конечную кратность.
Если $\lambda_j\to0$, то $|\mu_j|\to\infty$.
\end{proof}
\begin{example}[Вторая производная на окружности]
\at{pe5XWL9wtg4}{3176}{52:56}.
Пусть $T=-d^2/d\varphi^2$ на $S^1=\RR/(2\pi\ZZ)$.
В $\Hh=L^2(S^1)$ функции имеют вид
$\sum_{n\in\ZZ}a_n\ee^{\ii n\varphi}$, где
$(a_n)_{n\in\ZZ}\in\ell^2(\ZZ)$.
Область определения $T$ состоит из функций, чья вторая
обобщённая производная лежит в $L^2(S^1)$. Поэтому
\[
 \Dom(T)=H^2(S^1)=
 \left\{\sum_{n\in\ZZ}a_n\ee^{\ii n\varphi}:
            (n^2a_n)_{n\in\ZZ}\in\ell^2(\ZZ)\right\}.
\]
Резольвента действует непрерывно из $L^2(S^1)$ в $H^2(S^1)$.
По теореме вложения Соболева на окружности
$H^2(S^1)$ компактно вкладывается в $H^0(S^1)=L^2(S^1)$.
Композиция этих отображений компактна, поэтому $T$
имеет чисто дискретный спектр.

Для вложения Соболева $H^s\hookrightarrow C^k$ в размерности $n$
условие имеет вид $s>n/2+k$ (\at{pe5XWL9wtg4}{3565}{59:25}).
Собственные значения $T$~--- $0,1,4,9,\ldots$,
собственные функции~--- $\ee^{\ii n\varphi}$, $n\in\ZZ$.
\end{example}

\begin{remark}
\at{pe5XWL9wtg4}{3745}{1:02:25}.
Если одна резольвента компактна, то компактны все.
По первому резольвентному тождеству
\[
 R_\lambda(T)=R_{\lambda_0}(T)
 +(\lambda-\lambda_0)R_\lambda(T)R_{\lambda_0}(T).
\]
Компактные операторы образуют идеал в алгебре ограниченных
операторов, поэтому правая часть компактна.
\end{remark}

\topic{Симметричность и полуограниченность}{\at{pe5XWL9wtg4}{4027}{1:07:07}}
\begin{definition}[Симметрический оператор]
Линейный оператор $T$ называется симметрическим, если
$(Tx,y)=(x,Ty)$ для всех $x,y\in\Dom(T)$.
\end{definition}
Замкнутость и плотность области определения не требуются.
\begin{lemma}
Симметричность равносильна вещественности $(Tx,x)$ на
$\Dom(T)$.
\end{lemma}
\begin{proof}
Если $T$ симметричен, то $(Tx,x)=(x,Tx)=\overline{(Tx,x)}$.
Обратно, положим $\Phi(x)=(Tx,x)$ и $\Phi(x,y)=(Tx,y)$.
Формула поляризации в комплексном пространстве имеет вид
\[
 \Phi(x,y)=\tfrac14\bigl[
 \Phi(x+y)-\Phi(x-y)+\ii\Phi(x+\ii y)-\ii\Phi(x-\ii y)\bigr].
\]
При перестановке $x$ и $y$ вещественные слагаемые сохраняются,
а мнимые меняют знак. Значит,
$\Phi(y,x)=\overline{\Phi(x,y)}$, то есть $(Ty,x)=(y,Tx)$
(\at{pe5XWL9wtg4}{4161}{1:09:21}).
\end{proof}

\begin{definition}[Полуограниченность снизу]
Симметрический оператор $T$ полуограничен снизу, если существует
$m\in\RR$ такое, что $(Tx,x)\ge m\norm{x}^2$ при всех $x\in\Dom(T)$.
Точная нижняя грань равна
\[
 m_T=\inf_{x\in\Dom(T),\ x\ne0}\frac{(Tx,x)}{\norm{x}^2}.
\]
\end{definition}
\begin{example}
\at{pe5XWL9wtg4}{4622}{1:17:02}.
Пусть $T=-d^2/d\varphi^2$ на окружности,
$T_0=T|_{C^\infty(S^1)}$. Гладкие функции плотны в $L^2(S^1)$,
и $T=\overline{T_0}$.
Дважды интегрируя по частям, получаем
\[
 (T_0x,y)=\int_0^{2\pi}-x''(\varphi)\overline{y(\varphi)}\dd\varphi
 =\int_0^{2\pi}x(\varphi)\overline{-y''(\varphi)}\dd\varphi
 =(x,T_0y).
\]
Граничные члены сокращаются по периодичности, поэтому
$T_0$ симметричен. Одно интегрирование по частям даёт
\[
 (T_0x,x)=\int_0^{2\pi}|x'(\varphi)|^2\dd\varphi\ge0.
\]
Точная нижняя грань равна нулю.
\end{example}
\begin{example}[Симметрический незамыкаемый оператор]
\at{pe5XWL9wtg4}{4922}{1:22:02}.
Пусть $S$~--- всюду определённый незамыкаемый оператор в $\Hh_1$.
В $\Hh=\Hh_1\oplus\Hh_1$ положим
\[
 \Dom(T)=\Hh_1\oplus\{0\},\qquad
 T\binom{x}{0}=\binom{0}{Sx}.
\]
Квадратичная форма $T$ равна нулю, поэтому $T$ симметричен.
График $T$ есть график $S$ в первом слагаемом аргумента
и втором слагаемом значения. Из замыкаемости $T$ следовала бы
замыкаемость $S$. Область определения $T$ неплотна.
\end{example}
\begin{corollary}
Плотно определённый симметрический оператор $T$ замыкаем,
и $T\subset\overline T=T^{**}\subset T^*$.
\end{corollary}
\begin{proof}
Для $y\in\Dom(T)$ симметричность даёт $(Tx,y)=(x,Ty)$
при всех $x\in\Dom(T)$. Функционал $x\mapsto(Tx,y)$
непрерывен, поэтому $y\in\Dom(T^*)$ и $T^*y=Ty$.
Это означает $T\subset T^*$.
Так как $\Dom(T)$ плотна, $\Dom(T^*)$ тоже плотна.
По теореме~\ref{thm:closability-via-adjoint}
$T^{**}=\overline T$. Замкнутый график $G(T^*)$ содержит
$G(T)$, а значит, и $\overline{G(T)}=G(\overline T)$.
\end{proof}
\sessionsection{Лекция 7. Самосопряжённость и преобразование Кэли}
\record{16 октября}{DdQgQF1YTs8}

\begin{lemma}[Регулярность симметрического оператора]
Для симметрического $A$ в $\Pi(A)$ содержатся все
невещественные точки. Если $A$ ограничен снизу числом
$m$, то $(-\infty,m)\subset\Pi(A)$.
\end{lemma}
\begin{proof}
Пишем $\lambda=a+\ii b$, $a,b\in\RR$. Поскольку
$((A-a\id)x,x)$ вещественно, смешанные члены при
раскрытии скалярного квадрата сокращаются:
\[
 \norm{(A-\lambda\id)x}^2
 =\norm{(A-a\id)x}^2+b^2\norm x^2 .
\]
При $b\ne0$ это оценка снизу с константой $\abs b$.
При $\lambda<m$ имеем
$(m-\lambda)\norm x^2\le((A-\lambda\id)x,x)
\le\norm{(A-\lambda\id)x}\norm x$.
\end{proof}

\topic{Сюръективность сопряжённого сдвига}{\at{DdQgQF1YTs8}{710}{11:50}}
\begin{proposition}
Если $A$ плотно определён и $\lambda\in\Pi(A)$, то
$\Ran(A^*-\bar\lambda\id)=\Hh$.
\end{proposition}
\begin{proof}
Если $\lambda\in\Pi(A)$, то для любого $y\in\Hh$ на
$\Ran(A-\lambda)$ определён функционал
$(A-\lambda)x\mapsto(x,y)$. Оценка регулярности делает его
ограниченным, поскольку
$\abs{(x,y)}\le c_\lambda^{-1}\norm{(A-\lambda)x}\norm y$.
Продолжаем функционал по непрерывности на замыкание
образа и представляем его в виде скалярного произведения
с некоторым $v$ по теореме Рисса. Равенство
$((A-\lambda)x,v)=(x,y)$ на $\Dom(A)$ означает
$v\in\Dom(A^*)$ и $(A^*-\bar\lambda)v=y$.
\end{proof}

\topic{Критерий самосопряжённости}{\at{DdQgQF1YTs8}{1795}{29:55}}
\begin{definition}
Плотно определённый оператор $A$ самосопряжён, если
$A=A^*$ с равенством областей определения.
\end{definition}
Из $A\subset B$ следует $B^*\subset A^*$, что видно
из определения сопряжённого либо из формулы для его графика.
Если $A=A^*$ и $A\subset B\subset B^*$, то
$B^*\subset A$ и все включения обращаются в равенства.
\begin{theorem}[Критерий самосопряжённости]
Для плотно определённого замкнутого симметрического $A$
следующие условия равносильны:
\[
 A=A^*,\qquad
 \Ran(A-\ii)=\Ran(A+\ii)=\Hh,\qquad
 \ker(A^*-\ii)=\ker(A^*+\ii)=0.
\]
\end{theorem}
\begin{proof}
При $A=A^*$ у него нет невещественных собственных чисел:
из $Ax=\pm\ii x$ и вещественности $(Ax,x)$
следует $x=0$. Условие о ядрах эквивалентно тому,
что оба образа плотны, по формуле дефектного пространства.
Образы замкнуты по оценке $\norm{(A\pm\ii)x}^2=\norm{Ax}^2+\norm x^2$.
Если они равны $\Hh$, то для $y\in\Dom(A^*)$ можно решить
$(A-\ii)x=(A^*-\ii)y$. Разность $y-x$ лежит в нулевом ядре
$A^*-\ii$, откуда $y\in\Dom(A)$.
\end{proof}

\topic{Разложение области определения}{\at{DdQgQF1YTs8}{3230}{53:50}}
\begin{theorem}[Разложение области сопряжённого]
Пусть $A$ плотно определён, замкнут и симметричен.
Обозначим $N_\pm=\ker(A^*\mp\ii)$. Тогда
\[
 \Dom(A^*)=\Dom(A)\dotplus N_+\dotplus N_-.
\]
Сумма ортогональна в скалярном произведении графика $A^*$.
\end{theorem}
\begin{proof}
Для $y\in\Dom(A^*)$ раскладываем
$(A^*-\ii)y=(A-\ii)x+v_-$,
где $x\in\Dom(A)$ и $v_-\in N_-$, поскольку
$N_-=\Ran(A-\ii)^\perp$.
Положим $y_-=v_-/(-2\ii)$. Тогда
$(A^*-\ii)y_-=v_-$, а
$y_+=y-x-y_-$ лежит в $\ker(A^*-\ii)=N_+$.
Для $z\in N_\pm$ и $x\in\Dom(A)$ имеем
\[
 (x,z)_{A^*}=(x,z)+(Ax,\pm\ii z)
             =(x,z)\mp\ii(x,\pm\ii z)=0 .
\]
Кроме того, $(u_+,u_-)_{A^*}=(u_+,u_-)+(\ii u_+,-\ii u_-)=0$.
Это даёт ортогональность и прямоту суммы.
\end{proof}

\topic{Преобразование Кэли}{\at{DdQgQF1YTs8}{4349}{1:12:29}}
\begin{definition}
Преобразование Кэли замкнутого симметрического оператора
$A$ есть оператор $U=(A+\ii)(A-\ii)^{-1}$ из
$\Ran(A-\ii)$ на $\Ran(A+\ii)$.
\end{definition}
\begin{theorem}[Прямое и обратное преобразования Кэли]
Преобразование Кэли плотно определённого замкнутого
симметрического оператора изометрично на замкнутой области.
У него $\ker(U-\id)=0$ и $\Ran(U-\id)$ плотен.
Обратно, для изометрии $U$ на замкнутом линеале с
плотным $\Ran(U-\id)$ формулы
\[
 \Dom(A)=\Ran(U-\id),\qquad
 A=\ii(U+\id)(U-\id)^{-1}.
\]
задают плотно определённый замкнутый симметрический
оператор. Самосопряжённости соответствует унитарность
$U$ на всём пространстве.
\end{theorem}
\begin{proof}
При $y=(A-\ii)x$ имеем $Uy=(A+\ii)x$. Равенство норм
$(A-\ii)x$ и $(A+\ii)x$ даёт изометричность. Кроме того,
$Uy-y=2\ii x$ и $Uy+y=2Ax$, откуда получаются формулы
для образа $U-\id$ и действия $A$.

На графиках выполняется замена
$(x,y)\mapsto(u,v)=(y-\ii x,y+\ii x)$ с обратной
\[
 x=(v-u)/(2\ii),\qquad y=(u+v)/2.
\]
В обе стороны это непрерывное линейное отображение
$\Hh\oplus\Hh$. Для обратного перехода надо проверить
однозначность $x\mapsto y$. Если $Uz=z$ на $\Dom(U)$,
то $(Uy,z)=(Uy,Uz)=(y,z)$ для всех $y\in\Dom(U)$.
Значит, $z$ ортогонален $\Ran(U-\id)$ и равен нулю.
Поэтому $v=Uu$ и знание $x=(U-\id)u/(2\ii)$
однозначно задают $y$. График $U$ замкнут, поскольку
область замкнута, а $U$ непрерывен на ней. Следовательно,
график $A$ также замкнут. Из $\norm{y-\ii x}=\norm{y+\ii x}$
следует $\operatorname{Im}(y,x)=0$. По критерию
вещественности квадратичной формы $A$ симметричен.
Унитарность $U$ требует, чтобы оба образа $A\pm\ii$
равнялись $\Hh$, что эквивалентно самосопряжённости.
\end{proof}

\sessionsection{Лекция 8. Самосопряжённые расширения}
\record{23 октября}{Hsak3UepdSQ}

\topic{Продолжение изометрии}{\at{Hsak3UepdSQ}{285}{4:45}}
\begin{theorem}[Фон Неймана о самосопряжённых расширениях]
Пусть $A$ плотно определён, замкнут и симметричен.
Самосопряжённое расширение существует тогда и только
тогда, когда $\dim N_-=\dim N_+$.
Каждому унитарному $U_1:N_-\to N_+$ соответствует
единственное расширение $B$:
\[
 \begin{aligned}
 \Dom(B)&=\Dom(A)\dotplus\{U_1z-z:z\in N_-\},\\
 B(x+U_1z-z)&=Ax+\ii(U_1z+z)\quad (x\in\Dom(A)).
 \end{aligned}
\]
Все самосопряжённые расширения получаются таким способом.
\end{theorem}
\begin{proof}
Пусть $U$~--- преобразование Кэли для $A$. Имеются
ортогональные разложения
\[
 \Hh=\Ran(A-\ii)\oplus N_-,\qquad
 \Hh=\Ran(A+\ii)\oplus N_+ .
\]
Если $B$~--- самосопряжённое расширение $A$, то для
$u\in\Dom(A)$ равенства $(B-\ii)u=(A-\ii)u$ и
$(B+\ii)u=(A+\ii)u$ показывают, что преобразование
Кэли для $B$ продолжает $U$. Оно унитарно на $\Hh$.
Сохранение скалярного произведения переводит
ортогональное дополнение к $\Dom(U)$ в ортогональное
дополнение к $\Ran(U)$, а сюръективность делает это
отображение унитарным между $N_-$ и $N_+$.

Обратно, пусть $U_1:N_-\to N_+$ унитарно.
Оператор $\widetilde U=U\oplus U_1$ определён на всей
$\Hh$, изометричен по ортогональности слагаемых и имеет
образ $\Hh$. Образ $\widetilde U-\id$ содержит
$\Ran(U-\id)=\Dom(A)$, значит, плотен. Поэтому
существует обратное преобразование Кэли. В точках
первоначальной области оно даёт $A$, а замкнутость
и самосопряжённость следуют из унитарности $\widetilde U$.

Чтобы найти область, берём $h\in\Hh$ и раскладываем
$h=h_0+z$, где $h_0\in\Ran(A-\ii)$, $z\in N_-$.
Тогда
$(\widetilde U-\id)h=(U-\id)h_0+U_1z-z$.
Первое слагаемое пробегает $\Dom(A)$, второе~---
$\{U_1z-z:z\in N_-\}$. Сумма прямая по разложению
$\Dom(A^*)$. Поскольку $B\subset A^*$, на $U_1z-z$
оператор $B$ действует как $A^*$, давая $\ii U_1z+\ii z$.
Взаимная обратимость преобразований Кэли даёт
биекцию между расширениями $B$ и продолжениями
$\widetilde U=U\oplus U_1$.
\end{proof}

\begin{figure}[ht]
\centering
\begin{tikzpicture}[>=Stealth]
 \node (left) at (0,1) {$\Ran(A-\ii)$};
 \node (right) at (4.8,1) {$\Ran(A+\ii)$};
 \node (lker) at (0,-.35) {$\ker(A^*+\ii)$};
 \node (rker) at (4.8,-.35) {$\ker(A^*-\ii)$};
 \node at (0,.35) {$\oplus$};
 \node at (4.8,.35) {$\oplus$};
 \draw[->] (left)--node[above] {$U$} (right);
 \draw[->] (lker)--node[above] {$U_1$} (rker);
\end{tikzpicture}
\caption{Продолжение преобразования Кэли на дефектных
пространствах. По доске в \at{Hsak3UepdSQ}{2360}{39:20}.}
\end{figure}

\topic{Замкнутые симметрические расширения}{\at{Hsak3UepdSQ}{573}{9:33}}
\begin{proposition}
Замкнутые симметрические расширения $A$ соответствуют
изометриям $V:E_-\to E_+$, где $E_\pm\subset N_\pm$
замкнуты, а $V$ отображает $E_-$ на $E_+$. Преобразование
Кэли расширения есть $U\oplus V$.
\end{proposition}
\begin{proof}
Пусть $\widetilde U$~--- изометрическое расширение
$U$ с замкнутой областью. Область $\widetilde U$
распадается на $\Dom(U)$ и ортогональный остаток
$E_-\subset N_-$. Этот остаток переходит изометрически
на $E_+\subset N_+$, поскольку $\widetilde U$
сохраняет скалярное произведение и совпадает с $U$
на $\Dom(U)$. Обратно, для $U\oplus V$ образ
$(U\oplus V)-\id$ содержит $\Dom(A)$, поэтому
обратное преобразование даёт плотно определённое
симметрическое расширение. Замкнутость следует из
замкнутости $\Dom(U)\oplus E_-$. Его дефектные
пространства равны $N_-\ominus E_-$ и $N_+\ominus E_+$,
так как это ортогональные дополнения к области
и образу новой изометрии соответственно.
\end{proof}
В обсуждении области при \at{Hsak3UepdSQ}{4605}{1:16:45}
подчёркивается роль прямой суммы $N_-\dotplus N_+$:
если оба дефектных пространства $n$-мерны, в
$2n$-мерной сумме добавляется $n$-мерный график
$\{U_1z-z\}$. Вся сумма дефектных пространств добавляется
к $\Dom(A)$ в области $A^*$.

\sessionsection{Лекция 9. Нормальные операторы и замкнутые формы}
\record{6 ноября}{gTOamZ5pJdY}

\topic{Оператор $T^*T$}{\at{gTOamZ5pJdY}{26}{0:26}}
\begin{theorem}
Пусть $T$ замкнут и плотно определён. Тогда оператор
$T^*T$ неотрицателен и самосопряжён, а
$\Dom(T^*T)$ плотна в $\Dom(T)$ относительно
графовой нормы $\norm{\cdot}_T$.
\end{theorem}
\begin{proof}
По формуле для графика сопряжённого и замкнутости $T$
имеем $\Hh\oplus\Hh=G(T)\oplus WG(T^*)$.
Для $h\in\Hh$ запишем ортогональное разложение
\[
 (h,0)=(x,Tx)+(-T^*y,y),
 \qquad x\in\Dom(T),\quad y\in\Dom(T^*).
\]
Из второй координаты $y=-Tx$. Поэтому
$x\in\Dom(T^*T)$ и $x+T^*Tx=h$ по первой координате.
Имеем $(T^*Tx,x)=\norm{Tx}^2$. Отсюда
\[
 \norm{x}^2+\norm{Tx}^2=(h,x)\le\norm h\norm x .
\]
Поэтому решение единственно, а
$S=(\id+T^*T)^{-1}$ ограничен и всюду определён.
Кроме того, $(Sh,h)=(x,x+T^*Tx)=\norm x^2+\norm{Tx}^2$
вещественно, откуда $S$ симметричен и самосопряжён.
Он инъективен, а его образ плотен, так как
$(\Ran S)^\perp=\ker S^*=0$.

Обратный $S^{-1}$ самосопряжён на $\Ran S$.
Действительно, для $y\in\Dom((S^{-1})^*)$ найдётся
$z$ такое, что $(S^{-1}x,y)=(x,z)$ для всех
$x\in\Ran S$. Подстановка $x=Sw$ даёт
$(w,y)=(Sw,z)=(w,Sz)$ для всех $w\in\Hh$.
Следовательно, $y=Sz$ и $S^{-1}y=z$. Это доказывает
равенство сопряжённого с обратным. Вычитание $\id$
сохраняет самосопряжённость с той же областью.

Для плотности рассмотрим $v\in\Dom(T)$, ортогональный
$\Dom(T^*T)$ в скалярном произведении графика. Для
$x\in\Dom(T^*T)$
\[
 0=(v,x)_T=(v,x)+(Tv,Tx)=(v,(\id+T^*T)x).
\]
Образ $\id+T^*T$ есть $\Hh$, значит, $v=0$.
В полном пространстве $\Dom(T)$ с графовой нормой
это доказывает плотность $\Dom(T^*T)$.
\end{proof}

\topic{Нормальность}{\at{gTOamZ5pJdY}{2145}{35:45}}
\begin{definition}[Нормальный оператор]
Замкнутый плотно определённый $T$ называется нормальным,
если $T^*T=TT^*$ с равенством областей определения.
\end{definition}
Самосопряжённый оператор является нормальным.
В частности, для самосопряжённого $T$ в силу
предыдущей теоремы $T^2$ самосопряжён и плотно
определён. Повторяя рассуждение для степеней $2^k$,
получаем плотную область для любого многочлена от $T$:
$\Dom(T^j)$ содержит $\Dom(T^{2^k})$ при $2^k\ge j$.

\begin{theorem}
Для нормального $T$ справедливы
$\Dom(T)=\Dom(T^*)$ и $\norm{Tx}=\norm{T^*x}$ для $x$
из общей области. Кроме того,
\[
 \norm{(T-\lambda)x}=\norm{(T^*-\bar\lambda)x},\qquad
 \ker(T-\lambda)=\ker(T^*-\bar\lambda).
\]
Собственные подпространства при разных собственных значениях
ортогональны и $\Pi(T)=\rho(T)$.
\end{theorem}
\begin{proof}
Пусть $L=\Dom(T^*T)=\Dom(TT^*)$. На $L$
равенство произведений даёт
$\norm{Tx}^2=(x,T^*Tx)=(x,TT^*x)=\norm{T^*x}^2$.
По теореме о $T^*T$ линеал $L$ плотен и в $\Dom(T)$,
и в $\Dom(T^*)$ в соответствующих графовых нормах.
Для второго применения используем $T^{**}=T$.
Для $x\in\Dom(T)$ возьмём $x_n\in L$ с
$x_n\to x$ по норме $\norm{\cdot}_T$.
На $L$ эта норма равна $\norm{\cdot}_{T^*}$, поэтому
$x_n$ сходится в полной области для $T^*$.
Предел в $\Hh$ один и тот же. Значит, $x\in\Dom(T^*)$
и равенство норм переносится на $x$.
Обратное включение получается перестановкой $T$ и $T^*$.

Теперь раскрытие квадратов норм показывает, что
смешанные слагаемые совпадают благодаря
$(Tx,x)=(x,T^*x)$. Отсюда равенство для сдвигов и ядер.
Если $Tx=\lambda x$ и $Ty=\mu y$, то $T^*y=\bar\mu y$,
поэтому
$\lambda(x,y)=(Tx,y)=(x,T^*y)=\mu(x,y)$.
При $\lambda\ne\mu$ это даёт ортогональность.
Если $\lambda\in\Pi(T)$, то образ $T-\lambda$
замкнут. Его ортогональное дополнение равно
$\ker(T^*-\bar\lambda)=\ker(T-\lambda)=0$.
Следовательно, образ есть $\Hh$ и $\lambda\in\rho(T)$.
\end{proof}

\topic{Форма, согласованная с оператором}{\at{gTOamZ5pJdY}{4369}{1:12:49}}
\begin{definition}[Положительно определённая замкнутая форма]
Полуторалинейная форма $a[x,y]$ задана на плотном
линеале $\Dom[a]$ и линейна по $x$, антилинейна по $y$.
Эрмитовость означает $a[x,y]=\overline{a[y,x]}$.
Точная нижняя грань определяется как
$m_a=\inf_{0\ne x\in\Dom[a]} a[x,x]/\norm x^2$.
Если $m_a>0$, то форма задаёт на $\Dom[a]$
скалярное произведение с нормой $\abs x_a=a[x,x]^{1/2}$.
Форма замкнута, когда эта норма делает $\Dom[a]$ полным.
\end{definition}
Из $\abs x_a\ge\sqrt{m_a}\norm x$ следует непрерывность
вложения пространства формы в $\Hh$. Поэтому предел
по норме формы, если он существует, совпадает с
пределом в $\Hh$.

\begin{definition}[Согласованность]
Самосопряжённый $A$ согласован с формой $a$, если
\[
 \Dom(A)\subset\Dom[a],\qquad
 a[x,y]=(Ax,y)\quad (x\in\Dom(A),\ y\in\Dom[a]).
\]
\end{definition}
\begin{lemma}[Плотность операторной области]
Если $a$ замкнута, $m_a>0$ и $A$ согласован с $a$,
то $\Dom(A)$ плотна в $\Dom[a]$ по норме $\abs{\cdot}_a$.
Точные нижние грани совпадают: $m_A=m_a$.
\end{lemma}
\begin{proof}
Беря $y=x$ в условии согласования, получаем
$(Ax,x)=a[x,x]\ge m_a\norm x^2$ на $\Dom(A)$,
поэтому $m_A\ge m_a>0$. Точка $0$ принадлежит
$\Pi(A)$ по оценке снизу, а для самосопряжённого
оператора $\Pi(A)=\rho(A)$, так что $\Ran A=\Hh$.
Если $v\in\Dom[a]$ ортогонален $\Dom(A)$ по форме,
то $0=a[x,v]=(Ax,v)$ для всех $x\in\Dom(A)$.
Значит, $v=0$ и плотность доказана.
Для $y\in\Dom[a]\setminus\{0\}$ берём
$x_n\in\Dom(A)$ с $\abs{x_n-y}_a\to0$.
Тогда $a[x_n,x_n]\to a[y,y]$ и $\norm{x_n}\to\norm y$,
поэтому $m_A\le a[y,y]/\norm y^2$. Берём инфимум по $y$.
\end{proof}

\sessionsection{Лекция 10. Положительные операторы и формы}
\record{13 ноября}{1OVSIVm2iks}

\topic{От оператора к форме}{\at{1OVSIVm2iks}{380}{6:20}}
\begin{theorem}[Форма положительного оператора]
Пусть $A$ самосопряжён и $m_A>0$. Тогда существует
единственная замкнутая форма $a$ со строго
положительной нижней гранью, согласованная с $A$.
Она задаётся формулой
\[
 \Dom[a]=\Dom(A^{1/2}),\qquad
 a[x,y]=(A^{1/2}x,A^{1/2}y).
\]
\end{theorem}
\begin{proof}
Так как $0\in\rho(A)$, оператор $A^{-1}$ ограничен
и всюду определён. Он самосопряжён и неотрицателен,
поскольку при $h=Ax$ имеем
$(A^{-1}h,h)=(x,Ax)\ge0$. По непрерывному исчислению
определим $B=(A^{-1})^{1/2}$; $B$ ограничен,
самосопряжён и $B^2=A^{-1}$. Из инъективности $B^2$
следует инъективность $B$, а из самосопряжённости~---
плотность его образа. Обратный $B^{-1}$ на $\Ran B$
замкнут по перестановке координат графика и
самосопряжён по правилу сопряжения обратного.
Область $(B^{-1})^2$ состоит из $x=By$ таких, что
$y\in\Ran B$, то есть равна $\Ran B^2=\Dom(A)$.
На этой области $(B^{-1})^2=A$, поэтому берём
$A^{1/2}=B^{-1}$.

На $\Dom(A^{1/2})$ неравенство
$\norm x\le\norm B\norm{A^{1/2}x}$ показывает
эквивалентность нормы $\norm{A^{1/2}x}$ графовой норме
$A^{1/2}$. Отсюда полнота нормы формы и строго
положительная нижняя грань. Согласованность следует
из $(A^{1/2}x,A^{1/2}y)=(Ax,y)$ при
$x\in\Dom(A)$ и $y\in\Dom(A^{1/2})$.

Для единственности пусть $a_1$~--- другая такая форма.
На $\Dom(A)$ совпадают $a_1[x,y]$ и $a[x,y]$,
следовательно, и их нормы. Лемма о плотности
показывает, что обе области являются пополнениями
$\Dom(A)$ по этой норме. Эти пополнения вложены в
$\Hh$ и их пределы от одной последовательности
совпадают как векторы $\Hh$. Поэтому совпадают
и области, и значения форм.
\end{proof}

\topic{От формы к оператору}{\at{1OVSIVm2iks}{2360}{39:20}}
\begin{theorem}[Оператор положительной замкнутой формы]
Пусть $a$~--- плотно определённая замкнутая эрмитова
форма с $m_a>0$. Тогда существует единственный
самосопряжённый оператор $A$, согласованный с $a$.
Явное описание области имеет вид
\[
 \Dom(A)=\{x\in\Dom[a]:\exists h\in\Hh\ \forall y\in\Dom[a]\
                   a[x,y]=(h,y)\},\qquad Ax=h.
\]
\end{theorem}
\begin{proof}
Строим обратный к искомому оператору. Для каждого
$h\in\Hh$ функционал $l_h(y)=(h,y)$ антилинеен по $y$
и ограничен в норме формы, поскольку
$\abs{(h,y)}\le m_a^{-1/2}\norm h\,\abs y_a$.
По теореме Рисса, применённой с произведением $a$,
существует единственное $x\in\Dom[a]$ с
\[
 a[x,y]=(h,y)\qquad (y\in\Dom[a]).
\]
Пишем $x=Bh$. Линейность следует из единственности
представляющего вектора. Беря $y=Bh$, получаем
\[
 m_a\norm{Bh}^2\le a[Bh,Bh]=(h,Bh)\le \norm h\norm{Bh},
 \qquad \norm{B}\le m_a^{-1}.
\]
При $Bh=0$ равенство $a[Bh,y]=(h,y)$ означает, что
$h\perp\Dom[a]$, откуда $h=0$.
Эрмитовость даёт
$(h,Bk)=a[Bh,Bk]=\overline{a[Bk,Bh]}=(Bh,k)$,
поэтому $B$ самосопряжён. Его образ плотен, так как
ортогональное дополнение к образу равно $\ker B^*=0$.
Следовательно, обратный $A=B^{-1}$ самосопряжён
и согласован с $a$ по построению. Для $x$ из выписанной
в теореме области равенство $a[x,y]=(h,y)$ в силу
единственности в теореме Рисса означает $x=Bh$, что
проверяет явное описание. Если $A'$ также согласован
с $a$, то для $x\in\Dom(A')$ получаем $x=B A'x$.
Значит, $A'\subset A$, а оба оператора самосопряжённы,
так что $A'=A$.
\end{proof}

\sessionsection{Лекция 11. Расширение Фридрихса и функции от унитарного оператора}
\record{20 ноября}{E-jM7sTmSKc}

\topic{Полуограниченные формы}{\at{E-jM7sTmSKc}{50}{0:50}}
\begin{definition}[Замкнутая полуограниченная форма]
Для плотно определённой эрмитовой формы с конечной
нижней гранью $m_a$ рассматривают форму на той же области
$a_\alpha[x,y]=a[x,y]+(\alpha-m_a)(x,y)$ при $\alpha>0$.
Её точная нижняя грань равна $\alpha$. Форма $a$
называется замкнутой, если $\Dom[a]$ полна по
$a_\alpha[x,x]^{1/2}$.
\end{definition}
\begin{lemma}
Это определение не зависит от выбора $\alpha>0$.
Замкнутая полуограниченная форма $a$ имеет единственный
согласованный самосопряжённый оператор.
\end{lemma}
\begin{proof}
При $\alpha\ge\beta>0$ имеем
$a_\beta[x,x]\le a_\alpha[x,x]\le
(\alpha/\beta)a_\beta[x,x]$, так как
$a_\beta[x,x]\ge\beta\norm x^2$. Значит, нормы
эквивалентны. После построения $A_\alpha$ для положительной
замкнутой формы $a_\alpha$ полагаем
$A=A_\alpha-(\alpha-m_a)\id$ на той же области.
Тогда $(Ax,y)=a_\alpha[x,y]-(\alpha-m_a)(x,y)=a[x,y]$.
Любой другой согласованный оператор после такого же
сдвига был бы согласован с $a_\alpha$ и совпал бы с
$A_\alpha$ по предыдущей теореме.
\end{proof}

\topic{Конструкция Фридрихса}{\at{E-jM7sTmSKc}{1384}{23:04}}
\begin{theorem}[Расширение Фридрихса]
У плотно определённого симметрического полуограниченного
снизу оператора $A_0$ существует самосопряжённое
расширение $A_F$ с той же точной нижней гранью.
Если $a$~--- замыкание формы $a_0[x,y]=(A_0x,y)$
на $\Dom(A_0)$, то $A_F$~--- единственное
самосопряжённое расширение $A_0$ с областью,
содержащейся в $\Dom[a]$.
\end{theorem}
\begin{proof}
Сдвигом $A_0$ можно сделать нижнюю грань строго
положительной; в конце сдвиг убирается. Обозначим
через $V$ абстрактное гильбертово пополнение
линеала $\Dom(A_0)$ по $a_0[x,x]^{1/2}$.
Оценка $\norm x\le m_{A_0}^{-1/2}a_0[x,x]^{1/2}$
означает, что включение $\Dom(A_0)\to\Hh$ продолжается
до непрерывного $j:V\to\Hh$. Для отождествления
$V$ с подмножеством $\Hh$ требуется инъективность $j$.
Пусть $x_n$ сходится в $V$ и стремится к нулю в $\Hh$.
Для каждого $y\in\Dom(A_0)$ по симметричности
$a_0[x_n,y]=(x_n,A_0y)\to0$. Значит, предел $x_n$ в $V$
ортогонален плотному в $V$ подпространству $\Dom(A_0)$
и равен нулю. Следовательно, $j$ инъективно.

На $j(V)=\Dom[a]$ задаём
форму $a$ скалярным произведением пространства $V$.
Она замкнута, потому что $V$ полно. Переход к пределу
сохраняет $a[x,x]\ge m_{A_0}\norm x^2$, а на
$\Dom(A_0)$ форма совпадает с $a_0$. Отсюда
$m_a=m_{A_0}$.

Пусть $A_F$~--- согласованный с $a$ оператор.
Если $x\in\Dom(A_0)$, то равенство $a[x,y]=(A_0x,y)$
верно при $y\in\Dom(A_0)$ и продолжается на
$y\in\Dom[a]$. Слева это непрерывность по норме формы,
справа~--- непрерывность вложения в $\Hh$.
Критерий области согласованного оператора даёт
$x\in\Dom(A_F)$ и $A_F x=A_0x$.
Точная нижняя грань $A_F$ равна $m_a$.

Для единственности пусть $B$~--- самосопряжённое
расширение с $\Dom(B)\subset\Dom[a]$.
Если $x\in\Dom(B)$ и $y\in\Dom(A_0)$, берём
$x_n\in\Dom(A_0)$ с $x_n\to x$ по норме формы.
Тогда
\[
 a[x,y]=\lim_n a_0[x_n,y]=\lim_n(x_n,A_0y)
       =(x,By)=(Bx,y).
\]
Снова по плотности равенство выполняется для всех
$y\in\Dom[a]$, откуда $B\subset A_F$. Самосопряжённость
обоих даёт $B=A_F$.
\end{proof}

\topic{Непрерывное исчисление для унитарного оператора}{\at{E-jM7sTmSKc}{5102}{1:25:02}}
\begin{theorem}[Исчисление непрерывных функций на окружности]
Пусть $U$ унитарен. Существует непрерывное унитальное
исчисление $f\mapsto f(U)$ для $f\in C(S^1)$,
которое переводит $z\mapsto z$ в $U$ и удовлетворяет
\[
 \begin{gathered}
 f(U)^*=\bar f(U),\qquad f(U)g(U)=(fg)(U),\\
 f(U)\ge0\text{ при }f\ge0,\qquad
 \norm{f(U)}\le\norm f_\infty .
 \end{gathered}
\]
\end{theorem}
\begin{proof}
Спектр $U$ лежит на $S^1$. Для гладкой функции пишем
$f(\ee^{\ii\theta})=\sum_{k\in\ZZ}\widehat f(k)\ee^{\ii k\theta}$
и определяем $f(U)=\sum_k\widehat f(k)U^k$.
Коэффициенты убывают быстрее каждой степени $|k|^{-1}$,
а $\norm{U^k}=1$ для всех $k$, так что ряд сходится
абсолютно по операторной норме. Линейность и значения
на $1,z$ очевидны из определения.
Так как $(U^k)^*=U^{-k}$ и
$\widehat{\bar f}(-k)=\overline{\widehat f(k)}$, сопряжение
ряда даёт $\bar f(U)$. Для умножения перегруппировываем
абсолютно сходящийся двойной ряд:
\[
 f(U)g(U)
 =\sum_{m\in\ZZ} U^m\sum_{k\in\ZZ}\widehat f(k)\widehat g(m-k)
 =\sum_{m\in\ZZ} U^m\widehat{fg}(m).
\]
Если гладкая $f\ge0$, то гладкая вещественная
$h=\sqrt{f+\varepsilon}$ удовлетворяет
$f(U)=h(U)^*h(U)-\varepsilon\id$. Поэтому
$(f(U)v,v)=\norm{h(U)v}^2-\varepsilon\norm v^2\ge
-\varepsilon\norm v^2$. Левая часть не зависит от
$\varepsilon$, значит, $f(U)\ge0$. Применив положительность
к $\norm f_\infty^2-\abs f^2$, получаем
\[
 \norm{f(U)v}^2=\bigl((|f|^2)(U)v,v\bigr)
 \le\norm f_\infty^2\norm v^2.
\]

Если $f$ лишь непрерывна, берём гладкие $f_n\to f$
равномерно. Тогда
$\norm{f_n(U)-f_m(U)}\le\norm{f_n-f_m}_\infty$,
так что $f_n(U)$ сходится по норме, и предел
не зависит от аппроксимации. Сопряжение,
умножение и оценка нормы переносятся по
непрерывности. Для неотрицательной непрерывной
$f$ можно сначала приблизить $\sqrt f$ гладкими
функциями, а затем использовать
$f(U)=(\sqrt f(U))^*\sqrt f(U)$.
\end{proof}

\Needspace{10\baselineskip}
\part{Семинары}
\setcounter{section}{0}
\renewcommand{\thetheorem}{С\arabic{section}.\arabic{theorem}}

\sessionsection{Семинар 1. Оператор интегрирования}
\record{4 сентября}{xOAr776deQI}

\topic{Степени и резольвента}{\at{xOAr776deQI}{3}{0:03}}
\begin{exercise}[Оператор Вольтерры]
На $C[0,1]$ рассмотрим $(Ax)(t)=\int_0^t x(s)\dd s$.
Найти $A^n$, резольвенту и $f(A)$ для $f$,
голоморфной в окрестности спектра.
\end{exercise}
\begin{proof}[Решение]
При больших $|\lambda|$ выносим $\lambda$ и применяем
ряд Неймана \at{xOAr776deQI}{142}{2:22}:
\[
 R_\lambda(A)=(\lambda\id-A)^{-1}
 =\frac{1}{\lambda}\left(\id-\frac{A}{\lambda}\right)^{-1}
 =\sum_{n=0}^\infty\frac{A^n}{\lambda^{n+1}} .
\]
Для вычисления ряда понадобятся степени
\[
 (A^{n+1}x)(t)=\frac{1}{n!}\int_0^t(t-s)^n x(s)\dd s,
 \qquad n\ge0 .
\]
Для $A^2$ меняем порядок интегрирования по треугольнику
$0\le\tau\le t$, $0\le s\le\tau$. При внешнем
интегрировании по $s$ получаем $0\le s\le t$,
$s\le\tau\le t$:
\[
 (A^2x)(t)=\int_0^t\int_0^\tau x(s)\dd s\,\dd\tau
 =\int_0^t x(s)\int_s^t 1\dd\tau\,\dd s
 =\int_0^t(t-s)x(s)\dd s .
\]
\begin{center}
\begin{minipage}{\linewidth}
\centering
\begin{tikzpicture}[>=Stealth,scale=.84]
 \fill[pattern=vertical lines] (0,0)--(1.75,0)--(1.75,1.75)--cycle;
 \draw[->] (-.15,0)--(2.72,0) node[right] {$\tau$};
 \draw[->] (0,-.15)--(0,2.48) node[above] {$s$};
 \draw (-.15,-.15)--(2.05,2.05);
 \draw (1.75,0)--(1.75,1.75)--(0,1.75);
 \node[below] at (1.75,0) {$t$};
 \node[left] at (0,1.75) {$t$};
\end{tikzpicture}
\captionof{figure}{Область интегрирования при вычислении $A^2$.
По доске в \at{xOAr776deQI}{570}{9:30}.}
\end{minipage}
\end{center}
Для остальных степеней переставляем интегралы так же.
При переходе от $A^{n+1}$ к $A^{n+2}$ внутренний
интеграл по $\tau$ равен
$\int_s^t(\tau-s)^n\dd\tau/n!=(t-s)^{n+1}/(n+1)!$.

В ряде для резольвенты выделяем $n=0$ и сдвигаем индекс
\at{xOAr776deQI}{636}{10:36}:
\[
 R_\lambda(A)=\frac{\id}{\lambda}
      +\sum_{n=0}^\infty\frac{A^{n+1}}{\lambda^{n+2}}.
\]
Подставляем формулу для $A^{n+1}$ и суммируем
экспоненциальный ряд:
\[
 \begin{aligned}
 (R_\lambda(A)x)(t)
 &=\frac{x(t)}{\lambda}
     +\frac{1}{\lambda^2}\sum_{n=0}^\infty\frac{1}{n!}
         \int_0^t\left(\frac{t-s}{\lambda}\right)^n x(s)\dd s\\
 &=\frac{x(t)}{\lambda}
     +\frac{1}{\lambda^2}\int_0^t x(s)\ee^{(t-s)/\lambda}\dd s .
 \end{aligned}
\]
По оценке $\norm{A^{n+1}}\le1/(n+1)!$ ряд резольвенты
сходится при всех $\lambda\ne0$. Спектр ограниченного
оператора непуст, значит, $\sigma(A)=\{0\}$
\at{xOAr776deQI}{901}{15:01}.

Если $f$ аналитична в окрестности нуля, берём малую
положительно ориентированную окружность $|\lambda|=r$.
Подставляем в формулу для $f(A)$ найденную резольвенту
и сразу применяем оператор к $x$
\at{xOAr776deQI}{950}{15:50}:
\[
 \begin{aligned}
 (f(A)x)(t)
 &=\frac{1}{2\pi\ii}\int_{|\lambda|=r} f(\lambda)
   \left(\frac{x(t)}{\lambda}+
    \frac{1}{\lambda^2}\int_0^t x(s)\ee^{(t-s)/\lambda}\dd s\right)
     \dd\lambda\\
 &=f(0)x(t)+
    \int_0^t
    \underbrace{\left(\frac{1}{2\pi\ii}\int_{|\lambda|=r}
        \frac{f(\lambda)}{\lambda^2}\ee^{(t-s)/\lambda}\dd\lambda
        \right)}_{F(t-s)}x(s)\dd s .
 \end{aligned}
\]
В первом слагаемом $x(t)$ выносим за интеграл, а интеграл
от $f(\lambda)/\lambda$ даёт $f(0)$. Во втором меняем
порядок интегрирования и обозначаем внутренний
интеграл по $\lambda$ через $F(t-s)$.

Пусть $f(\lambda)=\sum_{n=0}^\infty a_n\lambda^n$.
Перемножаем этот ряд Тейлора с рядом экспоненты
\at{xOAr776deQI}{1387}{23:07}:
\[
 \frac{f(\lambda)}{\lambda^2}\ee^{(t-s)/\lambda}
 =\sum_{n,m\ge0}a_n\,\frac{(t-s)^m}{m!}\,\lambda^{n-m-2}.
\]
Это ряд Лорана подынтегральной функции. В интеграле
остаётся коэффициент при $\lambda^{-1}$, то есть
$n-m-2=-1$ или $n=m+1$. Поэтому
\[
 F(\xi)=\sum_{m=0}^\infty a_{m+1}\frac{\xi^m}{m!},
 \qquad
 (f(A)x)(t)=f(0)x(t)+\int_0^t F(t-s)x(s)\dd s .
\]
\end{proof}

\topic{Решение уравнения с нулевыми начальными данными}{\at{xOAr776deQI}{1764}{29:24}}
\begin{exercise}
Решить задачу Коши
$y^{(n)}+a_1y^{(n-1)}+\cdots+a_ny=x$,
$y(0)=y'(0)=\cdots=y^{(n-1)}(0)=0$,
с постоянными коэффициентами, используя оператор интегрирования.
\end{exercise}
\begin{proof}[Решение]
При указанных начальных значениях
$A^n y^{(n)}=y$ и $A^ny^{(n-j)}=A^jy$.
Поэтому после интегрирования получаем
$f(A)y=A^nx$, где $f(z)=1+a_1z+\cdots+a_nz^n$.
Поскольку $f(0)=1$, можно обратить $f(A)$ в исчислении и
получить $y=h(A)x$ с $h(z)=z^n/f(z)$. Для ядра решения
замена $z=1/\lambda$ даёт
$\lambda^{-2}\dd\lambda=-\dd z$ и меняет ориентацию.
Знак от дифференциала компенсируется заменой ориентации.
Кроме того, $h(0)=0$, так что остаётся только интегральный
член с ядром
\[
 F(r)=\frac{1}{2\pi\ii}\int_{|z|=R}
          \frac{\ee^{rz}}{z^n+a_1z^{n-1}+\cdots+a_n}\dd z,\qquad
 y(t)=\int_0^t F(t-s)x(s)\dd s .
\]
\end{proof}

\topic{Совпадение трёх записей исчисления}{\at{xOAr776deQI}{3004}{50:04}}
На семинаре сравниваются степенной ряд, интеграл Коши
и формула Фурье для целой функции с быстрым убыванием на
$\RR$, например $f(z)=\exp(-z^2)$, и ограниченного
самосопряжённого $A$. На контуре $|\lambda|=R>\norm A$
резольвента раскладывается в ряд Неймана, после чего
\[
 \frac{1}{2\pi\ii}\int_{|\lambda|=R} f(\lambda)
            \sum_{n\ge0}\lambda^{-n-1}A^n\dd\lambda
     =\sum_{n\ge0} A^n\frac{f^{(n)}(0)}{n!}
\]
по формуле Коши для производной.
В записи Фурье, с соглашением
$\widehat f(t)=\int_{\RR}\ee^{-\ii tx}f(x)\dd x$,
пишем $\ee^{\ii tA}=\sum_{n\ge0}(\ii t)^nA^n/n!$.
При допустимом почленном интегрировании коэффициент
при $A^n$ равен
\[
 \frac{1}{2\pi n!}\int_\RR (\ii t)^n\widehat f(t)\dd t
 =\frac{1}{2\pi n!}\int_\RR\widehat{f^{(n)}}(t)\dd t
 =\frac{f^{(n)}(0)}{n!}.
\]
Первый переход получается интегрированием по частям
в определении преобразования Фурье, второй~---
формулой обращения в точке $0$. Для гауссовой функции
почленный переход оправдан интегрируемостью
$|\widehat f(t)|\exp(\norm A|t|)$.
\topic{Домашнее задание}{\at{xOAr776deQI}{4079}{1:07:59}}
А.\,А.~Кириллов, А.\,Д.~Гвишиани,
\emph{Теоремы и задачи функционального анализа}, с.~274.

\begin{description}
\item[644*.] Пусть $A$~--- самосопряжённый оператор,
$a\id\le A\le b\id$, а многочлен $p$ неотрицателен на
$[a,b]$. Доказать, что $p(A)\ge0$.

\item[645.] При $a\id\le A\le b\id$ доказать
непрерывность отображения $p\mapsto p(A)$, если
многочлены снабжены нормой $C[a,b]$, а операторы~---
операторной нормой.

\item[646.] Пусть $A$~--- ограниченный самосопряжённый
оператор. Доказать, что $U(t)=\ee^{\ii tA}$ унитарен
при каждом $t\in\RR$ и что
\[
 U(t)U(s)=U(t+s),\qquad U(t)^*=U(-t).
\]

\item[647.] В условиях задачи~646 доказать
дифференцируемость операторной функции $U(t)$ и равенства
\[
 U'(t)=\ii AU(t)=\ii U(t)A.
\]

\item[648*.] Доказать, что непрерывная в операторной
норме функция $U(t)$, удовлетворяющая равенствам
$U(t)U(s)=U(t+s)$ и $U(t)^*=U(-t)$, имеет вид
$U(t)=\ee^{\ii tA}$ с ограниченным самосопряжённым $A$.
На семинаре~2 в условие также включено $U(0)=\id$.
\end{description}

\sessionsection{Семинар 2. Положительность и унитарные группы}
\record{11 сентября}{uLlyRG0VRnY}

\topic{Положительность многочлена от оператора}{\at{uLlyRG0VRnY}{0}{0:00}}
\begin{exercise}[№~644* из сборника Кириллова и Гвишиани]
Пусть $A$ ограничен, самосопряжён и $a\id\le A\le b\id$.
Пусть $p$~--- вещественный многочлен, неотрицательный на
$[a,b]$. Доказать, что $p(A)\ge0$.
\end{exercise}
\begin{proof}[Решение]
Так как $A=A^*$, его спектр лежит в $[-\norm A,\norm A]$.
По теореме о границах спектра самосопряжённого оператора
(\at{uLlyRG0VRnY}{326}{5:26}) имеем $\sigma(A)\subset[m,M]$, где
\[
 m=\inf_{\norm v=1}(Av,v),\qquad M=\sup_{\norm v=1}(Av,v).
\]
Из $a\id\le A\le b\id$ следует $a\norm v^2\le(Av,v)\le b\norm v^2$.
При $\norm v=1$ получаем $a\le m\le M\le b$.

\begin{center}
\begin{minipage}{\linewidth}
\centering
\begin{tikzpicture}[>=Stealth]
 \draw[->] (-2,0)--(2,0);
 \draw[->] (0,-.55)--(0,.75);
 \draw[line width=1pt] (-1.22,0)--(1.15,0);
 \draw (-1.22,-.09)--(-1.22,.09) node[below=5pt] {$a$};
 \draw (1.15,-.09)--(1.15,.09) node[below=5pt] {$b$};
\end{tikzpicture}
\captionof{figure}{Отрезок $[a,b]$ на вещественной оси
(\at{uLlyRG0VRnY}{426}{7:06}).}
\end{minipage}
\end{center}

В частности, $\sigma(A)\subset[a,b]$. По теореме о спектральном
отображении $\sigma(p(A))=p(\sigma(A))\subset[0,\infty)$.
Для самосопряжённого $B$ со спектром в $[0,\infty)$
спектральное разложение даёт
\[
 B=\int_{\sigma(B)}\lambda\,\Pi(\dd\lambda),\qquad
 (Bv,v)=\int_{\sigma(B)}\lambda\,\Pi_{v,v}(\dd\lambda)\ge0.
\]
Действительно, $\lambda\ge0$ на спектре, а
$\Pi_{v,v}(S)=(\Pi(S)v,v)=\norm{\Pi(S)v}^2\ge0$.
Так как $p$ вещественный, $B=p(A)$ самосопряжён.
\end{proof}

\topic{Задача № 645}{\at{uLlyRG0VRnY}{832}{13:52}}
Пусть $A=A^*$ и $a\id\le A\le b\id$. Доказать непрерывность
$p\mapsto p(A)$ по норме $C[a,b]$.
\begin{proof}[Решение]
Оператор $p(A)$ нормален, поэтому его норма равна спектральному
радиусу. По предыдущей задаче $\sigma(A)\subset[a,b]$, откуда
\[
 \norm{p(A)}=\max_{\lambda\in\sigma(A)}\abs{p(\lambda)}
 \le\max_{\lambda\in[a,b]}\abs{p(\lambda)}=\norm p_{C[a,b]}.
\]
Применяем оценку к $p-q$ и получаем
$\norm{p(A)-q(A)}\le\norm{p-q}_{C[a,b]}$.
\end{proof}

\topic{Задачи № 646 и 647}{\at{uLlyRG0VRnY}{1122}{18:42}}
Унитарность, групповые тождества для $U(t)=\ee^{\ii tA}$ и формулу
$U'(t)=\ii AU(t)=\ii U(t)A$ получают из ряда экспоненты.

\topic{Непрерывная однопараметрическая группа}{\at{uLlyRG0VRnY}{1194}{19:54}}
\begin{exercise}[№~648* из сборника Кириллова и Гвишиани]
Пусть $U(t+s)=U(t)U(s)$, $U(t)^*=U(-t)$, $U(0)=\id$
и $t\mapsto U(t)$ непрерывно по операторной норме.
Доказать, что $U(t)=\ee^{\ii tA}$ с ограниченным
самосопряжённым $A$.
\end{exercise}
\begin{proof}[Решение]
Из $U(t)U(-t)=U(-t)U(t)=\id$ и $U(t)^*=U(-t)$ следует
унитарность $U(t)$. Представление $U(t)=\ee^{\ii A(t)}$ при
каждом $t$ само по себе не даёт коммутирования $A(t)$ и $A(-t)$
(\at{uLlyRG0VRnY}{1650}{27:30}).

Докажем дифференцируемость $U$ по указанию из сборника
(\at{uLlyRG0VRnY}{2056}{34:16}). Рассмотрим несобственный интеграл
\[
 R_n=\int_0^\infty\ee^{-s}U(s/n)\dd s
    =n\int_0^\infty\ee^{-nt}U(t)\dd t .
\]
Интеграл сходится, поскольку $\norm{U(s/n)}=1$. Вычитаем
$\id=\int_0^\infty\ee^{-s}\id\dd s$ и разбиваем интеграл в точке
$S>0$. Оценивая $\norm{U(s/n)-\id}\le2$ при $s\ge S$, получаем
\[
 \norm{R_n-\id}
 \le \sup_{0\le t\le S/n}\norm{U(t)-\id}+2\ee^{-S}.
\]
Второе слагаемое делаем сколь угодно малым выбором $S$.
При фиксированном $S$ первое стремится к нулю по непрерывности
$U$ в нуле. Значит, $R_n\to\id$ по норме.

Из группового тождества после замены $r=t+\tau$ следует
\[
 R_n U(\tau)=n\int_0^\infty\ee^{-nt}U(t+\tau)\dd t
            =n\ee^{n\tau}\int_\tau^\infty\ee^{-nr}U(r)\dd r .
\]
Правая часть дифференцируема по $\tau$, поскольку производная
интеграла по нижнему пределу равна $-\ee^{-n\tau}U(\tau)$.
Зафиксируем $n$ так, чтобы $\norm{\id-R_n}<1$. Тогда
\[
 R_n^{-1}=\sum_{k=0}^{\infty}(\id-R_n)^k.
\]
Ряд сходится по норме. При умножении суммы первых $m+1$ членов
на $R_n$ с любой стороны получаем $\id-(\id-R_n)^{m+1}\to\id$
(\at{uLlyRG0VRnY}{2586}{43:06}).
Умножая $R_nU(\tau)$ слева на $R_n^{-1}$, получаем
дифференцируемость $U(\tau)$.

Положим $B=\lim_{t\to0}(U(t)-\id)/t$. Тогда
\[
 (Bu,v)=\lim_{t\to0}\left(u,\frac{U(-t)-\id}{t}v\right)
       =-(u,Bv),
\]
Последнее равенство следует из замены $s=-t$. Значит,
$B^*=-B$ и $B=\ii A$, где $A=A^*$. Дифференцируя
$U(t+h)=U(h)U(t)$ по $h$ в нуле, получаем $U'(t)=\ii AU(t)$.
Единственное решение с $U(0)=\id$ имеет вид $U(t)=\ee^{\ii tA}$.
\end{proof}

\topic{Произведение матричных экспонент}{\at{uLlyRG0VRnY}{3206}{53:26}}
Формула Бейкера~--- Кэмпбелла~--- Хаусдорфа выражает
$C=\log(\ee^B\ee^A)$, то есть произведение в матричной группе
Ли $G$ в экспоненциальных координатах.

\begin{center}
\begin{minipage}{\linewidth}
\centering
\begin{tikzpicture}[>=Stealth]
 \draw (-2,0) -- (1.15,-.1) -- (2.1,1.12) -- (-.98,1.1) -- cycle;
 \draw (-2,-.48) .. controls (-1.05,-.3) and (-.08,-.6) .. (.93,-1.1)
   .. controls (1.28,-.68) and (1.81,-.42) .. (2.24,-.36)
   .. controls (1.92,-.02) and (1.81,.17) .. (1.73,.36);
 \draw (-2,-.48) -- (-1.76,-.01);
 \draw (-.64,.54) .. controls (-.3,.91) and (.53,.9) .. (.83,.57)
   .. controls (1.18,.18) and (.32,.1) .. (-.31,.17)
   .. controls (-.62,.22) and (-.89,.32) .. cycle;
 \draw (-1.04,-.11) .. controls (-1.34,-.41) and (-.29,-.63) .. (.36,-.4)
   .. controls (.63,-.29) and (.76,-.2) .. (.82,-.08);
 \draw[->] (-.73,.38) .. controls (-1.03,.22) and (-1.1,.02) .. (-1.04,-.11);
 \draw[->] (.92,.59) .. controls (1.24,.54) and (1.32,.3) .. (.82,-.08);
 \fill (-.08,.48) circle (1.2pt);
 \node[above] at (-.08,.48) {$0$};
 \fill (-.08,-.29) circle (1.2pt);
 \node[right] at (-.08,-.29) {$E$};
 \node[above] at (1.56,1.12) {$T_EG=\mathfrak g$};
 \node[right] at (1.8,-.53) {$G$};
\end{tikzpicture}
\captionof{figure}{Экспонента переводит окрестность нуля в $T_EG$
в окрестность единицы $E$ в $G$ (\at{uLlyRG0VRnY}{3403}{56:43}).}
\end{minipage}
\end{center}

\begin{exercise}
Если $\ee^C=\ee^B\ee^A$, доказать
$\ee^{\ad_C}=\ee^{\ad_B}\ee^{\ad_A}$.
В локальной задаче о логарифме
$\ee^{D(t)}=\ee^{tB}\ee^A$, $D(0)=A$,
выразить $D'(t)$ и $C=D(1)$.
\end{exercise}
Из $\ad_B=L_B-R_B$ получаем
(\at{uLlyRG0VRnY}{3442}{57:22})
\[
 f(L_B,R_B)X=f(\overset3 B,\overset1 B)\overset2 X,
 \qquad
 \ee^{\ad_B}X=\ee^{\overset3 B-\overset1 B}\overset2 X.
\]
Во внешней экспоненте ставим индексы $4$ и $0$. Тогда для
\[
 F(x,y,z,w,\alpha)=\ee^{x-y}\ee^{z-w}\alpha
                 =\ee^{x+z}\ee^{-(y+w)}\alpha
\]
имеем
\[
 \begin{aligned}
  \ee^{\ad_B}\ee^{\ad_A}X
    &=F(\overset4 B,\overset0 B,\overset3 A,\overset1 A,\overset2 X)\\
    &=\overset3{\boxed{\ee^{\overset4 B+\overset3 A}}}
      \overset1{\boxed{\ee^{-(\overset0 B+\overset1 A)}}}\overset2 X\\
    &=\ee^{\overset3 C}\ee^{-\overset1 C}\overset2 X
      =\ee^{\ad_C}X.
 \end{aligned}
\]
По правилу для разложимых символов первая скобка даёт
$\ee^B\ee^A=\ee^C$, вторая~--- $\ee^{-A}\ee^{-B}=\ee^{-C}$.

Из $\ee^{D(t)}=\ee^{tB}\ee^A$ получаем
$C=D(1)=A+\int_0^1D'(t)\dd t$
(\at{uLlyRG0VRnY}{4188}{1:09:48}). Дифференцируем равенство
для экспонент и умножаем справа на $\ee^{-D(t)}$. Пишем $D=D(t)$.
\[
 \begin{aligned}
  B&=\overset{2}{D'}\,
        \frac{\ee^{\overset3 D}-\ee^{\overset1 D}}
             {\overset3 D-\overset1 D}\,\ee^{-\overset0 D}\\
   &=\overset{2}{D'}\,
        \frac{\ee^{\overset3 D-\overset0 D}-\ee^{\overset1 D-\overset0 D}}
             {\overset3 D-\overset1 D}\\
   &=\overset{2}{D'}\,
        \frac{\ee^{\overset3 D-\overset1 D}-1}
             {\overset3 D-\overset1 D}
     =f(\ad_D)D',\qquad f(x)=\frac{\ee^x-1}{x}.
 \end{aligned}
\]
В последней строке склеены соседние индексы $0$ и $1$
(\at{uLlyRG0VRnY}{4696}{1:18:16}). Для достаточно малых $A,B$
обращаем символ $f$ и используем ветвь логарифма у единицы.
\[
 D'=g(\ad_D)B=h(\ee^{\ad_D})B,\qquad
 g(x)=\frac{x}{\ee^x-1},\qquad h(y)=\frac{\log y}{y-1}.
\]
Здесь $f(0)=g(0)=h(1)=1$. По первой части
$\ee^{\ad_{D(t)}}=\ee^{t\ad_B}\ee^{\ad_A}$, откуда
\[
 C=A+\int_0^1 h\!\left(\ee^{t\ad_B}\ee^{\ad_A}\right)B\,\dd t .
\]
Для разложения по коммутаторам используем ряд
\[
 h(y)=\sum_{n=0}^{\infty}\frac{(-1)^n}{n+1}(y-1)^n,
 \qquad \abs{y-1}<1 .
\]

\topic{Домашнее задание}{\at{uLlyRG0VRnY}{5060}{1:24:20}}
\begin{enumerate}
 \item Собрать коммутаторы первого и второго порядка в ряде для $C$.
 \item Решить задачи 1.4 и 1.11 из «Задачника по операторным
 методам» Карасёва.
\end{enumerate}

\sessionsection{Семинар 3. Функции от нескольких операторов}
\record{18 сентября}{ZDZiS1ShThI}

\topic{Формулы Фурье и Коши}{\at{ZDZiS1ShThI}{67}{1:07}}
\begin{exercise}[Задача 1.4]
Пусть гладкая $f$ задана в окрестности
$K=\sigma(A_1)\times\cdots\times\sigma(A_n)$, а
$\rho\in C_c^\infty(\RR^n)$ равна единице в окрестности $K$ и $\supp \rho \subseteq \Dom(f)$.
Для $g=\rho f$ доказать независимость формулы Фурье от $\rho$.
При голоморфной $f$ сравнить её с формулой Коши на полиобласти,
то есть на декартовом произведении плоских областей.
Для самосопряжённых $A_j$ выразить результат через спектральные меры.
\end{exercise}
В начале уточняют условие на производящие элементы
(\at{ZDZiS1ShThI}{119}{1:59}). Операторные экспоненты растут
не быстрее многочлена по вещественному параметру. Условие нужно
для интеграла с быстро убывающей $\widehat g$.
\[
 \begin{aligned}
 \widehat g(x)&=\frac{1}{(2\pi)^{n/2}}
       \int_{\RR^n}g(\lambda)\ee^{\ii x\cdot\lambda}\dd\lambda,\\
 T(g)&=\frac{1}{(2\pi)^{n/2}}\int_{\RR^n}
   \widehat g(x)\ee^{-\ii x_n A_n}\cdots\ee^{-\ii x_1 A_1}\dd x .
 \end{aligned}
\]
Пусть $\Gamma_j$~--- положительно ориентированный контур,
охватывающий $\sigma(A_j)$. По теореме~\ref{thm:holomorphic-functional-calculus}
о голоморфном функциональном исчислении
(\at{ZDZiS1ShThI}{177}{2:57})
\[
 \ee^{-\ii x_jA_j}=\frac{1}{2\pi\ii}
 \int_{\Gamma_j}\ee^{-\ii x_j\lambda_j}
            (\lambda_j\id-A_j)^{-1}\dd\lambda_j .
\]
Перемножая эти формулы в порядке от $n$ к $1$, получаем
\[
 \begin{aligned}
 &\ee^{-\ii x_nA_n}\cdots\ee^{-\ii x_1A_1}\\
 &\quad=\frac{1}{(2\pi\ii)^n}
   \int_{\Gamma_n}\!\cdots\!\int_{\Gamma_1}
   \ee^{-\ii x\cdot\lambda}
   (\lambda_n\id-A_n)^{-1}\cdots(\lambda_1\id-A_1)^{-1}
   \dd\lambda_1\cdots\dd\lambda_n .
 \end{aligned}
\]
Подставим произведение в $T(g)$ и вынесем контурные интегралы
наружу (\at{ZDZiS1ShThI}{525}{8:45}). Перед произведением
резольвент остаётся выражение
\[
 \frac{1}{(2\pi)^{n/2}}\int_{\RR^n}
              \widehat g(x)\ee^{-\ii x\cdot\lambda}\dd x,
\]
Заменим его на $g(\lambda)$ и получим
\[
 T(g)=\frac{1}{(2\pi\ii)^n}
 \int_{\Gamma_n}\!\cdots\!\int_{\Gamma_1}g(\lambda)
 (\lambda_n\id-A_n)^{-1}\cdots(\lambda_1\id-A_1)^{-1}
 \dd\lambda_1\cdots\dd\lambda_n .
\]

\topic{Независимость от срезающей функции}{\at{ZDZiS1ShThI}{752}{12:32}}
Для двух срезающих функций $\rho_1,\rho_2$ покажем, что $T((\rho_1-\rho_2)f)=0$. Обе функции равны единице
в некоторой общей окрестности $K$, поэтому там $(\rho_1-\rho_2)f=0$.
Приблизим $\Gamma_j$ к $\sigma(A_j)$,
сохраняя спектр внутри. Если $\Gamma_1 \times \ldots \times \Gamma_n$ лежит в
окрестности, где $(\rho_1-\rho_2)f=0$, $T((\rho_1-\rho_2)f)=0$.

\begin{center}
\begin{minipage}{\linewidth}
\centering
\begin{tikzpicture}[>=Stealth]
 \draw (-1.13,-.53) .. controls (-1.65,.35) and (-.93,1.29) .. (.08,1.34)
   .. controls (1.08,1.4) and (1.45,.48) .. (1.33,-.26)
   .. controls (1.23,-1.08) and (.52,-1.34) .. (-.19,-1.29)
   .. controls (-.56,-1.27) and (-.93,-.92) .. cycle;
 \draw (-.94,-.46) .. controls (-1.37,.29) and (-.76,1.11) .. (.08,1.12)
   .. controls (.88,1.18) and (1.23,.43) .. (1.13,-.25)
   .. controls (1.03,-.9) and (.38,-1.14) .. (-.17,-1.08)
   .. controls (-.51,-1.03) and (-.77,-.78) .. cycle;
 \draw (-.74,-.36) .. controls (-1.05,.1) and (-.71,.78) .. (-.19,.89)
   .. controls (.14,1.01) and (.69,.83) .. (.8,.51)
   .. controls (1.09,.27) and (.96,-.48) .. (.58,-.71)
   .. controls (.08,-.98) and (-.49,-.8) .. cycle;
 \draw[->] (.44,-1.26) .. controls (.77,-1.19) and (1.02,-.99) .. (1.19,-.67);
 \fill (-.18,.18) circle (1.2pt) (.18,-.13) circle (1.2pt);
 \node at (.04,-.48) {$\sigma(A_j)$};
 \node[above] at (.1,1.34) {$\Gamma_j$};
\end{tikzpicture}
\captionof{figure}{Контуры вокруг $\sigma(A_j)$
(\at{ZDZiS1ShThI}{912}{15:12}).}
\end{minipage}
\end{center}

При голоморфной $f$ берут срезающую функцию с носителем в
полиобласти голоморфности, а произведение контуров~--- в
окрестности $K$, где $\rho=1$. Тогда под интегралом вместо
$g$ стоит $f$, как в формуле Коши
(\at{ZDZiS1ShThI}{1125}{18:45}).
\begin{remark}
В записи остаётся пробел. Быстрое убывание $\widehat g$ обеспечивает
сходимость исходного $T(g)$ при полиномиальном росте операторных
экспонент. На комплексном контуре $\ee^{-\ii x\cdot\lambda}$ может
расти экспоненциально, а обращение Фурье восстанавливает $g$ на
$\RR^n$. Для перехода к контурной записи требуется обоснование.
\end{remark}

\topic{Самосопряжённые аргументы}{\at{ZDZiS1ShThI}{1210}{20:10}}
Пусть $E^{(j)}$~--- спектральная мера самосопряжённого $A_j$.
Тогда
\[
 \ee^{-\ii x_j A_j}=\int_{\RR}\ee^{-\ii x_j\psi_j}\dd E^{(j)}(\psi_j).
\]
Подставляем спектральные интегралы в порядке от $n$ к $1$.
Обращение Фурье применяем при вещественных $\psi_j$.
\[
 \begin{aligned}
 T(\rho f)&=\frac{1}{(2\pi)^{n/2}}\int_{\RR^n}\widehat{\rho f}(x)
 \left[\int_{\RR^n}\ee^{-\ii x\cdot\psi}
       \dd E^{(n)}(\psi_n)\cdots\dd E^{(1)}(\psi_1)\right]\dd x\\
 &=\int_{\RR^n}(\rho f)(\psi)\dd E^{(n)}(\psi_n)\cdots\dd E^{(1)}(\psi_1)\\
 &=\int_K f(\psi)\dd E^{(n)}(\psi_n)\cdots\dd E^{(1)}(\psi_1).
 \end{aligned}
\]
В последнем равенстве используем $\supp E^{(j)}=\sigma(A_j)$
и $\rho=1$ на $K$.

\topic{Домашнее задание}{\at{ZDZiS1ShThI}{1701}{28:21}}
Задачи 1.4 и 1.11 из задачника Карасёва, с.~33--34.

\sessionsection{Семинар 4. Производная функции от оператора и замыкаемость}
\record{25 сентября}{RgVedijsA-I}

\topic{Дифференцирование синуса}{\at{RgVedijsA-I}{85}{1:25}}
\begin{exercise}[Задача 1.4, пункт 4]
Пусть $\ad_A=L_A-R_A$, $\operatorname{an}_A=L_A+R_A$ и
$\rho(u)=\sin u/u$, $\rho(0)=1$. Доказать равенство
\[
 \mathcal D(\sin A)=\cos\!\left(\frac{\operatorname{an}_A}{2}\right)
      \rho\!\left(\frac{\ad_A}{2}\right)\mathcal DA.
\]
\end{exercise}
\begin{proof}[Решение]
По формуле дифференцирования сложной функции
\[
 \mathcal D(\sin A)=\ddiff{\sin}(\overset3 A,\overset1 A)
                     \overset2{\mathcal DA}.
\]
Запишем правую часть с теми же индексами. Для любого $X$
\[
 f(\ad_A)X=f(\overset3 A-\overset1 A)\overset2 X,
 \qquad
 f(\operatorname{an}_A)X=f(\overset3 A+\overset1 A)\overset2 X.
\]
Поскольку $L_A$ и $R_A$ коммутируют, получаем
\[
 \cos\!\left(\frac{\operatorname{an}_A}{2}\right)
      \rho\!\left(\frac{\ad_A}{2}\right)\mathcal DA
 =\cos\!\left(\frac{\overset3 A+\overset1 A}{2}\right)
      \rho\!\left(\frac{\overset3 A-\overset1 A}{2}\right)
      \overset2{\mathcal DA}.
\]
Символ справа совпадает с разностной производной синуса по тождеству
из первого пункта задачи
\[
 \frac{\sin x-\sin y}{x-y}
 =\cos\!\left(\frac{x+y}{2}\right)\rho\!\left(\frac{x-y}{2}\right).
\]
\end{proof}
Пятый пункт предлагают решить тем же способом
(\at{RgVedijsA-I}{480}{8:00}).

\topic{Формула для нормальной матрицы}{\at{RgVedijsA-I}{522}{8:42}}
\noindent\textbf{Задача 1.11, пункт 4.}
Пусть $A=\sum_{j=1}^m\lambda_jE_j$~--- нормальная матрица,
$\lambda_j$~--- её различные собственные значения, $E_j$~---
спектральные проекторы, $f$ непрерывна в окрестности спектра и
$\mathcal D$~--- дифференцирование алгебры матриц. Доказать, что
\[
 \mathcal D(f(A))=\sum_{i\ne j}
   \frac{f(\lambda_i)-f(\lambda_j)}{\lambda_i-\lambda_j}
   E_i\mathcal D(A)E_j .
\]
\begin{proof}[Решение]
Дифференцирование матричной алгебры внутреннее, поэтому
$\mathcal D(X)=[M,X]$. На семинаре это выводят из теоремы
о дифференцированиях полупростых алгебр Ли в характеристике нуль
(\at{RgVedijsA-I}{1307}{21:47}). В разложении
$\mathfrak{gl}_n=\CC\id\oplus\mathfrak{sl}_n$ дифференцирование
обращает $\id$ в нуль и сохраняет
$\mathfrak{sl}_n=[\mathfrak{gl}_n,\mathfrak{gl}_n]$, поскольку
подчиняется правилу Лейбница для коммутатора.
На $\mathfrak{sl}_n$ оно имеет вид $[M,\cdot]$ по теореме.

Поскольку $\id=\sum_iE_i$, левую часть раскладываем по блокам
$E_i\mathcal D(f(A))E_j$. Подставляем $f(A)=\sum_k f(\lambda_k)E_k$.
\[
 \begin{aligned}
 E_i\mathcal D(f(A))E_j
 &=\sum_k f(\lambda_k)(E_iME_kE_j-E_iE_kME_j)\\
 &=(f(\lambda_j)-f(\lambda_i))E_iME_j.
 \end{aligned}
\]
Точно так же $E_i\mathcal D(A)E_j=(\lambda_j-\lambda_i)E_iME_j$.
Тогда \[
\begin{aligned}
\mathcal D(f(A))&=\sum_{i\ne j}E_i\mathcal D(f(A))E_j\\&=\sum_{i\ne j}(f(\lambda_j)-f(\lambda_i))E_iME_j\\&=\sum_{i\ne j}
   \frac{f(\lambda_i)-f(\lambda_j)}{\lambda_i-\lambda_j}
   E_i\mathcal D(A)E_j.
\end{aligned}
\]
\end{proof}
В конце отмечают, что формулу также можно получить из правила
дифференцирования сложной функции (\at{RgVedijsA-I}{1399}{23:19}).

\Needspace{8\baselineskip}
\topic{Значение в точке как оператор}{\at{RgVedijsA-I}{1487}{24:47}}
\begin{exercise}
В $L^2[0,1]$ рассмотрим оператор $Tf=f(0)\mathbf 1$ на
$\Dom(T)=C[0,1]$, где $\mathbf 1$~--- постоянная функция.
Доказать, что $T$ незамыкаем, и найти $T^*$.
\end{exercise}
\begin{proof}[Решение]
На языке последовательностей замыкаемость означает, что для
$f_n\in\Dom(T)$ из $f_n\to0$ и $Tf_n\to y$ следует $y=0$.
Возьмём $f_n(t)=1-nt$ при $0\le t\le 1/n$ и $f_n(t)=0$ при $t>1/n$.

\begin{center}
\begin{minipage}{\linewidth}
\centering
\begin{tikzpicture}[>=Stealth,scale=.85]
 \draw[->] (-.23,0)--(3.14,0) node[right] {$t$};
 \draw[->] (0,-.26)--(0,1.51);
 \draw[thick] (0,1.13)--(.88,0)--(2.79,0);
 \draw (-.071,1.13)--(.071,1.13) node[left=4pt] {$1$};
 \draw (.88,-.067)--(.88,.067) node[below=5pt] {$1/n$};
 \draw (2.79,-.067)--(2.79,.067) node[below=5pt] {$1$};
\end{tikzpicture}
\captionof{figure}{График $f_n$ (\at{RgVedijsA-I}{1744}{29:04}).}
\end{minipage}
\end{center}

Тогда $\norm{f_n}_2^2=\int_0^{1/n}(1-nt)^2\dd t=1/(3n)\to0$,
а $Tf_n=\mathbf 1$. Следовательно, $(0,\mathbf 1)\in\overline{G(T)}$
и $T$ незамыкаем.

Для $g\in\Dom(T^*)$ существует $c$ такое, что для всех $f\in C[0,1]$
\[
 \abs{(Tf,g)}=\abs{f(0)}\left|\int_0^1\overline{g(t)}\dd t\right|
              \le c\norm f_2 .
\]
Подставляем $f_n$ и переходим к пределу. Так как $f_n(0)=1$,
получаем $(\mathbf 1,g)=0$. При этом условии $(Tf,g)=0$ для всех
$f\in\Dom(T)$, поэтому $T^*g=0$. Таким образом,
$\Dom(T^*)=(\CC\mathbf 1)^\perp$~--- неплотное подпространство
коразмерности~1.
\end{proof}

\topic{Дифференцирование в $C[0,1]$}{\at{RgVedijsA-I}{2250}{37:30}}
Замкнутость определяется и в банаховом пространстве через график
в $X\oplus X$. Пусть $X=C[0,1]$ с равномерной нормой и
$Tf=f'$ на $\Dom(T)=C^1[0,1]$. Если $f_n\to f$ и $f_n'\to g$
равномерно, то по теореме о почленном дифференцировании
$f\in C^1[0,1]$ и $f'=g$. Следовательно, $T$ замкнут.

\topic{Дифференцирование в $L^2[0,1]$}{\at{RgVedijsA-I}{2527}{42:07}}
\noindent\textbf{Задача.} Проверить замкнутость оператора $Tx=x'$
в $L^2[0,1]$ на области
\[
 \Dom(T)=\{x\in AC[0,1]: x'\in L^2[0,1]\}.
\]
Здесь $AC[0,1]$~--- класс абсолютно непрерывных функций.
\begin{proof}[Решение]
Пусть $x_n\to x$ и $x_n'\to y$ в $L^2$. Для абсолютно непрерывных
$x_n$ выполнена формула Ньютона~--- Лейбница
$x_n(t)=x_n(0)+\int_0^t x_n'(s)\dd s$.
По неравенству Коши~--- Буняковского для всех $t\in[0,1]$
\[
 \left|\int_0^t(x_n'(s)-y(s))\dd s\right|
 \le\int_0^1\abs{x_n'(s)-y(s)}\dd s\le\norm{x_n'-y}_2\longrightarrow0.
\]
Значит, интегралы сходятся равномерно. Из той же формулы
Ньютона~--- Лейбница
\[
 x_n(0)-x_m(0)=x_n(t)-x_m(t)-\int_0^t(x_n'(s)-x_m'(s))\dd s.
\]
Берём норму $L^2[0,1]$ обеих частей.
\[
 \abs{x_n(0)-x_m(0)}
 \le\norm{x_n-x_m}_2+
    \left\|\int_0^{\,\cdot}(x_n'-x_m')\dd s\right\|_2
 \longrightarrow0.
\]
Первое слагаемое стремится к нулю по фундаментальности $x_n$
в $L^2$, второе~--- по равномерной сходимости интегралов.
Пусть $x_n(0)\to\alpha$. Тогда $x_n(t)\to\alpha+\int_0^t y(s)\dd s$
равномерно и, следовательно, в $L^2$. По единственности предела
эта функция представляет $x$. Она абсолютно непрерывна и имеет
производную $y$ почти всюду. Значит, $x\in\Dom(T)$ и $Tx=y$.
\end{proof}

Конкретные домашние задачи в записи не названы. Преподаватель
сказал, что пришлёт их по почте (\at{RgVedijsA-I}{3802}{1:03:22}).

\sessionsection{Семинар 5. Сопряжение производной и дефект на полупрямой}
\record{2 октября}{2s-Uo0FKoAE}

\topic{Сопряжение на отрезке}{\at{2s-Uo0FKoAE}{14}{0:14}}
\begin{exercise}
Найти сопряжённый к $Tf=f'$ в $L^2(0,1)$ при
$\Dom(T)=H^1(0,1)$, то есть на абсолютно непрерывных функциях
с производной из $L^2(0,1)$.
\end{exercise}
Замкнутость $T$ доказана на предыдущем семинаре. Если $g\in\Dom(T^*)$,
то подстановка финитных $f$ в $(f',g)=(f,T^*g)$ даёт
$g'=-T^*g$ в смысле распределений. В частности, $g\in H^1(0,1)$.
Интегрируем по частям:
\[
 (Tf,g)=f(1)\overline{g(1)}-f(0)\overline{g(0)}-(f,g').
\]
Непрерывность функционала $f\mapsto(Tf,g)$ по норме $L^2$
равносильна исчезновению граничного члена. Действительно,
при подстановке $f_n(t)=\max(1-nt,0)$, а затем $f_n(1-t)$
$L^2$-норма стремится к нулю, а значение на одном конце равно $1$
и на другом~--- $0$. Поэтому $g(0)=g(1)=0$ и
\[
 T^*g=-g',\qquad \Dom(T^*)=\{g\in H^1(0,1):g(0)=g(1)=0\}.
\]
Эта область плотна в $L^2(0,1)$, так как содержит $\sch{(0,1)}$.
На вопрос о знаке в формуле для $T^*$ приводят пример из квантовой
механики, где берут дифференциальное выражение $-\ii\,\mathrm d/\mathrm dt$
(\at{2s-Uo0FKoAE}{360}{6:00}). При формальном сопряжении оно сохраняется.

\topic{Минимальный оператор на полупрямой}{\at{2s-Uo0FKoAE}{489}{8:09}}
\begin{exercise}
Пусть $T_0=-\ii\,\mathrm d/\mathrm dt$ на $\sch{(0,\infty)}$
в $L^2(0,\infty)$ и $T=\overline{T_0}$. Найти $T$, $T^*$,
индексы дефекта и поле регулярности $\Pi(T)$.
\end{exercise}
\paragraph*{Сопряжённый оператор.}
По определению $u\in\Dom(T_0^*)$ тогда и только тогда, когда
существует $h\in L^2(0,\infty)$ такое, что для всех
$v\in\sch{(0,\infty)}$ выполнено равенство
\[
 (T_0v,u)=(v,h),\qquad
 -\ii\int_0^\infty v'(t)\overline{u(t)}\dd t
 =\int_0^\infty v(t)\overline{h(t)}\dd t.
\]
Это определение обобщённой производной $u'=\ii h$. Поэтому
$T_0^*u=h=-\ii u'$, а $\Dom(T_0^*)$ состоит из функций $u\in L^2(0,\infty)$,
у которых обобщённая производная принадлежит $L^2(0,\infty)$
(\at{2s-Uo0FKoAE}{640}{10:40}). Такие функции имеют абсолютно
непрерывных представителей на каждом $[0,b]$. Обозначаем
эту область через $H^1(0,\infty)$.

\begin{lemma}[Значение на бесконечности]
Если $u\in H^1(0,\infty)$, то $u(t)\to0$ при $t\to+\infty$
(\at{2s-Uo0FKoAE}{947}{15:47}).
\end{lemma}
\begin{proof}
По формуле Ньютона~--- Лейбница
\[
 \abs{u(t)}^2=\abs{u(0)}^2+
 \int_0^t\bigl(u'(s)\overline{u(s)}+u(s)\overline{u'(s)}\bigr)\dd s.
\]
Так как $u,u'\in L^2(0,\infty)$, при $t\to\infty$ интегралы
от двух слагаемых стремятся к $(u',u)$ и $(u,u')$ соответственно.
Следовательно, $\abs{u(t)}^2$ имеет предел. Из $u\in L^2(0,\infty)$
следует, что этот предел равен нулю.
\end{proof}

\paragraph*{Замыкание.}
Берём $u,v\in\Dom(T_0^*)$ и интегрируем по частям. По лемме
$u(t),v(t)\to0$ при $t\to\infty$, поэтому остаётся только
значение на левом конце (\at{2s-Uo0FKoAE}{1265}{21:05}).
\[
 \begin{aligned}
 (T_0^*u,v)&=\int_0^\infty(-\ii u'(t))\overline{v(t)}\dd t\\
 &=-\ii\bigl[u(t)\overline{v(t)}\bigr]_0^\infty
   +\ii\int_0^\infty u(t)\overline{v'(t)}\dd t\\
 &=\ii u(0)\overline{v(0)}+(u,T_0^*v).
 \end{aligned}
\]
Для финитных $u,v$ граничный член равен нулю. Следовательно,
$(T_0u,v)=(u,T_0v)$, то есть $T_0$ симметричен и $T_0\subset T_0^*$
(\at{2s-Uo0FKoAE}{1670}{27:50}). Оператор $T_0^*$ замкнут, поэтому
он содержит замыкание $T_0$. У $T_0$ и $\overline{T_0}$ один и тот же
сопряжённый оператор. По теореме~\ref{thm:closability-via-adjoint}
получаем
\[
 T_0\subset T=\overline{T_0}=T_0^{**}\subset T_0^*=T^*.
\]

Если $v\in\Dom(T)$, то по включению $T\subset T^*$ имеем
$v\in\Dom(T^*)$ и $Tv=T^*v$. По определению двойного сопряжённого
для каждого $u\in\Dom(T^*)$ выполнено
\[
 (T^*u,v)=(u,Tv)=(u,T^*v).
\]
Сравнивая с формулой интегрирования по частям, получаем
$\ii u(0)\overline{v(0)}=0$ при всех $u\in\Dom(T^*)$. Значение
$u(0)$ на $\Dom(T^*)$ произвольно, поэтому $v(0)=0$.
Наоборот, если $v\in\Dom(T^*)$ и $v(0)=0$, то та же формула даёт
$(T^*u,v)=(u,T^*v)$ для всех $u\in\Dom(T^*)$. Поэтому
$v\in\Dom((T^*)^*)=\Dom(T)$ (\at{2s-Uo0FKoAE}{2025}{33:45}). Итак,
\[
 Tv=-\ii v',\qquad \Dom(T)=\{v\in H^1(0,\infty):v(0)=0\}.
\]
У $T^*$ область шире, так что $T$ симметричен, но не самосопряжён.

\topic{Поле регулярности на полупрямой}{\at{2s-Uo0FKoAE}{2071}{34:31}}
Для $u\in\Dom(T)$ имеем $u(0)=0$ и $u(t)\to0$ при $t\to\infty$.
По формуле интегрирования по частям $(Tu,u)=(u,Tu)$, поэтому
числовая область значений вещественна,
и по лемме~\ref{lem:numerical-range-regularity} обе открытые
полуплоскости входят в $\Pi(T)$.

Дефектное число в точке регулярности показывает, какова коразмерность
области, на которой задан $(T-\lambda\id)^{-1}$. По
следствию~\ref{cor:regularity-decomposition} индексы дефекта равны
\[
 d_+(T)=\dim\ker(T^*+\ii\id),\qquad
 d_-(T)=\dim\ker(T^*-\ii\id).
\]
\Needspace{5\baselineskip}
Для $d_+(T)$ решаем уравнение для обобщённой производной
\[
 -\ii u'+\ii u=0,\qquad u(t)=C\ee^t.
\]
Условие $u\in L^2(0,\infty)$ требует $C=0$, поэтому $d_+(T)=0$.
Для $d_-(T)$ получаем
\[
 -\ii u'-\ii u=0,\qquad u(t)=C\ee^{-t}.
\]
Все решения принадлежат $\Dom(T^*)$ и образуют одномерное
подпространство. Значит, $d_-(T)=1$.

Если бы $\lambda_0\in\RR$ лежала в $\Pi(T)$, то вместе с ней
в $\Pi(T)$ лежал бы малый диск с центром $\lambda_0$.
По теореме~\ref{thm:krein-krasnoselskii} на диске дефектное число
постоянно. Но в верхней половине оно равно $0$, а в нижней~--- $1$.
Значит, $\Pi(T)=\CC\setminus\RR$.

\begin{center}
\begin{minipage}{\linewidth}
\centering
\begin{tikzpicture}[>=Stealth]
 \draw[->] (-2.5,0)--(2.08,0) node[right] {$\operatorname{Re}\lambda$};
 \draw[->] (0,-1.3)--(0,1.3) node[above] {$\operatorname{Im}\lambda$};
 \draw (-1.27,0) circle (.59);
 \fill (-1.27,0) circle (1.25pt);
 \node[below] at (-1.27,0) {$\lambda_0$};
 \node at (.98,.61) {$d_+=0$};
 \node at (.98,-.61) {$d_-=1$};
\end{tikzpicture}
\captionof{figure}{Диск около вещественной точки $\lambda_0$
пересекает полуплоскости с разными дефектами
(\at{2s-Uo0FKoAE}{2806}{46:46}).}
\end{minipage}
\end{center}

Если индексы совпадают, этого рассуждения недостаточно. Тогда
исследуют числовую область значений. Например, при полуограниченности
плоскость с вырезанным лучом входит в поле регулярности и соединяет
обе полуплоскости (\at{2s-Uo0FKoAE}{2656}{44:16}).

\topic{Домашнее задание}{\at{2s-Uo0FKoAE}{2906}{48:26}}
\begin{enumerate}
\item Для $T_0=-\ii\,\mathrm d/\mathrm dt$ на $\sch{\RR}$
в $L^2(\RR)$ найти поле регулярности и индексы дефекта замыкания.
\item Решить ту же задачу для $T_0=-\ii\,\mathrm d/\mathrm dt$ на
$\sch{(0,1)}$ в $L^2(0,1)$.
\item Задача 3.7 из задачника Карасёва, с.~26
(\at{2s-Uo0FKoAE}{3137}{52:17}). Выразить разность
\[
 f(\overset1 A,\overset2 B)-f(\overset2 A,\overset1 B).
\]
\end{enumerate}

\sessionsection{Семинар 6. Поле регулярности. Пустой спектр}
\record{9 октября}{9TaAAWmVph8}

\topic{Минимальная производная на отрезке и на прямой}{\at{9TaAAWmVph8}{64}{1:04}}
\begin{exercise}
Для замыкания оператора $-\ii\,d/dt$ на финитных гладких
функциях найти на отрезке $[-1,1]$ сопряжённый оператор,
индексы дефекта и поле регулярности. Найти поле регулярности
для того же оператора на прямой.
\end{exercise}

В $L^2[-1,1]$ оператор $T_0=-\ii\,d/dt$ на
$C_c^\infty(-1,1)$ симметричен и плотно определён.
Положим $T=\overline{T_0}$. Тогда
\[
 T_0\subset T=T_0^{**}\subset T_0^*,\qquad T^*=T_0^*.
\]
Вектор $y\in L^2[-1,1]$ принадлежит $\Dom(T_0^*)$,
если $(T_0x,y)$~--- непрерывный линейный функционал от $x$.
По теореме Рисса это означает, что существует $h\in L^2[-1,1]$
такое, что
\[
 \int_{-1}^1-\ii x'(t)\overline{y(t)}\dd t
 =\int_{-1}^1 x(t)\overline{h(t)}\dd t
 \qquad(x\in C_c^\infty(-1,1)).
\]
Таким образом, $h=-\ii y'$ в смысле обобщённой производной и
\[
 \Dom(T_0^*)=H^1(-1,1)
 =\{f\in AC[-1,1]:f'\in L^2[-1,1]\}.
\]
При втором сопряжении интегрирование по частям даёт
граничный член $-\ii[x(t)\overline{y(t)}]_{-1}^{1}$.
Значения $x$ на концах произвольны, а функционал
от этих значений не ограничен в норме $L^2$, если
$y$ не обращается в нуль на обоих концах. Поэтому
\[
 \Dom(T)=
 \{y\in\Dom(T_0^*):y(-1)=y(1)=0\},\qquad Ty=-\ii y'.
\]

Для $d_\pm=\dim\ker(T^*\pm\ii\id)$ решаем уравнения
$y'\mp y=0$. Решения~--- кратные $\ee^{\pm t}$,
и обе экспоненты лежат в $L^2[-1,1]$. Значит, $d_+=d_-=1$.

\Needspace{7\baselineskip}
\textbf{Поле регулярности на отрезке.}
Обе открытые полуплоскости лежат в $\Pi(T)$ в силу
симметричности. Для $\lambda\in\RR$ подставляем
$x(t)=\ee^{\ii\lambda t}y(t)$ в оценку
$\norm{(T-\lambda\id)x}_2\ge c\norm{x}_2$.
Тогда
\[
 -\ii x'(t)-\lambda x(t)=\ee^{\ii\lambda t}(-\ii y'(t)),\qquad
 \norm{(T-\lambda\id)x}_2=\norm{y'}_2,\quad
 \norm{x}_2=\norm{y}_2.
\]
Здесь $y$ удовлетворяет тем же граничным условиям,
а $\bigl|\ee^{\ii\lambda t}\bigr|=1$.
Остаётся доказать неравенство Фридрихса
$\norm{y'}_2\ge c\norm{y}_2$ (\at{9TaAAWmVph8}{1271}{21:11}).

Для ненулевой константы неравенство не выполнено.
Пусть $f\in\Dom(T)$. По формуле Ньютона--Лейбница
$f(t)=\int_{-1}^t f'(\tau)\dd\tau$.
К квадрату модуля применяем неравенство Коши--Буняковского
на квадрате $[-1,t]^2$. Его площадь не превосходит $4$, поэтому
\[
 \begin{aligned}
 |f(t)|^2
 &=\int_{-1}^t\!\!\int_{-1}^t f'(\tau)\overline{f'(s)}
     \dd\tau\,\dd s\\
 &\le 4\left(
   \int_{-1}^t\!\!\int_{-1}^t |f'(\tau)|^2|f'(s)|^2
   \dd\tau\,\dd s\right)^{1/2}
 =4\int_{-1}^t|f'(\tau)|^2\dd\tau
 \le4\norm{f'}_2^2.
 \end{aligned}
\]
Получили оценку равномерной нормы через норму производной
для вложения в $C[-1,1]$.
Интегрируя по $t$ от $-1$ до $1$, получаем
$\norm f_2^2\le8\norm{f'}_2^2$. Следовательно, $\Pi(T)=\CC$.

\textbf{Поле регулярности на прямой.}
Для $\lambda\in\RR$ та же подстановка сводит вопрос
к оценке $\norm{y'}_2\ge c\norm{y}_2$.
Берём финитные гладкие функции $y_n$, равные единице на
$[-n,n]$ и нулю вне $[-n-1,n+1]$ (рис.~\ref{fig:sem6-cutoff}).
При изменении $n$ участки возрастания и убывания только
сдвигаются, поэтому $\norm{y_n'}_2$ не зависит от $n$.
При этом $\norm{y_n}_2\to\infty$.
Оценка снизу невозможна ни для какого $\lambda\in\RR$.
Невещественные точки регулярны в силу симметричности,
так что $\Pi(T)=\CC\setminus\RR$.

\begin{figure}[ht]
\centering
\begin{tikzpicture}[>=Stealth]
 \draw[->] (-4.3,0)--(4.3,0) node[right] {$t$};
 \draw[thick] (-4.1,0)--(-3.4,0)
  .. controls (-2.8,0) and (-2.8,1) .. (-2.2,1)
  --(2.2,1) .. controls (2.8,1) and (2.8,0) .. (3.4,0)
  --(4.1,0);
 \foreach \x/\tick in {-3.4/{-n-1},-2.2/{-n},2.2/n,3.4/{n+1}}{
  \draw (\x,-.04)--(\x,.04);
  \node[below] at (\x,0) {$\tick$};
 }
 \node[above] at (0,1) {$1$};
 \node at (-3.65,.55) {$y_n$};
\end{tikzpicture}
\caption{Функции $y_n$ на прямой (\at{9TaAAWmVph8}{2180}{36:20}).}
\label{fig:sem6-cutoff}
\end{figure}

\topic{Одно граничное условие}{\at{9TaAAWmVph8}{2399}{39:59}}
\begin{exercise}
Построить примеры неограниченных операторов с пустым
спектром и со спектром, равным $\CC$.
\end{exercise}
На $[-1,1]$ рассмотрим $T\subset\widetilde T_\theta\subset T_0^*$,
где
\[
 \Dom(\widetilde T_\theta)
 =\{u\in\Dom(T_0^*):u(1)=\theta u(-1)\},\qquad \theta\in\CC.
\]
При $\theta=\infty$ берём условие $u(-1)=0$.
Чтобы найти $v=R_\lambda(\widetilde T_\theta)h$ для
$h\in L^2[-1,1]$, решаем $-\ii v'-\lambda v=h$.
Переписываем его как $v'-\ii\lambda v=\ii h$ и умножаем
на $\ee^{-\ii\lambda t}$. Получаем
\[
 \frac{d}{dt}\bigl(\ee^{-\ii\lambda t}v(t)\bigr)
 =\ii h(t)\ee^{-\ii\lambda t}.
\]
Интегрируя от $-1$ до $t$, получаем
\[
 \begin{aligned}
 \ee^{-\ii\lambda t}v(t)&=\ee^{\ii\lambda}v(-1)
      +\ii\int_{-1}^t h(\tau)\ee^{-\ii\lambda\tau}\dd\tau,\\
 v(t)&=v(-1)\ee^{\ii\lambda(t+1)}
      +\ii\int_{-1}^t h(\tau)\ee^{\ii\lambda(t-\tau)}\dd\tau.
 \end{aligned}
\]
При конечном $\theta$ подставляем $t=1$ и получаем
из краевого условия
\[
 \theta v(-1)=v(-1)\ee^{2\ii\lambda}
       +\ii\int_{-1}^1 h(\tau)\ee^{\ii\lambda(1-\tau)}\dd\tau.
\]
При $\theta\ne\ee^{2\ii\lambda}$ получаем
\[
 v(-1)=
 \frac{\ii\int_{-1}^1 h(\tau)\ee^{\ii\lambda(1-\tau)}\dd\tau}
      {\theta-\ee^{2\ii\lambda}}.
\]
Подстановка в формулу для $v(t)$ даёт
\[
 \bigl(R_\lambda(\widetilde T_\theta)h\bigr)(t)
 =\frac{\ii\ee^{\ii\lambda(t+2)}
            (h,\ee^{\ii\bar\lambda\,\cdot})}
        {\theta-\ee^{2\ii\lambda}}
   +\ii\int_{-1}^t h(\tau)\ee^{\ii\lambda(t-\tau)}\dd\tau
 \quad\text{(\at{9TaAAWmVph8}{3331}{55:31})}.
\]
Второе слагаемое~--- интегральный оператор, ограниченный
по норме ядра. Первое ограничено при ненулевом знаменателе.
При $\theta=\ee^{2\ii\lambda}$ условие на $v(-1)$
выполняется не для всех $h$. Следовательно,
\[
 \sigma(\widetilde T_\theta)=
   \{\lambda\in\CC:\ee^{2\ii\lambda}=\theta\},
   \qquad \theta\in\CC.
\]
При $\theta=0$ спектр пуст. При $\theta\ne0$ корни
отличаются на целые кратные $\pi$, и спектр чисто дискретен.
При $\theta=\infty$ в формуле для $v(t)$ остаётся интеграл,
поскольку $v(-1)=0$.

\begin{exercise}[Задания в конце семинара]
\at{9TaAAWmVph8}{3665}{1:01:05}.
\begin{enumerate}
\item Классифицировать спектр $\widetilde T_\theta$ на отрезке.
\item Доказать, что спектр $T_0^*$ на отрезке совпадает со всей $\CC$.
\item Найти $(\widetilde T_\theta)^*$.
\item Найти и классифицировать спектр производной на прямой.
\end{enumerate}
\end{exercise}

\sessionsection{Семинар 7. Спектр производной и преобразование Кэли}
\record{16 октября}{qrpVbZfggnU}

\topic{Граничные условия для производной}{\at{qrpVbZfggnU}{645}{10:45}}
\begin{exercise}
Сопоставить спектры максимальной производной на отрезке
и производной на прямой, а также найти сопряжённый к
$T_\theta$ и проверить замкнутость $T_\theta$.
\end{exercise}
В начале семинара (\at{qrpVbZfggnU}{113}{1:53})
рассматривается максимальный оператор $T^*$ на отрезке.
В его области $H^1(-1,1)$ граничные условия отсутствуют.
Для каждого $\lambda\in\CC$ функция
$\ee^{\ii\lambda t}\in H^1(-1,1)$ является
собственным вектором. Поэтому точечный спектр
максимального оператора есть вся $\CC$.
На прямой нет ненулевых решений дефектных уравнений,
квадратично интегрируемых на обоих концах. Замыкание
минимальной производной самосопряжённо, а вещественная ось
составляет его спектр по проверке поля регулярности.
На отрезке проверка замкнутости для
$u(1)=\theta u(-1)$ устроена так. Если $u_n\to u$ и
$-\ii u_n'\to v$ в $L^2$, замкнутость $T^*$
даёт $u\in H^1$ и $-\ii u'=v$. Следовательно,
$u_n\to u$ по норме $H^1$ и граничные следы сходятся.
Условие при $\theta=\infty$ проходит к пределу
тем же способом. При интегрировании по частям остаётся
$-\ii[u\bar v]_{-1}^{1}$. Для конечного $\theta$ он
пропорционален $u(-1)(\theta\bar v(1)-\bar v(-1))$.
След $u(-1)$ произволен, откуда
$v(-1)=\bar\theta\,v(1)$. При $\theta=\infty$
произволен $u(1)$ и получаем $v(1)=0$.
Самосопряжённость в этом семействе получается при $\abs\theta=1$.
Для конечного ненулевого $\theta$ корни
$\ee^{2\ii\lambda}=\theta$ дают собственные значения
$\lambda_k=\arg(\theta)/2+\pi k-\tfrac{\ii}{2}\log|\theta|$
при $k\in\ZZ$,
а ограниченный обратный в остальных точках найден на семинаре 6.

\topic{Изометрия и её продолжения}{\at{qrpVbZfggnU}{2142}{35:42}}
\begin{exercise}
Найти преобразование Кэли минимальной производной
на $[-1,1]$, выписать его область и образ и описать
его унитарные расширения.
\end{exercise}
У минимальной производной на $[-1,1]$ для
$(A-\ii)x=h$ с $x(-1)=0$ имеем
\[
 x(t)=\ii\int_{-1}^t\ee^{s-t}h(s)\dd s.
\]
Условие $x(1)=0$ равносильно $h\perp\ee^t$. Поэтому
\[
 (Uh)(t)=h(t)-2\int_{-1}^t\ee^{s-t}h(s)\dd s,\qquad h\perp\ee^t.
\]
Здесь $(A+\ii)x=(A-\ii)x+2\ii x=h+2\ii x$,
что объясняет коэффициент $-2$ перед интегралом.
Образ $U$~--- ортогональное дополнение к $\ee^{-t}$.
Унитарные продолжения определяются добавочным условием
$\widetilde U\,\ee^t=\alpha\ee^{-t}$, $\abs\alpha=1$.
Нормы этих двух экспонент равны.

\begin{exercise}[Домашнее задание]
\at{qrpVbZfggnU}{3735}{1:02:15}.
Применить обратное преобразование Кэли к каждому
унитарному продолжению $\widetilde U$ и показать,
что получается оператор $A_\theta$ с граничным
условием $u(1)=\theta u(-1)$. Выразить $\theta$
через параметр продолжения $\alpha=\ee^{\ii\varphi}$.
\end{exercise}

\sessionsection{Семинар 8. Граничные условия для производной и лапласиана}
\record{23 октября}{RNiI2uAnRJU}

\topic{Два параметра одного расширения}{\at{RNiI2uAnRJU}{30}{0:30}}
\begin{exercise}
Пусть $\widetilde U\,\ee^t=\alpha\ee^{-t}$,
$|\alpha|=1$, для продолжения Кэли с предыдущего
семинара. Найти параметр $\theta$ граничного
условия $u(1)=\theta u(-1)$ полученного расширения.
\end{exercise}
Дополнительное направление в области обратного преобразования
Кэли порождается $(\widetilde U-\id)\ee^t=\alpha\ee^{-t}-\ee^t$.
Остальные векторы области относятся к $\Dom(A)$
и имеют нулевые граничные следы. Поэтому вся область
имеет ту же пропорцию следов, что и дополнительный
вектор. Его след при $-1$ ненулевой, поскольку
$\alpha\ee-\ee^{-1}=0$ невозможно при $|\alpha|=1$.
Подстановка в $u(1)=\theta u(-1)$ даёт
\[
 \theta=\frac{\alpha\ee^{-1}-\ee}{\alpha\ee-\ee^{-1}}.
\]
При $\abs\alpha=1$ числитель и знаменатель одинаковы по
модулю: квадрат каждого модуля равен
$\ee^2+\ee^{-2}-2\operatorname{Re}\alpha$.
Обратно, из формулы
$\alpha=(\theta\ee^{-1}-\ee)/(\theta\ee-\ee^{-1})$
при $|\theta|=1$ вновь получается $|\alpha|=1$.
Таким образом, обе параметризации покрывают все
самосопряжённые расширения.

\topic{Лапласиан на отрезке}{\at{RNiI2uAnRJU}{1890}{31:30}}
\begin{exercise}
Для минимального $A_0=-d^2/dt^2$ выписать
максимальный сопряжённый оператор, область замыкания
и дефектные пространства. Выразить симметричные
расширения через граничные данные.
\end{exercise}
Для $A_0=-d^2/dt^2$ на $\sch{(0,1)}$ максимальный
сопряжённый действует на $H^2(0,1)$.
Минимальная область после замыкания задаётся нулевыми
значениями функций и их первых производных на концах.
При первом сопряжении финитные тестовые функции показывают,
что вторая слабая производная лежит в $L^2$, а это
даёт $H^2(0,1)$. При втором сопряжении используется
двукратное интегрирование по частям и независимость
четырёх следов пробной функции на концах.
В каждом дефектном уравнении $-u''=\pm\ii u$ решения
порождаются $\exp(\kappa t)$ и $\exp(-\kappa t)$,
где $\kappa^2=\mp\ii$. Оба решения лежат в $L^2(0,1)$,
поэтому дефектные индексы равны $(2,2)$.
На занятии сопоставляются два
способа описывать самосопряжённые расширения: добавлять
график унитарного отображения между дефектными
пространствами или задавать условия на векторы из четырёх
граничных значений. Формула интегрирования по частям
\[
 (A^*u,v)-(u,A^*v)
  =[u\overline{v'}-u'\bar v]_0^1
\]
выражает симметричность через зануление граничной формы.
Ядро отображения
$u\mapsto(u(0),u(1),-u'(0),u'(1))$ есть минимальная
область. Поэтому линейное подпространство
граничных данных задаёт расширение, а при
переходе к сопряжённому условия определяются
аннулятором относительно формулы Грина. Этот
переход подробно разбирается на семинаре~9.
Вычисление преобразования Кэли минимального лапласиана
поставлено в конце семинара как отдельная задача.

\sessionsection{Семинар 9. Лагранжевы граничные условия}
\record{6 ноября}{hQx8UC4B5TY}

\topic{Формула Грина и сопряжённое граничное условие}{\at{hQx8UC4B5TY}{240}{4:00}}
\begin{exercise}
Доказать, что самосопряжённым расширениям минимального
лапласиана на отрезке соответствуют лагранжевы плоскости
граничных данных. Описать эти плоскости с помощью
унитарных матриц.
\end{exercise}
Запишем
$\Gamma^{(1)}u=(u(0),u(1))$,
$\Gamma^{(2)}u=(-u'(0),u'(1))$.
На $\CC^2\oplus\CC^2$ получается косоэрмитова форма
\[
 [(\xi,\eta),(\xi',\eta')]
  =(\xi,\eta')-(\eta,\xi').
\]
Область расширения $A_L$ равна $\Gamma^{-1}(L)$,
где $\Gamma u=(\Gamma^{(1)}u,\Gamma^{(2)}u)$.
Если $A_0$~--- замкнутый минимальный оператор, то
$A_0\subset A_L\subset A_0^*$, следовательно,
$A_0\subset A_L^*\subset A_0^*$. По формуле Грина
\[
 (A_Lu,v)-(u,A_0^*v)=[\Gamma u,\Gamma v]
 \qquad(u\in\Dom(A_L),\ v\in H^2).
\]
Поэтому $v\in\Dom(A_L^*)$ тогда и только тогда,
когда $[\Gamma u,\Gamma v]=0$ для всех
$u\in\Dom(A_L)$. Ограничение $\Gamma$ на многочлены
степени не выше трёх инъективно: четыре нулевых следа
дали бы два двукратных корня. По размерности это
ограничение биективно, поэтому $\Gamma$ сюръективно.
Следовательно, $\Dom(A_L^*)=\Gamma^{-1}(L^{\perp_{[\ ,\ ]}})$.
Поэтому расширение самосопряжённо ровно тогда, когда
$L$ совпадает со своим косоортогональным дополнением.
Это лагранжевы плоскости в $\CC^2\oplus\CC^2$.
Граничные данные условия $u(0)=u(1)$, $u'(0)=u'(1)=0$
образуют прямую $L=\{((c,c),(0,0)):c\in\CC\}$.
В сопряжённой области остаётся условие $v'(1)=v'(0)$,
полученное из ортогональности этому $L$ по граничной форме.

\topic{Соответствие лагранжевых плоскостей и $U(2)$}{\at{hQx8UC4B5TY}{2365}{39:25}}
В координатах $z_\pm=\xi\pm\ii\eta$ граничная форма
выражается разностью скалярных произведений:
\[
 (z_+,z'_+)-(z_-,z'_-)
    =-2\ii\big((\xi,\eta')-(\eta,\xi')\big).
\]
На изотропной
плоскости $\norm{z_+}=\norm{z_-}$.
Проекция на $z_+$ инъективна, поскольку из $z_+=0$ следует
$z_-=0$, а затем $\xi=\eta=0$. На лагранжевой плоскости
размерности $2$ она также сюръективна. Поэтому плоскость
задаётся графиком унитарного оператора между этими копиями
$\CC^2$. Если выбрать запись $z_+=Uz_-$, уравнение принимает вид
\[
 (U-\id)\xi-\ii(U+\id)\eta=0.
\]
Обратная проверка получается из сохранения скалярного
произведения унитарным $U$: соответствующий график
изотропен и двумерен, значит, лагранжев.
Явное вычисление унитарного преобразования для
заданной плоскости оставлено на дом
(\at{hQx8UC4B5TY}{4491}{1:14:51}).

\sessionsection{Семинар 10. Лапласиан и форма на отрезке}
\record{13 ноября}{oDIn4ERoCkc}

\topic{Существенная самосопряжённость на $\RR^n$}{\at{oDIn4ERoCkc}{20}{0:20}}
\begin{exercise}
Доказать существенную самосопряжённость $-\Delta$
на $\sch{\RR^n}$ в $L^2(\RR^n)$ и описать область
замыкания.
\end{exercise}
Для $-\Delta$ на $\sch{\RR^n}$ преобразование Фурье
переводит дифференциальное выражение в умножение на $\abs p^2$.
Естественный домен оператора умножения состоит из
$g\in L^2$ с $\abs p^2g\in L^2$. Нужно показать, что образ
первоначальной области является ядром по графовой норме,
то есть замыкание преобразованного графика даёт весь
график оператора умножения. На доске эта норма записывается
как норма $L^2$ с весом $1+\abs p^4$:
\[
 \norm g_{M_{|p|^2}}^2=\int_{\RR^n}(1+|p|^4)|g(p)|^2\dd p .
\]
Вес сравним с $(1+|p|^2)^2$. Если $g$ лежит
в максимальной области умножения, то
$\mathcal F^{-1}g\in H^2$. Приближаем эту функцию
функциями $f_j\in\sch{\RR^n}$ в $H^2$.
Тогда $\mathcal Ff_j\to g$ и
$|p|^2\mathcal Ff_j\to |p|^2g$ в $L^2$.
Обратное включение следует из замкнутости
оператора умножения: для последовательностей
$g_j\to g$ и $|p|^2g_j\to h$ в $L^2$ общая
подпоследовательность сходится почти всюду и
даёт $h=|p|^2g$. Поэтому область замыкания
на стороне $x$ есть $H^2(\RR^n)$.
Самосопряжённость максимального оператора умножения
доводится на следующем семинаре.

\begin{figure}[ht]
\centering
\begin{tikzpicture}[>=Stealth]
 \node (l) at (0,1.45) {$\sch{\RR^n}$};
 \node (r) at (5.3,1.45) {$L^2(\RR^n)$};
 \node (bl) at (0,0) {$\mathcal F(\sch{\RR^n})$};
 \node (br) at (5.3,0) {$L^2(\RR^n)$};
 \draw[->] (l)--node[above] {$-\Delta$} (r);
 \draw[->] (l)--node[left] {$\mathcal F$} (bl);
 \draw[->] (r)--node[right] {$\mathcal F$} (br);
 \draw[->] (bl)--node[below] {$M_{|p|^2}$} (br);
\end{tikzpicture}
\caption{Преобразование графика лапласиана по Фурье.
По схеме на доске в \at{oDIn4ERoCkc}{2790}{46:30}.}
\end{figure}

\topic{Как выделить член формы на границе}{\at{oDIn4ERoCkc}{3170}{52:50}}
\begin{exercise}
Найти форму для самосопряжённой реализации
$-d^2/dt^2$ на $[a,b]$ с унитарным граничным условием.
\end{exercise}
В этом вычислении ориентируем производные внутрь отрезка:
$\xi=\Gamma^{(1)}f=(f(a),f(b))$,
$\eta=\Gamma^{(2)}f=(f'(a),-f'(b))$.
Тогда в формуле для $(-f'',g)$ граничный член
есть $+(\eta,\Gamma^{(1)}g)$. Записываем условие в виде
\[
 \ii(U-\id)\xi+(U+\id)\eta=0.
\]
Пусть $P_D$, $P_N$~--- ортогональные проекторы на
собственные пространства $U$ с собственными значениями
$-1$ и $+1$, а $P_R=\id-P_D-P_N$. Подставляя $U=-\id$
на $\Ran P_D$ и $U=\id$ на $\Ran P_N$, получаем
$P_D\xi=0$ и $P_N\eta=0$. На $\Ran P_R$ операторы
$U\pm\id$ обратимы, поэтому
\[
 P_R\eta=M P_R\xi,\qquad M=-\ii(U+\id)^{-1}(U-\id)
 \quad\text{на }\Ran P_R.
\]
Числа $-\ii(z-1)/(z+1)$
вещественны для $\abs z=1$, $z\ne-1$,
и ненулевые при $z\ne1$, так что $M$ эрмитов и
обратим на $\Ran P_R$. При интегрировании по частям
\[
 (-f'',g)=\int_a^b f'\overline{g'}\dd t+
                        (\Gamma^{(2)}f,\Gamma^{(1)}g).
\]
Поскольку $P_D\Gamma^{(1)}g=0$ и
$P_N\Gamma^{(2)}f=0$, в последнем скалярном
произведении остаётся только $P_R$-часть.
Подставляя туда $M P_R\Gamma^{(1)}f$, получаем
\[
 \begin{aligned}
 h[f,g]&=\int_a^b f'\overline{g'}\dd t
           +(M P_R\Gamma^{(1)}f,P_R\Gamma^{(1)}g),\\
 \Dom[h]&=\{f\in H^1(a,b):P_D\Gamma^{(1)}f=0\}.
 \end{aligned}
\]
К концу записи поставлен вопрос о малой относительной
оценке граничного члена, которая нужна для
полуограниченности и замкнутости.

\sessionsection{Семинар 11. Замкнутость формы и восстановление оператора}
\record{20 ноября}{Whuc-LrLTZ0}

\topic{Самосопряжённость оператора умножения}{\at{Whuc-LrLTZ0}{25}{0:25}}
Умножение на вещественную функцию $m(p)=\abs p^2$
на естественном домене симметрично, поскольку
$\int mf\overline g=\int f\overline{mg}$.
Если $g$ лежит в области сопряжённого, то
$(mf,g)=(f,h)$ для некоторого $h\in L^2$ и всех
$f$ из естественного домена. На каждом шаре
множитель $m$ ограничен. Подставляя
$f=(h-mg)\mathbf 1_{\{|p|\le R\}}$, получаем
$\|(h-mg)\mathbf 1_{\{|p|\le R\}}\|_2^2=0$.
Так как $R$ произволен, $h=mg$ почти всюду.
Поэтому $mg\in L^2$ и $g$ принадлежит исходному домену.
Сочетание с замыканием
преобразованного графика с семинара 10 даёт
самосопряжённость замыкания $-\Delta$ на $\RR^n$.

\topic{Следы, нижняя грань и замкнутость}{\at{Whuc-LrLTZ0}{1090}{18:10}}
\begin{lemma}[Оценка следов с малым параметром]
Для каждого $\varepsilon>0$ существует $C_\varepsilon$
такое, что для всех $f\in H^1(a,b)$
\[
 \abs{f(a)}^2+\abs{f(b)}^2
 \le \varepsilon\norm{f'}_2^2+C_\varepsilon\norm f_2^2 .
\]
\end{lemma}
\begin{proof}
Если для некоторого $\varepsilon>0$ нет такой
константы, выбираем $f_n$, для которых оценка с
$C_\varepsilon=n$ нарушается. Неравенство однородно,
так что можно нормировать левую часть на $1$:
\[
 |f_n(a)|^2+|f_n(b)|^2=1,\qquad
 \varepsilon\norm{f_n'}_2^2+n\norm{f_n}_2^2<1.
\]
Значит, $f_n$ ограничена в $H^1(a,b)$ и стремится
к нулю в $L^2(a,b)$. По компактности вложения
$H^1(a,b)\hookrightarrow C[a,b]$ выделяем
равномерно сходящуюся подпоследовательность.
Её предел нулевой, так как её $L^2$-предел нулевой.
Тогда сумма квадратов граничных значений стремится
к нулю, что противоречит нормировке.
\end{proof}
\begin{proposition}
Форма $h$ предыдущего семинара полуограничена снизу
и замкнута на $\Dom[h]$.
\end{proposition}
\begin{proof}
По оценке следов
\[
 |(M P_R\xi,P_R\xi)|\le
 \norm M(\varepsilon\norm{f'}_2^2+C_\varepsilon\norm f_2^2).
\]
Выбираем $\varepsilon$ так, чтобы
$\norm M\varepsilon\le1/2$. Получаем
$h[f,f]\ge\tfrac12\norm{f'}_2^2-\norm M C_\varepsilon\norm f_2^2$.
После прибавления $(1+\norm M C_\varepsilon)\norm f_2^2$
эта оценка и оценка сверху дают эквивалентность
нормы формы и нормы $H^1$. Условие
$P_D\Gamma^{(1)}f=0$ задаёт ядро непрерывного
отображения из $H^1$ в $\CC^2$. Поэтому область
замкнута в $H^1$ и полна по норме сдвинутой формы.
\end{proof}

\topic{Оператор, согласованный с формой}{\at{Whuc-LrLTZ0}{3437}{57:17}}
\begin{exercise}
Восстановить оператор по замкнутой форме $h$,
то есть доказать действие $Af=-f''$ и вывести
три граничных условия на его область.
\end{exercise}
Проверяем $h[f,g]=(Af,g)$ сначала для $g\in\sch{(a,b)}$.
Граничные значения такого $g$ исчезают. По определению
слабой производной $Af=-f''$, и $f''\in L^2$, то есть
$f\in H^2$. Затем берём все $g\in\Dom[h]$. После
интегрирования по частям требуется
\[
 (\Gamma^{(2)}f,\Gamma^{(1)}g)
      =(M P_R\Gamma^{(1)}f,P_R\Gamma^{(1)}g).
\]
На $\Dom[h]$ вектор $\Gamma^{(1)}g$ пробегает
$\Ran P_N\oplus\Ran P_R$, поскольку разрешённые значения
на концах можно реализовать, например, линейной функцией.
Берём след $\Gamma^{(1)}g=P_N\Gamma^{(2)}f$.
Слева стоит квадрат нормы этого вектора, справа
нуль. Значит, $P_N\Gamma^{(2)}f=0$. Если же
$\Gamma^{(1)}g$ произвольно пробегает $\Ran P_R$,
получаем $P_R\Gamma^{(2)}f=M P_R\Gamma^{(1)}f$.
Условие $P_D\Gamma^{(1)}f=0$ уже входит в $\Dom[h]$.
Обратно, из этих условий следует нужное равенство для всех
$g$ из области формы.

\begin{figure}[ht]
\centering
\begin{tikzpicture}
 \draw (-.55,0)--(4.65,0);
 \draw[thick] (0,0)--(.55,0)
  .. controls (.86,0) and (.72,.99) .. (1.12,.99)
  .. controls (1.51,.99) and (1.38,-.88) .. (1.94,-.70)
  .. controls (2.36,-.58) and (2.22,1.05) .. (2.78,1.05)
  .. controls (3.19,1.05) and (3.22,0) .. (3.56,0)--(4.1,0);
 \foreach \x/\label in {0/a,4.1/b} {
   \draw (\x,.1)--(\x,-.1);
   \node[below] at (\x,-.1) {$\label$};
 }
\end{tikzpicture}
\caption{Финитная пробная функция, равная нулю в окрестностях
концов. По доске в \at{Whuc-LrLTZ0}{3970}{1:06:10}.}
\end{figure}

\topic{Вопросы в конце семинара}{\at{Whuc-LrLTZ0}{5088}{1:24:48}}
\noindent
\begin{minipage}{.69\linewidth}
Первая задача~--- найти расширение
Фридрихса замыкания минимального лапласиана на отрезке.
Вторая~--- найти форму для $-d^2/dt^2$ с условиями
$f(a)=f(b)$ и $f'(b)-f'(a)=\lambda f(a)$. Концы отрезка
склеиваются, а производная имеет скачок в точке склейки.
Решения этих двух задач в записи не разбираются.
\end{minipage}\hfill
\begin{minipage}{.27\linewidth}
 \centering
 \begin{tikzpicture}[scale=.62]
  \draw[thick] (0,0) .. controls (-1.35,.48) and (-1.15,1.93) .. (0,2)
   .. controls (1.3,2.08) and (1.28,.45) .. (0,0);
  \fill (0,0) circle (1.4pt);
 \end{tikzpicture}
 \captionof{figure}{Склейка концов.
 \at{Whuc-LrLTZ0}{5260}{1:27:40}.}
\end{minipage}

\end{document}
